% Options for packages loaded elsewhere
% Options for packages loaded elsewhere
\PassOptionsToPackage{unicode}{hyperref}
\PassOptionsToPackage{hyphens}{url}
%
\documentclass[
  letterpaper,
]{scrbook}
\usepackage{xcolor}
\usepackage{amsmath,amssymb}
\setcounter{secnumdepth}{5}
\usepackage{iftex}
\ifPDFTeX
  \usepackage[T1]{fontenc}
  \usepackage[utf8]{inputenc}
  \usepackage{textcomp} % provide euro and other symbols
\else % if luatex or xetex
  \usepackage{unicode-math} % this also loads fontspec
  \defaultfontfeatures{Scale=MatchLowercase}
  \defaultfontfeatures[\rmfamily]{Ligatures=TeX,Scale=1}
\fi
\usepackage{lmodern}
\ifPDFTeX\else
  % xetex/luatex font selection
\fi
% Use upquote if available, for straight quotes in verbatim environments
\IfFileExists{upquote.sty}{\usepackage{upquote}}{}
\IfFileExists{microtype.sty}{% use microtype if available
  \usepackage[]{microtype}
  \UseMicrotypeSet[protrusion]{basicmath} % disable protrusion for tt fonts
}{}
\makeatletter
\@ifundefined{KOMAClassName}{% if non-KOMA class
  \IfFileExists{parskip.sty}{%
    \usepackage{parskip}
  }{% else
    \setlength{\parindent}{0pt}
    \setlength{\parskip}{6pt plus 2pt minus 1pt}}
}{% if KOMA class
  \KOMAoptions{parskip=half}}
\makeatother
% Make \paragraph and \subparagraph free-standing
\makeatletter
\ifx\paragraph\undefined\else
  \let\oldparagraph\paragraph
  \renewcommand{\paragraph}{
    \@ifstar
      \xxxParagraphStar
      \xxxParagraphNoStar
  }
  \newcommand{\xxxParagraphStar}[1]{\oldparagraph*{#1}\mbox{}}
  \newcommand{\xxxParagraphNoStar}[1]{\oldparagraph{#1}\mbox{}}
\fi
\ifx\subparagraph\undefined\else
  \let\oldsubparagraph\subparagraph
  \renewcommand{\subparagraph}{
    \@ifstar
      \xxxSubParagraphStar
      \xxxSubParagraphNoStar
  }
  \newcommand{\xxxSubParagraphStar}[1]{\oldsubparagraph*{#1}\mbox{}}
  \newcommand{\xxxSubParagraphNoStar}[1]{\oldsubparagraph{#1}\mbox{}}
\fi
\makeatother

\usepackage{color}
\usepackage{fancyvrb}
\newcommand{\VerbBar}{|}
\newcommand{\VERB}{\Verb[commandchars=\\\{\}]}
\DefineVerbatimEnvironment{Highlighting}{Verbatim}{commandchars=\\\{\}}
% Add ',fontsize=\small' for more characters per line
\usepackage{framed}
\definecolor{shadecolor}{RGB}{241,243,245}
\newenvironment{Shaded}{\begin{snugshade}}{\end{snugshade}}
\newcommand{\AlertTok}[1]{\textcolor[rgb]{0.68,0.00,0.00}{#1}}
\newcommand{\AnnotationTok}[1]{\textcolor[rgb]{0.37,0.37,0.37}{#1}}
\newcommand{\AttributeTok}[1]{\textcolor[rgb]{0.40,0.45,0.13}{#1}}
\newcommand{\BaseNTok}[1]{\textcolor[rgb]{0.68,0.00,0.00}{#1}}
\newcommand{\BuiltInTok}[1]{\textcolor[rgb]{0.00,0.23,0.31}{#1}}
\newcommand{\CharTok}[1]{\textcolor[rgb]{0.13,0.47,0.30}{#1}}
\newcommand{\CommentTok}[1]{\textcolor[rgb]{0.37,0.37,0.37}{#1}}
\newcommand{\CommentVarTok}[1]{\textcolor[rgb]{0.37,0.37,0.37}{\textit{#1}}}
\newcommand{\ConstantTok}[1]{\textcolor[rgb]{0.56,0.35,0.01}{#1}}
\newcommand{\ControlFlowTok}[1]{\textcolor[rgb]{0.00,0.23,0.31}{\textbf{#1}}}
\newcommand{\DataTypeTok}[1]{\textcolor[rgb]{0.68,0.00,0.00}{#1}}
\newcommand{\DecValTok}[1]{\textcolor[rgb]{0.68,0.00,0.00}{#1}}
\newcommand{\DocumentationTok}[1]{\textcolor[rgb]{0.37,0.37,0.37}{\textit{#1}}}
\newcommand{\ErrorTok}[1]{\textcolor[rgb]{0.68,0.00,0.00}{#1}}
\newcommand{\ExtensionTok}[1]{\textcolor[rgb]{0.00,0.23,0.31}{#1}}
\newcommand{\FloatTok}[1]{\textcolor[rgb]{0.68,0.00,0.00}{#1}}
\newcommand{\FunctionTok}[1]{\textcolor[rgb]{0.28,0.35,0.67}{#1}}
\newcommand{\ImportTok}[1]{\textcolor[rgb]{0.00,0.46,0.62}{#1}}
\newcommand{\InformationTok}[1]{\textcolor[rgb]{0.37,0.37,0.37}{#1}}
\newcommand{\KeywordTok}[1]{\textcolor[rgb]{0.00,0.23,0.31}{\textbf{#1}}}
\newcommand{\NormalTok}[1]{\textcolor[rgb]{0.00,0.23,0.31}{#1}}
\newcommand{\OperatorTok}[1]{\textcolor[rgb]{0.37,0.37,0.37}{#1}}
\newcommand{\OtherTok}[1]{\textcolor[rgb]{0.00,0.23,0.31}{#1}}
\newcommand{\PreprocessorTok}[1]{\textcolor[rgb]{0.68,0.00,0.00}{#1}}
\newcommand{\RegionMarkerTok}[1]{\textcolor[rgb]{0.00,0.23,0.31}{#1}}
\newcommand{\SpecialCharTok}[1]{\textcolor[rgb]{0.37,0.37,0.37}{#1}}
\newcommand{\SpecialStringTok}[1]{\textcolor[rgb]{0.13,0.47,0.30}{#1}}
\newcommand{\StringTok}[1]{\textcolor[rgb]{0.13,0.47,0.30}{#1}}
\newcommand{\VariableTok}[1]{\textcolor[rgb]{0.07,0.07,0.07}{#1}}
\newcommand{\VerbatimStringTok}[1]{\textcolor[rgb]{0.13,0.47,0.30}{#1}}
\newcommand{\WarningTok}[1]{\textcolor[rgb]{0.37,0.37,0.37}{\textit{#1}}}

\usepackage{longtable,booktabs,array}
\newcounter{none} % for unnumbered tables
\usepackage{calc} % for calculating minipage widths
% Correct order of tables after \paragraph or \subparagraph
\usepackage{etoolbox}
\makeatletter
\patchcmd\longtable{\par}{\if@noskipsec\mbox{}\fi\par}{}{}
\makeatother
% Allow footnotes in longtable head/foot
\IfFileExists{footnotehyper.sty}{\usepackage{footnotehyper}}{\usepackage{footnote}}
\makesavenoteenv{longtable}
\usepackage{graphicx}
\makeatletter
\newsavebox\pandoc@box
\newcommand*\pandocbounded[1]{% scales image to fit in text height/width
  \sbox\pandoc@box{#1}%
  \Gscale@div\@tempa{\textheight}{\dimexpr\ht\pandoc@box+\dp\pandoc@box\relax}%
  \Gscale@div\@tempb{\linewidth}{\wd\pandoc@box}%
  \ifdim\@tempb\p@<\@tempa\p@\let\@tempa\@tempb\fi% select the smaller of both
  \ifdim\@tempa\p@<\p@\scalebox{\@tempa}{\usebox\pandoc@box}%
  \else\usebox{\pandoc@box}%
  \fi%
}
% Set default figure placement to htbp
\def\fps@figure{htbp}
\makeatother





\setlength{\emergencystretch}{3em} % prevent overfull lines

\providecommand{\tightlist}{%
  \setlength{\itemsep}{0pt}\setlength{\parskip}{0pt}}



 


\usepackage[brazil]{babel}
\usepackage{fontspec}
\setmainfont{Latin Modern Roman}
\makeatletter
\@ifpackageloaded{tcolorbox}{}{\usepackage[skins,breakable]{tcolorbox}}
\@ifpackageloaded{fontawesome5}{}{\usepackage{fontawesome5}}
\definecolor{quarto-callout-color}{HTML}{909090}
\definecolor{quarto-callout-note-color}{HTML}{0758E5}
\definecolor{quarto-callout-important-color}{HTML}{CC1914}
\definecolor{quarto-callout-warning-color}{HTML}{EB9113}
\definecolor{quarto-callout-tip-color}{HTML}{00A047}
\definecolor{quarto-callout-caution-color}{HTML}{FC5300}
\definecolor{quarto-callout-color-frame}{HTML}{acacac}
\definecolor{quarto-callout-note-color-frame}{HTML}{4582ec}
\definecolor{quarto-callout-important-color-frame}{HTML}{d9534f}
\definecolor{quarto-callout-warning-color-frame}{HTML}{f0ad4e}
\definecolor{quarto-callout-tip-color-frame}{HTML}{02b875}
\definecolor{quarto-callout-caution-color-frame}{HTML}{fd7e14}
\makeatother
\makeatletter
\@ifpackageloaded{bookmark}{}{\usepackage{bookmark}}
\makeatother
\makeatletter
\@ifpackageloaded{caption}{}{\usepackage{caption}}
\AtBeginDocument{%
\ifdefined\contentsname
  \renewcommand*\contentsname{Table of contents}
\else
  \newcommand\contentsname{Table of contents}
\fi
\ifdefined\listfigurename
  \renewcommand*\listfigurename{List of Figures}
\else
  \newcommand\listfigurename{List of Figures}
\fi
\ifdefined\listtablename
  \renewcommand*\listtablename{List of Tables}
\else
  \newcommand\listtablename{List of Tables}
\fi
\ifdefined\figurename
  \renewcommand*\figurename{Figure}
\else
  \newcommand\figurename{Figure}
\fi
\ifdefined\tablename
  \renewcommand*\tablename{Table}
\else
  \newcommand\tablename{Table}
\fi
}
\@ifpackageloaded{float}{}{\usepackage{float}}
\floatstyle{ruled}
\@ifundefined{c@chapter}{\newfloat{codelisting}{h}{lop}}{\newfloat{codelisting}{h}{lop}[chapter]}
\floatname{codelisting}{Listing}
\newcommand*\listoflistings{\listof{codelisting}{List of Listings}}
\makeatother
\makeatletter
\makeatother
\makeatletter
\@ifpackageloaded{caption}{}{\usepackage{caption}}
\@ifpackageloaded{subcaption}{}{\usepackage{subcaption}}
\makeatother
\usepackage{bookmark}
\IfFileExists{xurl.sty}{\usepackage{xurl}}{} % add URL line breaks if available
\urlstyle{same}
\hypersetup{
  pdftitle={Curso R para Análise de Dados},
  pdfauthor={Ítalo Cegatta},
  hidelinks,
  pdfcreator={LaTeX via pandoc}}


\title{Curso R para Análise de Dados}
\usepackage{etoolbox}
\makeatletter
\providecommand{\subtitle}[1]{% add subtitle to \maketitle
  \apptocmd{\@title}{\par {\large #1 \par}}{}{}
}
\makeatother
\subtitle{Do zero ao relatório técnico em 4 dias}
\author{Ítalo Cegatta}
\date{2025-12-31}
\begin{document}
\frontmatter
\maketitle

\renewcommand*\contentsname{Table of contents}
{
\setcounter{tocdepth}{1}
\tableofcontents
}

\mainmatter
\bookmarksetup{startatroot}

\chapter{Curso R para Análise de
Dados}\label{curso-r-para-anuxe1lise-de-dados}

Do zero ao relatório técnico

\hfill\break

\bookmarksetup{startatroot}

\chapter*{Bem-vindo(a)}\label{bem-vindoa}
\addcontentsline{toc}{chapter}{Bem-vindo(a)}

\markboth{Bem-vindo(a)}{Bem-vindo(a)}

Este livro reúne o material de três cursos de introdução ao R para
análise de dados. O conteúdo parte do zero --- sem necessidade de
programação prévia --- e avança até a produção de análises completas com
modelagem estatística.

\section{Escolha seu curso}\label{escolha-seu-curso}

{\def\LTcaptype{none} % do not increment counter
\begin{longtable}[]{@{}
  >{\raggedright\arraybackslash}p{(\linewidth - 6\tabcolsep) * \real{0.0625}}
  >{\centering\arraybackslash}p{(\linewidth - 6\tabcolsep) * \real{0.2812}}
  >{\centering\arraybackslash}p{(\linewidth - 6\tabcolsep) * \real{0.3438}}
  >{\raggedright\arraybackslash}p{(\linewidth - 6\tabcolsep) * \real{0.3125}}@{}}
\toprule\noalign{}
\begin{minipage}[b]{\linewidth}\raggedright
\end{minipage} & \begin{minipage}[b]{\linewidth}\centering
Duração
\end{minipage} & \begin{minipage}[b]{\linewidth}\centering
Capítulos
\end{minipage} & \begin{minipage}[b]{\linewidth}\raggedright
Conteúdo
\end{minipage} \\
\midrule\noalign{}
\endhead
\bottomrule\noalign{}
\endlastfoot
\textbf{\href{01-dia/index.qmd}{R Express --- 1 dia}} & 8h & 3 &
Primeiro contato com R: importar, manipular, visualizar \\
\textbf{\href{02-dias/index.qmd}{R Fundamentos --- 2 dias}} & 16h & 5 &
Dominar o fluxo completo: dplyr, joins, ggplot2 \\
\textbf{\href{04-dias/index.qmd}{R Completo --- 4 dias}} & 32h & 8 &
Dominar R --- da importação à modelagem estatística \\
\end{longtable}
}

Todos compartilham a mesma base. O curso de \textbf{4 dias} é o mais
completo. O de \textbf{2 dias} condensa os tópicos principais. O de
\textbf{1 dia} é uma introdução rápida e prática.

\section{O que você vai aprender?}\label{o-que-vocuxea-vai-aprender}

\begin{itemize}
\tightlist
\item
  Usar o R e o RStudio
\item
  Importar dados de planilhas, CSVs e outras fontes
\item
  Transformar, filtrar e resumir dados com \texttt{dplyr}
\item
  Criar gráficos profissionais com \texttt{ggplot2}
\item
  Cruzar tabelas com joins
\item
  Lidar com texto, datas e fatores
\item
  Automatizar tarefas com funções e \texttt{purrr}
\item
  Ajustar modelos estatísticos (regressão, ANOVA)
\item
  Produzir relatórios técnicos completos
\end{itemize}

\section{Pré-requisitos}\label{pruxe9-requisitos}

\begin{itemize}
\tightlist
\item
  Um computador com Windows, Mac ou Linux
\item
  \href{https://cran.r-project.org/}{R} (4.0 ou superior)
\item
  \href{https://posit.co/download/rstudio-desktop/}{RStudio}
\end{itemize}

\begin{tcolorbox}[enhanced jigsaw, arc=.35mm, bottomrule=.15mm, bottomtitle=1mm, breakable, colback=white, colbacktitle=quarto-callout-tip-color!10!white, colframe=quarto-callout-tip-color-frame, coltitle=black, left=2mm, leftrule=.75mm, opacityback=0, opacitybacktitle=0.6, rightrule=.15mm, title=\textcolor{quarto-callout-tip-color}{\faLightbulb}\hspace{0.5em}{Não sabe instalar?}, titlerule=0mm, toprule=.15mm, toptitle=1mm]

O \textbf{Capítulo 1} de qualquer curso ensina como instalar o R e o
RStudio passo a passo.

\end{tcolorbox}

\section{Dados utilizados}\label{dados-utilizados}

\begin{itemize}
\tightlist
\item
  🎬 \textbf{\texttt{filmes\_imdb}} --- filmes (nota, orçamento,
  faturamento, gênero). Dado universal para engajamento.
\item
  🌲 \textbf{\texttt{inv\_florestal.csv}},
  \textbf{\texttt{parcelas.csv}}, \textbf{\texttt{clima.xlsx}},
  \textbf{\texttt{talhoes.xlsx}}, \textbf{\texttt{experimento.xlsx}} ---
  aplicação profissional em dados florestais.
\end{itemize}

\section{Repositório}\label{reposituxf3rio}

👉
\href{https://github.com/italocegatta/curso_r}{github.com/italocegatta/curso\_r}

\section{Licença}\label{licenuxe7a}

\href{https://creativecommons.org/licenses/by-nc/4.0/}{Creative Commons
BY-NC 4.0}

\part{1 dia}

\chapter{R Express --- 1 dia}\label{r-express-1-dia}

\chapter*{R Express}\label{r-express}
\addcontentsline{toc}{chapter}{R Express}

\markboth{R Express}{R Express}

\textbf{8 horas} · \textbf{3 capítulos}

Introdução rápida e prática ao R. Ao final do dia, você terá importado
dados, criado tabelas e gráficos. O essencial para começar.

\section{Estrutura}\label{estrutura}

{\def\LTcaptype{none} % do not increment counter
\begin{longtable}[]{@{}
  >{\raggedright\arraybackslash}p{(\linewidth - 4\tabcolsep) * \real{0.1200}}
  >{\raggedright\arraybackslash}p{(\linewidth - 4\tabcolsep) * \real{0.4400}}
  >{\raggedright\arraybackslash}p{(\linewidth - 4\tabcolsep) * \real{0.4400}}@{}}
\toprule\noalign{}
\begin{minipage}[b]{\linewidth}\raggedright
\end{minipage} & \begin{minipage}[b]{\linewidth}\raggedright
Manhã (4h)
\end{minipage} & \begin{minipage}[b]{\linewidth}\raggedright
Tarde (4h)
\end{minipage} \\
\midrule\noalign{}
\endhead
\bottomrule\noalign{}
\endlastfoot
\textbf{Dia 1} & \href{d1-manha-parte1.qmd}{Dia 1 --- Manhã} ·
\href{d1-manha-parte2.qmd}{Dia 1 --- Manhã (continuação)} &
\href{d1-tarde.qmd}{Dia 1 --- Tarde} \\
\end{longtable}
}

\section{Capítulos}\label{capuxedtulos}

\begin{itemize}
\tightlist
\item
  \href{d1-manha-parte1.qmd}{Dia 1 --- Manhã} --- Introdução ao R e
  primeiros dados
\item
  \href{d1-manha-parte2.qmd}{Dia 1 --- Manhã (continuação)} ---
  Manipulação de dados com dplyr
\item
  \href{d1-tarde.qmd}{Dia 1 --- Tarde} --- Visualização de dados e
  Exercícios
\end{itemize}

\section{Datasets}\label{datasets}

\begin{itemize}
\tightlist
\item
  🎬 \texttt{filmes\_imdb.csv} --- filmes (nota, orçamento, faturamento,
  gênero)
\end{itemize}

\section{Pré-requisitos}\label{pruxe9-requisitos-1}

\begin{itemize}
\tightlist
\item
  \href{https://cran.r-project.org/}{R} (4.0+)
\item
  \href{https://posit.co/download/rstudio-desktop/}{RStudio}
\end{itemize}

\section{Repositório}\label{reposituxf3rio-1}

👉
\href{https://github.com/italocegatta/curso_r}{github.com/italocegatta/curso\_r}

\chapter{Dia 1 --- Manhã}\label{dia-1-manhuxe3}

\textbf{Manhã --- Parte 1 \textbar{} 2 horas}

\section*{O essencial}\label{o-essencial}
\addcontentsline{toc}{section}{O essencial}

\markright{O essencial}

\begin{Shaded}
\begin{Highlighting}[]
\CommentTok{\# R como calculadora}
\DecValTok{2} \SpecialCharTok{+} \DecValTok{3}\NormalTok{; }\FunctionTok{sqrt}\NormalTok{(}\DecValTok{81}\NormalTok{); }\DecValTok{2}\SpecialCharTok{\^{}}\DecValTok{3}
\end{Highlighting}
\end{Shaded}

\begin{verbatim}
[1] 5
\end{verbatim}

\begin{verbatim}
[1] 9
\end{verbatim}

\begin{verbatim}
[1] 8
\end{verbatim}

\begin{Shaded}
\begin{Highlighting}[]
\CommentTok{\# Atribuição}
\NormalTok{nota }\OtherTok{\textless{}{-}} \FloatTok{8.5}
\NormalTok{filmes }\OtherTok{\textless{}{-}} \FunctionTok{c}\NormalTok{(}\StringTok{"Matrix"}\NormalTok{, }\StringTok{"Interestelar"}\NormalTok{, }\StringTok{"Parasita"}\NormalTok{)}

\CommentTok{\# Vetores}
\NormalTok{notas }\OtherTok{\textless{}{-}} \FunctionTok{c}\NormalTok{(}\FloatTok{7.5}\NormalTok{, }\FloatTok{8.0}\NormalTok{, }\FloatTok{9.2}\NormalTok{, }\FloatTok{6.8}\NormalTok{, }\FloatTok{8.7}\NormalTok{)}
\FunctionTok{mean}\NormalTok{(notas)  }\CommentTok{\# média}
\end{Highlighting}
\end{Shaded}

\begin{verbatim}
[1] 8.04
\end{verbatim}

\begin{Shaded}
\begin{Highlighting}[]
\FunctionTok{sum}\NormalTok{(notas)   }\CommentTok{\# soma}
\end{Highlighting}
\end{Shaded}

\begin{verbatim}
[1] 40.2
\end{verbatim}

\begin{Shaded}
\begin{Highlighting}[]
\FunctionTok{length}\NormalTok{(notas) }\CommentTok{\# quantos}
\end{Highlighting}
\end{Shaded}

\begin{verbatim}
[1] 5
\end{verbatim}

\section*{Estruturas de dados}\label{estruturas-de-dados}
\addcontentsline{toc}{section}{Estruturas de dados}

\markright{Estruturas de dados}

\begin{Shaded}
\begin{Highlighting}[]
\CommentTok{\# Data frame}
\NormalTok{df }\OtherTok{\textless{}{-}} \FunctionTok{data.frame}\NormalTok{(}
  \AttributeTok{titulo =} \FunctionTok{c}\NormalTok{(}\StringTok{"Matrix"}\NormalTok{, }\StringTok{"Interestelar"}\NormalTok{, }\StringTok{"Parasita"}\NormalTok{),}
  \AttributeTok{ano =} \FunctionTok{c}\NormalTok{(}\DecValTok{1999}\NormalTok{, }\DecValTok{2014}\NormalTok{, }\DecValTok{2019}\NormalTok{),}
  \AttributeTok{nota =} \FunctionTok{c}\NormalTok{(}\FloatTok{8.7}\NormalTok{, }\FloatTok{8.6}\NormalTok{, }\FloatTok{8.5}\NormalTok{)}
\NormalTok{)}
\NormalTok{df}
\end{Highlighting}
\end{Shaded}

\begin{verbatim}
        titulo  ano nota
1       Matrix 1999  8.7
2 Interestelar 2014  8.6
3     Parasita 2019  8.5
\end{verbatim}

\begin{Shaded}
\begin{Highlighting}[]
\CommentTok{\# Fatores}
\NormalTok{genero }\OtherTok{\textless{}{-}} \FunctionTok{factor}\NormalTok{(}\FunctionTok{c}\NormalTok{(}\StringTok{"Ação"}\NormalTok{, }\StringTok{"Drama"}\NormalTok{, }\StringTok{"Comédia"}\NormalTok{, }\StringTok{"Drama"}\NormalTok{, }\StringTok{"Ação"}\NormalTok{))}
\FunctionTok{table}\NormalTok{(genero)}
\end{Highlighting}
\end{Shaded}

\begin{verbatim}
genero
   Ação Comédia   Drama 
      2       1       2 
\end{verbatim}

\section*{Carregando dados reais}\label{carregando-dados-reais}
\addcontentsline{toc}{section}{Carregando dados reais}

\markright{Carregando dados reais}

\begin{Shaded}
\begin{Highlighting}[]
\FunctionTok{library}\NormalTok{(dplyr)}
\FunctionTok{library}\NormalTok{(ggplot2)}
\FunctionTok{library}\NormalTok{(readr)}
\FunctionTok{library}\NormalTok{(tidyr)}
\FunctionTok{library}\NormalTok{(stringr)}
\FunctionTok{library}\NormalTok{(forcats)}
\FunctionTok{library}\NormalTok{(purrr)}
\FunctionTok{library}\NormalTok{(tibble)}
\NormalTok{filmes }\OtherTok{\textless{}{-}} \FunctionTok{read\_csv2}\NormalTok{(}\StringTok{"../data/filmes\_imdb.csv"}\NormalTok{)}
\FunctionTok{glimpse}\NormalTok{(filmes)}
\end{Highlighting}
\end{Shaded}

\begin{verbatim}
Rows: 3,875
Columns: 10
$ codigo_filme     <dbl> 1, 2, 3, 4, 5, 6, 7, 8, 9, 10, 11, 12, 13, 14, 15, 16~
$ titulo           <chr> "Avatar ", "Pirates of the Caribbean: At World's End ~
$ diretor          <chr> "James Cameron", "Gore Verbinski", "Sam Mendes", "Chr~
$ ano              <dbl> 2009, 2007, 2015, 2012, 2012, 2007, 2010, 2015, 2009,~
$ genero           <chr> "Ação", "Aventura", "Comédia", "Mistério", "Fantasia"~
$ duracao          <dbl> 81, 126, 88, 99, 92, 95, 39, 111, 82, 89, 77, 102, 98~
$ orcamento        <dbl> 237000000, 300000000, 245000000, 250000000, 263700000~
$ faturamento      <dbl> 760505847, 309404152, 200074175, 448130642, 73058679,~
$ nota             <dbl> 7.9, 7.1, 6.8, 8.5, 6.6, 6.2, 7.8, 7.5, 7.5, 6.9, 6.1~
$ elenco_principal <chr> "Viola Davis", "Cate Blanchett", "Natalie Portman", "~
\end{verbatim}

\begin{Shaded}
\begin{Highlighting}[]
\FunctionTok{dim}\NormalTok{(filmes)}
\end{Highlighting}
\end{Shaded}

\begin{verbatim}
[1] 3875   10
\end{verbatim}

\begin{Shaded}
\begin{Highlighting}[]
\FunctionTok{summary}\NormalTok{(filmes)}
\end{Highlighting}
\end{Shaded}

\begin{verbatim}
  codigo_filme       titulo            diretor               ano      
 Min.   :   1.0   Length:3875        Length:3875        Min.   :1920  
 1st Qu.: 971.5   Class :character   Class :character   1st Qu.:1999  
 Median :1942.0   Mode  :character   Mode  :character   Median :2005  
 Mean   :1941.6                                         Mean   :2003  
 3rd Qu.:2911.5                                         3rd Qu.:2010  
 Max.   :3881.0                                         Max.   :2016  
    genero             duracao        orcamento          faturamento       
 Length:3875        Min.   : 31.0   Min.   :      218   Min.   :      162  
 Class :character   1st Qu.: 89.0   1st Qu.: 10000000   1st Qu.:  6856116  
 Mode  :character   Median :105.0   Median : 24000000   Median : 28064226  
                    Mean   :104.9   Mean   : 38165464   Mean   : 51085555  
                    3rd Qu.:121.0   3rd Qu.: 50000000   3rd Qu.: 65517534  
                    Max.   :201.0   Max.   :553632000   Max.   :760505847  
      nota       elenco_principal  
 Min.   :1.600   Length:3875       
 1st Qu.:5.900   Class :character  
 Median :6.600   Mode  :character  
 Mean   :6.462                     
 3rd Qu.:7.200                     
 Max.   :9.300                     
\end{verbatim}

\section*{dplyr: seus primeiros
verbos}\label{dplyr-seus-primeiros-verbos}
\addcontentsline{toc}{section}{dplyr: seus primeiros verbos}

\markright{dplyr: seus primeiros verbos}

\begin{Shaded}
\begin{Highlighting}[]
\CommentTok{\# Filtrar}
\NormalTok{filmes }\SpecialCharTok{\%\textgreater{}\%} \FunctionTok{filter}\NormalTok{(nota }\SpecialCharTok{\textgreater{}=} \FloatTok{8.5}\NormalTok{)}
\end{Highlighting}
\end{Shaded}

\begin{verbatim}
# A tibble: 47 x 10
   codigo_filme titulo  diretor   ano genero duracao orcamento faturamento  nota
          <dbl> <chr>   <chr>   <dbl> <chr>    <dbl>     <dbl>       <dbl> <dbl>
 1            4 The Da~ Christ~  2012 Misté~      99 250000000   448130642   8.5
 2           65 The Da~ Christ~  2008 Avent~     151 185000000   533316061   9  
 3           94 Inters~ Christ~  2014 Roman~      90 165000000   187991439   8.6
 4           95 Incept~ Christ~  2010 Ação       103 160000000   292568851   8.8
 5          261 The Lo~ Peter ~  2001 Anima~     170  93000000   313837577   8.8
 6          273 Gladia~ Ridley~  2000 Drama      159 103000000   187670866   8.5
 7          278 Termin~ James ~  1991 Terror     119 102000000   204843350   8.5
 8          286 Django~ Quenti~  2012 Coméd~      91 100000000   162804648   8.5
 9          329 The Lo~ Peter ~  2003 Coméd~      80  94000000   377019252   8.9
10          330 The Lo~ Peter ~  2002 Anima~     121  94000000   340478898   8.7
# i 37 more rows
# i 1 more variable: elenco_principal <chr>
\end{verbatim}

\begin{Shaded}
\begin{Highlighting}[]
\CommentTok{\# Selecionar}
\NormalTok{filmes }\SpecialCharTok{\%\textgreater{}\%} \FunctionTok{select}\NormalTok{(titulo, ano, nota)}
\end{Highlighting}
\end{Shaded}

\begin{verbatim}
# A tibble: 3,875 x 3
   titulo                                      ano  nota
   <chr>                                     <dbl> <dbl>
 1 Avatar                                     2009   7.9
 2 Pirates of the Caribbean: At World's End   2007   7.1
 3 Spectre                                    2015   6.8
 4 The Dark Knight Rises                      2012   8.5
 5 John Carter                                2012   6.6
 6 Spider-Man 3                               2007   6.2
 7 Tangled                                    2010   7.8
 8 Avengers: Age of Ultron                    2015   7.5
 9 Harry Potter and the Half-Blood Prince     2009   7.5
10 Batman v Superman: Dawn of Justice         2016   6.9
# i 3,865 more rows
\end{verbatim}

\begin{Shaded}
\begin{Highlighting}[]
\CommentTok{\# Criar coluna}
\NormalTok{filmes }\SpecialCharTok{\%\textgreater{}\%} \FunctionTok{mutate}\NormalTok{(}\AttributeTok{lucro =}\NormalTok{ faturamento }\SpecialCharTok{{-}}\NormalTok{ orcamento)}
\end{Highlighting}
\end{Shaded}

\begin{verbatim}
# A tibble: 3,875 x 11
   codigo_filme titulo  diretor   ano genero duracao orcamento faturamento  nota
          <dbl> <chr>   <chr>   <dbl> <chr>    <dbl>     <dbl>       <dbl> <dbl>
 1            1 Avatar  James ~  2009 Ação        81 237000000   760505847   7.9
 2            2 Pirate~ Gore V~  2007 Avent~     126 300000000   309404152   7.1
 3            3 Spectr~ Sam Me~  2015 Coméd~      88 245000000   200074175   6.8
 4            4 The Da~ Christ~  2012 Misté~      99 250000000   448130642   8.5
 5            5 John C~ Andrew~  2012 Fanta~      92 263700000    73058679   6.6
 6            6 Spider~ Sam Ra~  2007 Drama       95 258000000   336530303   6.2
 7            7 Tangle~ Nathan~  2010 Terror      39 260000000   200807262   7.8
 8            8 Avenge~ Joss W~  2015 Misté~     111 250000000   458991599   7.5
 9            9 Harry ~ David ~  2009 Ficçã~      82 250000000   301956980   7.5
10           10 Batman~ Zack S~  2016 Terror      89 250000000   330249062   6.9
# i 3,865 more rows
# i 2 more variables: elenco_principal <chr>, lucro <dbl>
\end{verbatim}

\begin{Shaded}
\begin{Highlighting}[]
\CommentTok{\# Ordenar}
\NormalTok{filmes }\SpecialCharTok{\%\textgreater{}\%} \FunctionTok{arrange}\NormalTok{(}\FunctionTok{desc}\NormalTok{(nota)) }\SpecialCharTok{\%\textgreater{}\%} \FunctionTok{select}\NormalTok{(titulo, nota) }\SpecialCharTok{\%\textgreater{}\%} \FunctionTok{head}\NormalTok{(}\DecValTok{5}\NormalTok{)}
\end{Highlighting}
\end{Shaded}

\begin{verbatim}
# A tibble: 5 x 2
  titulo                                          nota
  <chr>                                          <dbl>
1 The Shawshank Redemption                         9.3
2 The Godfather                                    9.2
3 The Dark Knight                                  9  
4 The Godfather: Part II                           9  
5 The Lord of the Rings: The Return of the King    8.9
\end{verbatim}

\section*{Exercícios}\label{exercuxedcios}
\addcontentsline{toc}{section}{Exercícios}

\markright{Exercícios}

\begin{enumerate}
\def\labelenumi{\arabic{enumi}.}
\tightlist
\item
  Crie um vetor com 5 notas e calcule a média. \textbf{Resultado:}
  \texttt{8.4}\\
\item
  Carregue \texttt{filmes\_imdb.csv}. Quantas linhas e colunas?
  \textbf{Resultado:} 4826 × 11\\
\item
  Filtre filmes com nota ≥ 9.0. Quantos? \textbf{Resultado:} 18
\end{enumerate}

\chapter{Dia 1 --- Manhã
(continuação)}\label{dia-1-manhuxe3-continuauxe7uxe3o}

\textbf{Manhã --- Parte 2 \textbar{} 2 horas}

\section*{\texorpdfstring{\texttt{group\_by()} + \texttt{summarise()}
--- o coração da
análise}{group\_by() + summarise() --- o coração da análise}}\label{group_by-summarise-o-corauxe7uxe3o-da-anuxe1lise}
\addcontentsline{toc}{section}{\texttt{group\_by()} +
\texttt{summarise()} --- o coração da análise}

\markright{\texttt{group\_by()} + \texttt{summarise()} --- o coração da
análise}

\begin{Shaded}
\begin{Highlighting}[]
\CommentTok{\# Nota média por gênero}
\NormalTok{filmes }\SpecialCharTok{\%\textgreater{}\%} 
  \FunctionTok{group\_by}\NormalTok{(genero) }\SpecialCharTok{\%\textgreater{}\%} 
  \FunctionTok{summarise}\NormalTok{(}\AttributeTok{nota\_media =} \FunctionTok{mean}\NormalTok{(nota), }\AttributeTok{n =} \FunctionTok{n}\NormalTok{()) }\SpecialCharTok{\%\textgreater{}\%} 
  \FunctionTok{arrange}\NormalTok{(}\FunctionTok{desc}\NormalTok{(nota\_media))}
\end{Highlighting}
\end{Shaded}

\begin{verbatim}
# A tibble: 19 x 3
   genero            nota_media     n
   <chr>                  <dbl> <int>
 1 Musical                 6.75    39
 2 Esporte                 6.73    16
 3 História                6.56    31
 4 Guerra                  6.55    33
 5 Biografia               6.50   121
 6 Comédia                 6.50   452
 7 Ação                    6.49   566
 8 Drama                   6.48   719
 9 Romance                 6.48   245
10 Documentário            6.47   117
11 Ficção científica       6.47   306
12 Fantasia                6.47    82
13 Mistério                6.46   113
14 Aventura                6.45   204
15 Thriller                6.41   185
16 Terror                  6.40   311
17 Crime                   6.39   193
18 Animação                6.29   114
19 Família                 6.1     28
\end{verbatim}

\begin{Shaded}
\begin{Highlighting}[]
\CommentTok{\# Contagem rápida}
\NormalTok{filmes }\SpecialCharTok{\%\textgreater{}\%} \FunctionTok{count}\NormalTok{(genero, }\AttributeTok{sort =} \ConstantTok{TRUE}\NormalTok{)}
\end{Highlighting}
\end{Shaded}

\begin{verbatim}
# A tibble: 19 x 2
   genero                n
   <chr>             <int>
 1 Drama               719
 2 Ação                566
 3 Comédia             452
 4 Terror              311
 5 Ficção científica   306
 6 Romance             245
 7 Aventura            204
 8 Crime               193
 9 Thriller            185
10 Biografia           121
11 Documentário        117
12 Animação            114
13 Mistério            113
14 Fantasia             82
15 Musical              39
16 Guerra               33
17 História             31
18 Família              28
19 Esporte              16
\end{verbatim}

\begin{Shaded}
\begin{Highlighting}[]
\CommentTok{\# Diretores mais prolíficos}
\NormalTok{filmes }\SpecialCharTok{\%\textgreater{}\%} \FunctionTok{count}\NormalTok{(diretor, }\AttributeTok{sort =} \ConstantTok{TRUE}\NormalTok{) }\SpecialCharTok{\%\textgreater{}\%} \FunctionTok{head}\NormalTok{(}\DecValTok{10}\NormalTok{)}
\end{Highlighting}
\end{Shaded}

\begin{verbatim}
# A tibble: 10 x 2
   diretor               n
   <chr>             <int>
 1 Steven Spielberg     25
 2 Clint Eastwood       19
 3 Woody Allen          19
 4 Ridley Scott         17
 5 Martin Scorsese      16
 6 Steven Soderbergh    16
 7 Tim Burton           16
 8 Renny Harlin         15
 9 Spike Lee            15
10 Barry Levinson       13
\end{verbatim}

\section*{Pipeline completo}\label{pipeline-completo}
\addcontentsline{toc}{section}{Pipeline completo}

\markright{Pipeline completo}

\begin{Shaded}
\begin{Highlighting}[]
\CommentTok{\# Análise completa em um fluxo}
\NormalTok{filmes }\SpecialCharTok{\%\textgreater{}\%} 
  \FunctionTok{filter}\NormalTok{(ano }\SpecialCharTok{\textgreater{}=} \DecValTok{2000}\NormalTok{) }\SpecialCharTok{\%\textgreater{}\%} 
  \FunctionTok{group\_by}\NormalTok{(genero) }\SpecialCharTok{\%\textgreater{}\%} 
  \FunctionTok{summarise}\NormalTok{(}
    \AttributeTok{nota\_media =} \FunctionTok{mean}\NormalTok{(nota),}
    \AttributeTok{orcamento\_medio =} \FunctionTok{mean}\NormalTok{(orcamento, }\AttributeTok{na.rm =} \ConstantTok{TRUE}\NormalTok{),}
    \AttributeTok{n =} \FunctionTok{n}\NormalTok{()}
\NormalTok{  ) }\SpecialCharTok{\%\textgreater{}\%} 
  \FunctionTok{filter}\NormalTok{(n }\SpecialCharTok{\textgreater{}=} \DecValTok{10}\NormalTok{) }\SpecialCharTok{\%\textgreater{}\%} 
  \FunctionTok{arrange}\NormalTok{(}\FunctionTok{desc}\NormalTok{(nota\_media))}
\end{Highlighting}
\end{Shaded}

\begin{verbatim}
# A tibble: 19 x 4
   genero            nota_media orcamento_medio     n
   <chr>                  <dbl>           <dbl> <int>
 1 Musical                 6.70       41836667.    30
 2 Esporte                 6.59       28130769.    13
 3 Biografia               6.53       33198931.    84
 4 História                6.53       41181091.    22
 5 Aventura                6.48       44667589.   141
 6 Drama                   6.47       40643946.   519
 7 Ficção científica       6.44       38671878.   230
 8 Guerra                  6.44       49979167.    24
 9 Comédia                 6.44       45237743.   328
10 Romance                 6.41       42549107.   182
11 Fantasia                6.40       40976230.    61
12 Documentário            6.40       40269591.    88
13 Ação                    6.37       45278978.   406
14 Mistério                6.35       53388389.    85
15 Terror                  6.28       42478816.   228
16 Thriller                6.27       39948227.   142
17 Crime                   6.24       40600126.   141
18 Animação                6.14       42774907.    86
19 Família                 5.93       42855263.    19
\end{verbatim}

\section*{Exercícios}\label{exercuxedcios-1}
\addcontentsline{toc}{section}{Exercícios}

\markright{Exercícios}

\begin{enumerate}
\def\labelenumi{\arabic{enumi}.}
\tightlist
\item
  Nota média por gênero. Qual a maior? \textbf{Resultado:} Documentário
  (6.84)\\
\item
  Crie \texttt{lucro\ =\ faturamento\ -\ orcamento}. Qual o mais
  lucrativo? \textbf{Resultado:} ``The story of cap and trade''\\
\item
  Top 10 diretores com mais filmes. \textbf{Resultado:} Woody Allen
  lidera.
\end{enumerate}

\chapter{Dia 1 --- Tarde}\label{dia-1-tarde}

\textbf{Tarde \textbar{} 4 horas}

\section*{Seus primeiros gráficos 🎬}\label{seus-primeiros-gruxe1ficos}
\addcontentsline{toc}{section}{Seus primeiros gráficos 🎬}

\markright{Seus primeiros gráficos 🎬}

A receita é simples:
\texttt{ggplot(dados,\ aes(x,\ y))\ +\ geom\_tipo()}

\begin{Shaded}
\begin{Highlighting}[]
\CommentTok{\# Histograma — distribuição das notas}
\NormalTok{filmes }\SpecialCharTok{\%\textgreater{}\%} \FunctionTok{ggplot}\NormalTok{(}\FunctionTok{aes}\NormalTok{(}\AttributeTok{x =}\NormalTok{ nota)) }\SpecialCharTok{+} \FunctionTok{geom\_histogram}\NormalTok{(}\AttributeTok{bins =} \DecValTok{30}\NormalTok{, }\AttributeTok{fill =} \StringTok{"\#18BC9C"}\NormalTok{, }\AttributeTok{color =} \StringTok{"white"}\NormalTok{)}
\end{Highlighting}
\end{Shaded}

\pandocbounded{\includegraphics[keepaspectratio]{01-dia/d1-tarde_files/figure-pdf/unnamed-chunk-1-1.pdf}}

\begin{Shaded}
\begin{Highlighting}[]
\CommentTok{\# Dispersão — orçamento × faturamento}
\NormalTok{filmes }\SpecialCharTok{\%\textgreater{}\%} \FunctionTok{ggplot}\NormalTok{(}\FunctionTok{aes}\NormalTok{(}\AttributeTok{x =}\NormalTok{ orcamento, }\AttributeTok{y =}\NormalTok{ faturamento)) }\SpecialCharTok{+} \FunctionTok{geom\_point}\NormalTok{(}\AttributeTok{alpha =} \FloatTok{0.3}\NormalTok{)}
\end{Highlighting}
\end{Shaded}

\pandocbounded{\includegraphics[keepaspectratio]{01-dia/d1-tarde_files/figure-pdf/unnamed-chunk-1-2.pdf}}

\begin{Shaded}
\begin{Highlighting}[]
\CommentTok{\# Barras — top 10 gêneros}
\NormalTok{filmes }\SpecialCharTok{\%\textgreater{}\%} \FunctionTok{count}\NormalTok{(genero, }\AttributeTok{sort =} \ConstantTok{TRUE}\NormalTok{) }\SpecialCharTok{\%\textgreater{}\%} \FunctionTok{head}\NormalTok{(}\DecValTok{10}\NormalTok{) }\SpecialCharTok{\%\textgreater{}\%} 
  \FunctionTok{ggplot}\NormalTok{(}\FunctionTok{aes}\NormalTok{(}\AttributeTok{x =} \FunctionTok{reorder}\NormalTok{(genero, n), }\AttributeTok{y =}\NormalTok{ n)) }\SpecialCharTok{+} \FunctionTok{geom\_col}\NormalTok{(}\AttributeTok{fill =} \StringTok{"\#18BC9C"}\NormalTok{) }\SpecialCharTok{+} \FunctionTok{coord\_flip}\NormalTok{()}
\end{Highlighting}
\end{Shaded}

\pandocbounded{\includegraphics[keepaspectratio]{01-dia/d1-tarde_files/figure-pdf/unnamed-chunk-1-3.pdf}}

\begin{Shaded}
\begin{Highlighting}[]
\CommentTok{\# Boxplot — nota por gênero}
\NormalTok{filmes }\SpecialCharTok{\%\textgreater{}\%} \FunctionTok{filter}\NormalTok{(genero }\SpecialCharTok{\%in\%} \FunctionTok{c}\NormalTok{(}\StringTok{"Ação"}\NormalTok{, }\StringTok{"Drama"}\NormalTok{, }\StringTok{"Comédia"}\NormalTok{)) }\SpecialCharTok{\%\textgreater{}\%} 
  \FunctionTok{ggplot}\NormalTok{(}\FunctionTok{aes}\NormalTok{(}\AttributeTok{x =}\NormalTok{ genero, }\AttributeTok{y =}\NormalTok{ nota, }\AttributeTok{fill =}\NormalTok{ genero)) }\SpecialCharTok{+} \FunctionTok{geom\_boxplot}\NormalTok{() }\SpecialCharTok{+} \FunctionTok{theme}\NormalTok{(}\AttributeTok{legend.position =} \StringTok{"none"}\NormalTok{)}
\end{Highlighting}
\end{Shaded}

\pandocbounded{\includegraphics[keepaspectratio]{01-dia/d1-tarde_files/figure-pdf/unnamed-chunk-1-4.pdf}}

\section*{Com dados florestais 🌲}\label{com-dados-florestais}
\addcontentsline{toc}{section}{Com dados florestais 🌲}

\markright{Com dados florestais 🌲}

\begin{Shaded}
\begin{Highlighting}[]
\CommentTok{\# Distribuição diamétrica}
\NormalTok{inv\_florestal }\SpecialCharTok{\%\textgreater{}\%} \FunctionTok{ggplot}\NormalTok{(}\FunctionTok{aes}\NormalTok{(}\AttributeTok{x =}\NormalTok{ dap)) }\SpecialCharTok{+} \FunctionTok{geom\_histogram}\NormalTok{(}\AttributeTok{bins =} \DecValTok{30}\NormalTok{, }\AttributeTok{fill =} \StringTok{"\#18BC9C"}\NormalTok{)}
\end{Highlighting}
\end{Shaded}

\pandocbounded{\includegraphics[keepaspectratio]{01-dia/d1-tarde_files/figure-pdf/unnamed-chunk-2-1.pdf}}

\begin{Shaded}
\begin{Highlighting}[]
\CommentTok{\# Relação hipsométrica}
\NormalTok{inv\_florestal }\SpecialCharTok{\%\textgreater{}\%} \FunctionTok{ggplot}\NormalTok{(}\FunctionTok{aes}\NormalTok{(}\AttributeTok{x =}\NormalTok{ dap, }\AttributeTok{y =}\NormalTok{ altura, }\AttributeTok{color =}\NormalTok{ especie)) }\SpecialCharTok{+} \FunctionTok{geom\_point}\NormalTok{(}\AttributeTok{alpha =} \FloatTok{0.5}\NormalTok{)}
\end{Highlighting}
\end{Shaded}

\pandocbounded{\includegraphics[keepaspectratio]{01-dia/d1-tarde_files/figure-pdf/unnamed-chunk-2-2.pdf}}

\begin{Shaded}
\begin{Highlighting}[]
\CommentTok{\# Volume por espécie}
\NormalTok{inv\_florestal }\SpecialCharTok{\%\textgreater{}\%} \FunctionTok{ggplot}\NormalTok{(}\FunctionTok{aes}\NormalTok{(}\AttributeTok{x =}\NormalTok{ especie, }\AttributeTok{y =}\NormalTok{ volume, }\AttributeTok{fill =}\NormalTok{ especie)) }\SpecialCharTok{+} \FunctionTok{geom\_boxplot}\NormalTok{()}
\end{Highlighting}
\end{Shaded}

\pandocbounded{\includegraphics[keepaspectratio]{01-dia/d1-tarde_files/figure-pdf/unnamed-chunk-2-3.pdf}}

\section*{Exercícios}\label{exercuxedcios-2}
\addcontentsline{toc}{section}{Exercícios}

\markright{Exercícios}

\begin{enumerate}
\def\labelenumi{\arabic{enumi}.}
\tightlist
\item
  Histograma de \texttt{duracao} com 40 bins. \textbf{Resultado:} Pico
  em 90-100 min.\\
\item
  Gráfico de dispersão \texttt{dap} × \texttt{altura}, cor por
  espécie.\\
\item
  Boxplot de \texttt{volume} por \texttt{especie}. Qual a mediana de
  cada? \textbf{Resultado:} Eucalipto \textasciitilde0.35 m³, Pinus
  \textasciitilde0.20 m³.\\
\item
  Gráfico de barras: volume total por espécie.\\
\item
  Exporte seu melhor gráfico:
  \texttt{ggsave("meu\_grafico.png",\ width\ =\ 8,\ height\ =\ 5,\ dpi\ =\ 300)}
\end{enumerate}

\part{2 dias}

\chapter{R Fundamentos --- 2 dias}\label{r-fundamentos-2-dias}

\chapter*{R Fundamentos}\label{r-fundamentos}
\addcontentsline{toc}{chapter}{R Fundamentos}

\markboth{R Fundamentos}{R Fundamentos}

\textbf{16 horas} · \textbf{5 capítulos}

Domine o fluxo completo de análise de dados em R. Ao final dos dois
dias, você будет capaz de importar, manipular, visualizar e modelar
dados de forma independente.

\section{Estrutura}\label{estrutura-1}

{\def\LTcaptype{none} % do not increment counter
\begin{longtable}[]{@{}
  >{\raggedright\arraybackslash}p{(\linewidth - 4\tabcolsep) * \real{0.1200}}
  >{\raggedright\arraybackslash}p{(\linewidth - 4\tabcolsep) * \real{0.4400}}
  >{\raggedright\arraybackslash}p{(\linewidth - 4\tabcolsep) * \real{0.4400}}@{}}
\toprule\noalign{}
\begin{minipage}[b]{\linewidth}\raggedright
\end{minipage} & \begin{minipage}[b]{\linewidth}\raggedright
Manhã (4h)
\end{minipage} & \begin{minipage}[b]{\linewidth}\raggedright
Tarde (4h)
\end{minipage} \\
\midrule\noalign{}
\endhead
\bottomrule\noalign{}
\endlastfoot
\textbf{Dia 1} & \href{d1-manha.qmd}{Dia 1 --- Manhã} &
\href{d1-tarde.qmd}{Dia 1 --- Tarde} \\
\textbf{Dia 2} & \href{d2-manha.qmd}{Dia 2 --- Manhã} &
\href{d2-tarde.qmd}{Dia 2 --- Tarde} · \href{d2-tarde-parte2.qmd}{Dia 2
--- Tarde (continuação)} \\
\end{longtable}
}

\section{Capítulos}\label{capuxedtulos-1}

\begin{itemize}
\tightlist
\item
  \href{d1-manha.qmd}{Dia 1 --- Manhã} --- Introdução ao R e RStudio
\item
  \href{d1-tarde.qmd}{Dia 1 --- Tarde} --- Manipulação de dados com
  dplyr
\item
  \href{d2-manha.qmd}{Dia 2 --- Manhã} --- Joins e Visualização
\item
  \href{d2-tarde.qmd}{Dia 2 --- Tarde} --- Visualização de dados
  (aprofundamento)
\item
  \href{d2-tarde-parte2.qmd}{Dia 2 --- Tarde (continuação)} ---
  Exercícios integrados
\end{itemize}

\section{Datasets}\label{datasets-1}

\begin{itemize}
\tightlist
\item
  🎬 \texttt{filmes\_imdb.csv} --- filmes (nota, orçamento, faturamento,
  gênero)
\item
  🌲 \texttt{inv\_florestal.csv} --- inventário florestal
\end{itemize}

\section{Pré-requisitos}\label{pruxe9-requisitos-2}

\begin{itemize}
\tightlist
\item
  \href{https://cran.r-project.org/}{R} (4.0+)
\item
  \href{https://posit.co/download/rstudio-desktop/}{RStudio}
\end{itemize}

\section{Repositório}\label{reposituxf3rio-2}

👉
\href{https://github.com/italocegatta/curso_r}{github.com/italocegatta/curso\_r}

\chapter{Dia 1 --- Manhã}\label{dia-1-manhuxe3-1}

\textbf{Dia 1 --- Manhã \textbar{} 4 horas}

Este capítulo condensa o essencial para começar: o que é R, como
instalar, estruturas de dados e primeiros passos no console.

\section*{O que é o R?}\label{o-que-uxe9-o-r}
\addcontentsline{toc}{section}{O que é o R?}

\markright{O que é o R?}

O R é uma linguagem de programação gratuita para computação estatística
e gráficos, mantida pelo R-core team. Com mais de 20.000 pacotes, é uma
das ferramentas mais usadas no mundo para análise de dados.

\textbf{Por que R?} Gratuito, comunidade enorme, scripts reprodutíveis,
gráficos profissionais.

\section*{Instalação}\label{instalauxe7uxe3o}
\addcontentsline{toc}{section}{Instalação}

\markright{Instalação}

\begin{enumerate}
\def\labelenumi{\arabic{enumi}.}
\tightlist
\item
  Instale o \textbf{R} em
  \href{https://cran.r-project.org/}{cran.r-project.org}
\item
  Instale o \textbf{RStudio} em
  \href{https://posit.co/download/rstudio-desktop/}{posit.co/download/rstudio-desktop}
\end{enumerate}

\begin{tcolorbox}[enhanced jigsaw, arc=.35mm, bottomrule=.15mm, bottomtitle=1mm, breakable, colback=white, colbacktitle=quarto-callout-warning-color!10!white, colframe=quarto-callout-warning-color-frame, coltitle=black, left=2mm, leftrule=.75mm, opacityback=0, opacitybacktitle=0.6, rightrule=.15mm, title=\textcolor{quarto-callout-warning-color}{\faExclamationTriangle}\hspace{0.5em}{Warning}, titlerule=0mm, toprule=.15mm, toptitle=1mm]

Instale o R \textbf{primeiro}, depois o RStudio.

\end{tcolorbox}

\section*{Primeiros passos}\label{primeiros-passos}
\addcontentsline{toc}{section}{Primeiros passos}

\markright{Primeiros passos}

\begin{Shaded}
\begin{Highlighting}[]
\CommentTok{\# Calculadora}
\DecValTok{2} \SpecialCharTok{+} \DecValTok{3}
\end{Highlighting}
\end{Shaded}

\begin{verbatim}
[1] 5
\end{verbatim}

\begin{Shaded}
\begin{Highlighting}[]
\FunctionTok{sqrt}\NormalTok{(}\DecValTok{81}\NormalTok{)}
\end{Highlighting}
\end{Shaded}

\begin{verbatim}
[1] 9
\end{verbatim}

\begin{Shaded}
\begin{Highlighting}[]
\CommentTok{\# Atribuição}
\NormalTok{nota }\OtherTok{\textless{}{-}} \FloatTok{8.5}
\NormalTok{titulo }\OtherTok{\textless{}{-}} \StringTok{"O Poderoso Chefão"}

\CommentTok{\# Vetores}
\NormalTok{notas }\OtherTok{\textless{}{-}} \FunctionTok{c}\NormalTok{(}\FloatTok{7.5}\NormalTok{, }\FloatTok{8.0}\NormalTok{, }\FloatTok{9.2}\NormalTok{, }\FloatTok{6.8}\NormalTok{, }\FloatTok{8.7}\NormalTok{)}
\FunctionTok{mean}\NormalTok{(notas)}
\end{Highlighting}
\end{Shaded}

\begin{verbatim}
[1] 8.04
\end{verbatim}

\begin{Shaded}
\begin{Highlighting}[]
\FunctionTok{sum}\NormalTok{(notas)}
\end{Highlighting}
\end{Shaded}

\begin{verbatim}
[1] 40.2
\end{verbatim}

\begin{Shaded}
\begin{Highlighting}[]
\FunctionTok{length}\NormalTok{(notas)}
\end{Highlighting}
\end{Shaded}

\begin{verbatim}
[1] 5
\end{verbatim}

\section*{Estruturas de dados}\label{estruturas-de-dados-1}
\addcontentsline{toc}{section}{Estruturas de dados}

\markright{Estruturas de dados}

\begin{Shaded}
\begin{Highlighting}[]
\CommentTok{\# Fatores}
\NormalTok{genero }\OtherTok{\textless{}{-}} \FunctionTok{factor}\NormalTok{(}\FunctionTok{c}\NormalTok{(}\StringTok{"Ação"}\NormalTok{, }\StringTok{"Drama"}\NormalTok{, }\StringTok{"Comédia"}\NormalTok{, }\StringTok{"Drama"}\NormalTok{, }\StringTok{"Ação"}\NormalTok{))}
\FunctionTok{table}\NormalTok{(genero)}
\end{Highlighting}
\end{Shaded}

\begin{verbatim}
genero
   Ação Comédia   Drama 
      2       1       2 
\end{verbatim}

\begin{Shaded}
\begin{Highlighting}[]
\CommentTok{\# Data frame}
\NormalTok{df }\OtherTok{\textless{}{-}} \FunctionTok{data.frame}\NormalTok{(}
  \AttributeTok{titulo =} \FunctionTok{c}\NormalTok{(}\StringTok{"Matrix"}\NormalTok{, }\StringTok{"Interestelar"}\NormalTok{, }\StringTok{"Parasita"}\NormalTok{),}
  \AttributeTok{ano =} \FunctionTok{c}\NormalTok{(}\DecValTok{1999}\NormalTok{, }\DecValTok{2014}\NormalTok{, }\DecValTok{2019}\NormalTok{),}
  \AttributeTok{nota =} \FunctionTok{c}\NormalTok{(}\FloatTok{8.7}\NormalTok{, }\FloatTok{8.6}\NormalTok{, }\FloatTok{8.5}\NormalTok{)}
\NormalTok{)}
\NormalTok{df}
\end{Highlighting}
\end{Shaded}

\begin{verbatim}
        titulo  ano nota
1       Matrix 1999  8.7
2 Interestelar 2014  8.6
3     Parasita 2019  8.5
\end{verbatim}

\begin{Shaded}
\begin{Highlighting}[]
\CommentTok{\# Inspecionar dados}
\FunctionTok{library}\NormalTok{(dplyr)}
\FunctionTok{library}\NormalTok{(ggplot2)}
\FunctionTok{library}\NormalTok{(readr)}
\FunctionTok{library}\NormalTok{(tidyr)}
\FunctionTok{library}\NormalTok{(stringr)}
\FunctionTok{library}\NormalTok{(forcats)}
\FunctionTok{library}\NormalTok{(purrr)}
\FunctionTok{library}\NormalTok{(tibble)}
\NormalTok{filmes }\OtherTok{\textless{}{-}} \FunctionTok{read\_csv2}\NormalTok{(}\StringTok{"../data/filmes\_imdb.csv"}\NormalTok{)}
\FunctionTok{glimpse}\NormalTok{(filmes)}
\end{Highlighting}
\end{Shaded}

\begin{verbatim}
Rows: 3,875
Columns: 10
$ codigo_filme     <dbl> 1, 2, 3, 4, 5, 6, 7, 8, 9, 10, 11, 12, 13, 14, 15, 16~
$ titulo           <chr> "Avatar ", "Pirates of the Caribbean: At World's End ~
$ diretor          <chr> "James Cameron", "Gore Verbinski", "Sam Mendes", "Chr~
$ ano              <dbl> 2009, 2007, 2015, 2012, 2012, 2007, 2010, 2015, 2009,~
$ genero           <chr> "Ação", "Aventura", "Comédia", "Mistério", "Fantasia"~
$ duracao          <dbl> 81, 126, 88, 99, 92, 95, 39, 111, 82, 89, 77, 102, 98~
$ orcamento        <dbl> 237000000, 300000000, 245000000, 250000000, 263700000~
$ faturamento      <dbl> 760505847, 309404152, 200074175, 448130642, 73058679,~
$ nota             <dbl> 7.9, 7.1, 6.8, 8.5, 6.6, 6.2, 7.8, 7.5, 7.5, 6.9, 6.1~
$ elenco_principal <chr> "Viola Davis", "Cate Blanchett", "Natalie Portman", "~
\end{verbatim}

\begin{Shaded}
\begin{Highlighting}[]
\FunctionTok{head}\NormalTok{(filmes, }\DecValTok{5}\NormalTok{)}
\end{Highlighting}
\end{Shaded}

\begin{verbatim}
# A tibble: 5 x 10
  codigo_filme titulo   diretor   ano genero duracao orcamento faturamento  nota
         <dbl> <chr>    <chr>   <dbl> <chr>    <dbl>     <dbl>       <dbl> <dbl>
1            1 Avatar   James ~  2009 Ação        81 237000000   760505847   7.9
2            2 Pirates~ Gore V~  2007 Avent~     126 300000000   309404152   7.1
3            3 Spectre  Sam Me~  2015 Coméd~      88 245000000   200074175   6.8
4            4 The Dar~ Christ~  2012 Misté~      99 250000000   448130642   8.5
5            5 John Ca~ Andrew~  2012 Fanta~      92 263700000    73058679   6.6
# i 1 more variable: elenco_principal <chr>
\end{verbatim}

\begin{Shaded}
\begin{Highlighting}[]
\FunctionTok{summary}\NormalTok{(filmes)}
\end{Highlighting}
\end{Shaded}

\begin{verbatim}
  codigo_filme       titulo            diretor               ano      
 Min.   :   1.0   Length:3875        Length:3875        Min.   :1920  
 1st Qu.: 971.5   Class :character   Class :character   1st Qu.:1999  
 Median :1942.0   Mode  :character   Mode  :character   Median :2005  
 Mean   :1941.6                                         Mean   :2003  
 3rd Qu.:2911.5                                         3rd Qu.:2010  
 Max.   :3881.0                                         Max.   :2016  
    genero             duracao        orcamento          faturamento       
 Length:3875        Min.   : 31.0   Min.   :      218   Min.   :      162  
 Class :character   1st Qu.: 89.0   1st Qu.: 10000000   1st Qu.:  6856116  
 Mode  :character   Median :105.0   Median : 24000000   Median : 28064226  
                    Mean   :104.9   Mean   : 38165464   Mean   : 51085555  
                    3rd Qu.:121.0   3rd Qu.: 50000000   3rd Qu.: 65517534  
                    Max.   :201.0   Max.   :553632000   Max.   :760505847  
      nota       elenco_principal  
 Min.   :1.600   Length:3875       
 1st Qu.:5.900   Class :character  
 Median :6.600   Mode  :character  
 Mean   :6.462                     
 3rd Qu.:7.200                     
 Max.   :9.300                     
\end{verbatim}

\begin{Shaded}
\begin{Highlighting}[]
\FunctionTok{dim}\NormalTok{(filmes)}
\end{Highlighting}
\end{Shaded}

\begin{verbatim}
[1] 3875   10
\end{verbatim}

\section*{Exercícios}\label{exercuxedcios-3}
\addcontentsline{toc}{section}{Exercícios}

\markright{Exercícios}

\subsection{Exercício 1}\label{exercuxedcio-1}

Crie um vetor com notas de 5 filmes e calcule a média.

\textbf{Resultado esperado:} \texttt{{[}1{]}\ 8.4}

\subsection{Exercício 2}\label{exercuxedcio-2}

Carregue \texttt{filmes\_imdb.csv}. Quantas linhas e colunas? Qual a
nota média?

\textbf{Resultado esperado:} \texttt{{[}1{]}\ 4826\ \ \ 11} e
\texttt{{[}1{]}\ 6.37}

\chapter{Dia 1 --- Tarde}\label{dia-1-tarde-1}

\textbf{Dia 1 --- Tarde \textbar{} 4 horas}

\section*{\texorpdfstring{O pipe:
\texttt{\%\textgreater{}\%}}{O pipe: \%\textgreater\%}}\label{o-pipe}
\addcontentsline{toc}{section}{O pipe: \texttt{\%\textgreater{}\%}}

\markright{O pipe: \texttt{\%\textgreater{}\%}}

Atalho: \texttt{Ctrl/Cmd\ +\ Shift\ +\ M}. Encadeia operações da
esquerda para a direita:

\begin{Shaded}
\begin{Highlighting}[]
\CommentTok{\# Demonstração: o pipe encadeia operações}
\NormalTok{filmes }\SpecialCharTok{\%\textgreater{}\%} 
  \FunctionTok{filter}\NormalTok{(ano }\SpecialCharTok{\textgreater{}} \DecValTok{2000}\NormalTok{) }\SpecialCharTok{\%\textgreater{}\%} 
  \FunctionTok{arrange}\NormalTok{(}\FunctionTok{desc}\NormalTok{(nota))}
\end{Highlighting}
\end{Shaded}

\begin{verbatim}
# A tibble: 2,667 x 10
   codigo_filme titulo  diretor   ano genero duracao orcamento faturamento  nota
          <dbl> <chr>   <chr>   <dbl> <chr>    <dbl>     <dbl>       <dbl> <dbl>
 1           65 The Da~ Christ~  2008 Avent~     151 185000000   533316061   9  
 2          329 The Lo~ Peter ~  2003 Coméd~      80  94000000   377019252   8.9
 3           95 Incept~ Christ~  2010 Ação       103 160000000   292568851   8.8
 4          261 The Lo~ Peter ~  2001 Anima~     170  93000000   313837577   8.8
 5          330 The Lo~ Peter ~  2002 Anima~     121  94000000   340478898   8.7
 6         3398 City o~ Fernan~  2002 Ficçã~      56   3300000     7563397   8.7
 7           94 Inters~ Christ~  2014 Roman~      90 165000000   187991439   8.6
 8         2182 Spirit~ Hayao ~  2001 Roman~      93  19000000    10049886   8.6
 9            4 The Da~ Christ~  2012 Misté~      99 250000000   448130642   8.5
10          286 Django~ Quenti~  2012 Coméd~      91 100000000   162804648   8.5
# i 2,657 more rows
# i 1 more variable: elenco_principal <chr>
\end{verbatim}

\section*{Os 6 verbos do dplyr}\label{os-6-verbos-do-dplyr}
\addcontentsline{toc}{section}{Os 6 verbos do dplyr}

\markright{Os 6 verbos do dplyr}

{\def\LTcaptype{none} % do not increment counter
\begin{longtable}[]{@{}ll@{}}
\toprule\noalign{}
Verbo & O que faz \\
\midrule\noalign{}
\endhead
\bottomrule\noalign{}
\endlastfoot
\texttt{filter()} & Filtra linhas \\
\texttt{select()} & Seleciona colunas \\
\texttt{mutate()} & Cria/modifica colunas \\
\texttt{arrange()} & Ordena \\
\texttt{summarise()} & Resume \\
\texttt{group\_by()} & Agrupa \\
\end{longtable}
}

\subsection{\texorpdfstring{\texttt{filter()}}{filter()}}\label{filter}

\begin{Shaded}
\begin{Highlighting}[]
\CommentTok{\# Filmes com nota ≥ 8.5}
\NormalTok{filmes }\SpecialCharTok{\%\textgreater{}\%} \FunctionTok{filter}\NormalTok{(nota }\SpecialCharTok{\textgreater{}=} \FloatTok{8.5}\NormalTok{)}
\end{Highlighting}
\end{Shaded}

\begin{verbatim}
# A tibble: 47 x 10
   codigo_filme titulo  diretor   ano genero duracao orcamento faturamento  nota
          <dbl> <chr>   <chr>   <dbl> <chr>    <dbl>     <dbl>       <dbl> <dbl>
 1            4 The Da~ Christ~  2012 Misté~      99 250000000   448130642   8.5
 2           65 The Da~ Christ~  2008 Avent~     151 185000000   533316061   9  
 3           94 Inters~ Christ~  2014 Roman~      90 165000000   187991439   8.6
 4           95 Incept~ Christ~  2010 Ação       103 160000000   292568851   8.8
 5          261 The Lo~ Peter ~  2001 Anima~     170  93000000   313837577   8.8
 6          273 Gladia~ Ridley~  2000 Drama      159 103000000   187670866   8.5
 7          278 Termin~ James ~  1991 Terror     119 102000000   204843350   8.5
 8          286 Django~ Quenti~  2012 Coméd~      91 100000000   162804648   8.5
 9          329 The Lo~ Peter ~  2003 Coméd~      80  94000000   377019252   8.9
10          330 The Lo~ Peter ~  2002 Anima~     121  94000000   340478898   8.7
# i 37 more rows
# i 1 more variable: elenco_principal <chr>
\end{verbatim}

\begin{Shaded}
\begin{Highlighting}[]
\CommentTok{\# Ação OU Comédia}
\NormalTok{filmes }\SpecialCharTok{\%\textgreater{}\%} \FunctionTok{filter}\NormalTok{(genero }\SpecialCharTok{\%in\%} \FunctionTok{c}\NormalTok{(}\StringTok{"Ação"}\NormalTok{, }\StringTok{"Comédia"}\NormalTok{))}
\end{Highlighting}
\end{Shaded}

\begin{verbatim}
# A tibble: 1,018 x 10
   codigo_filme titulo  diretor   ano genero duracao orcamento faturamento  nota
          <dbl> <chr>   <chr>   <dbl> <chr>    <dbl>     <dbl>       <dbl> <dbl>
 1            1 Avatar  James ~  2009 Ação        81 237000000   760505847   7.9
 2            3 Spectr~ Sam Me~  2015 Coméd~      88 245000000   200074175   6.8
 3           17 The Av~ Joss W~  2012 Ação        53 220000000   623279547   8.1
 4           19 Men in~ Barry ~  2012 Ação        57 225000000   179020854   6.8
 5           29 Jurass~ Colin ~  2015 Ação        58 150000000   652177271   7  
 6           38 Oz the~ Sam Ra~  2013 Ação        95 215000000   234903076   6.4
 7           39 The Am~ Marc W~  2014 Ação        88 200000000   202853933   6.7
 8           40 TRON: ~ Joseph~  2010 Ação        82 170000000   172051787   6.8
 9           42 Green ~ Martin~  2011 Coméd~     117 200000000   116593191   5.6
10           43 Toy St~ Lee Un~  2010 Coméd~      48 200000000   414984497   8.3
# i 1,008 more rows
# i 1 more variable: elenco_principal <chr>
\end{verbatim}

\subsection{\texorpdfstring{\texttt{select()} e
\texttt{mutate()}}{select() e mutate()}}\label{select-e-mutate}

\begin{Shaded}
\begin{Highlighting}[]
\NormalTok{filmes }\SpecialCharTok{\%\textgreater{}\%} \FunctionTok{select}\NormalTok{(titulo, ano, nota)}
\end{Highlighting}
\end{Shaded}

\begin{verbatim}
# A tibble: 3,875 x 3
   titulo                                      ano  nota
   <chr>                                     <dbl> <dbl>
 1 Avatar                                     2009   7.9
 2 Pirates of the Caribbean: At World's End   2007   7.1
 3 Spectre                                    2015   6.8
 4 The Dark Knight Rises                      2012   8.5
 5 John Carter                                2012   6.6
 6 Spider-Man 3                               2007   6.2
 7 Tangled                                    2010   7.8
 8 Avengers: Age of Ultron                    2015   7.5
 9 Harry Potter and the Half-Blood Prince     2009   7.5
10 Batman v Superman: Dawn of Justice         2016   6.9
# i 3,865 more rows
\end{verbatim}

\begin{Shaded}
\begin{Highlighting}[]
\NormalTok{filmes }\SpecialCharTok{\%\textgreater{}\%} 
  \FunctionTok{mutate}\NormalTok{(}
    \AttributeTok{lucro =}\NormalTok{ faturamento }\SpecialCharTok{{-}}\NormalTok{ orcamento,}
    \AttributeTok{lucro\_milhoes =}\NormalTok{ lucro }\SpecialCharTok{/} \DecValTok{1000000}
\NormalTok{  ) }\SpecialCharTok{\%\textgreater{}\%} 
  \FunctionTok{arrange}\NormalTok{(}\FunctionTok{desc}\NormalTok{(lucro)) }\SpecialCharTok{\%\textgreater{}\%} 
  \FunctionTok{select}\NormalTok{(titulo, lucro\_milhoes) }\SpecialCharTok{\%\textgreater{}\%} 
  \FunctionTok{head}\NormalTok{(}\DecValTok{5}\NormalTok{)}
\end{Highlighting}
\end{Shaded}

\begin{verbatim}
# A tibble: 5 x 2
  titulo                              lucro_milhoes
  <chr>                                       <dbl>
1 Avatar                                       524.
2 Jurassic World                               502.
3 Titanic                                      459.
4 Star Wars: Episode IV - A New Hope           450.
5 E.T. the Extra-Terrestrial                   424.
\end{verbatim}

\subsection{\texorpdfstring{\texttt{group\_by()} +
\texttt{summarise()}}{group\_by() + summarise()}}\label{group_by-summarise}

\begin{Shaded}
\begin{Highlighting}[]
\CommentTok{\# Nota média por gênero}
\NormalTok{filmes }\SpecialCharTok{\%\textgreater{}\%} 
  \FunctionTok{group\_by}\NormalTok{(genero) }\SpecialCharTok{\%\textgreater{}\%} 
  \FunctionTok{summarise}\NormalTok{(}\AttributeTok{nota\_media =} \FunctionTok{mean}\NormalTok{(nota), }\AttributeTok{n =} \FunctionTok{n}\NormalTok{()) }\SpecialCharTok{\%\textgreater{}\%} 
  \FunctionTok{arrange}\NormalTok{(}\FunctionTok{desc}\NormalTok{(nota\_media))}
\end{Highlighting}
\end{Shaded}

\begin{verbatim}
# A tibble: 19 x 3
   genero            nota_media     n
   <chr>                  <dbl> <int>
 1 Musical                 6.75    39
 2 Esporte                 6.73    16
 3 História                6.56    31
 4 Guerra                  6.55    33
 5 Biografia               6.50   121
 6 Comédia                 6.50   452
 7 Ação                    6.49   566
 8 Drama                   6.48   719
 9 Romance                 6.48   245
10 Documentário            6.47   117
11 Ficção científica       6.47   306
12 Fantasia                6.47    82
13 Mistério                6.46   113
14 Aventura                6.45   204
15 Thriller                6.41   185
16 Terror                  6.40   311
17 Crime                   6.39   193
18 Animação                6.29   114
19 Família                 6.1     28
\end{verbatim}

\subsection{\texorpdfstring{\texttt{count()}}{count()}}\label{count}

\begin{Shaded}
\begin{Highlighting}[]
\NormalTok{filmes }\SpecialCharTok{\%\textgreater{}\%} \FunctionTok{count}\NormalTok{(genero, }\AttributeTok{sort =} \ConstantTok{TRUE}\NormalTok{)}
\end{Highlighting}
\end{Shaded}

\begin{verbatim}
# A tibble: 19 x 2
   genero                n
   <chr>             <int>
 1 Drama               719
 2 Ação                566
 3 Comédia             452
 4 Terror              311
 5 Ficção científica   306
 6 Romance             245
 7 Aventura            204
 8 Crime               193
 9 Thriller            185
10 Biografia           121
11 Documentário        117
12 Animação            114
13 Mistério            113
14 Fantasia             82
15 Musical              39
16 Guerra               33
17 História             31
18 Família              28
19 Esporte              16
\end{verbatim}

\section*{🎬 Primeiros gráficos}\label{primeiros-gruxe1ficos}
\addcontentsline{toc}{section}{🎬 Primeiros gráficos}

\markright{🎬 Primeiros gráficos}

\begin{Shaded}
\begin{Highlighting}[]
\CommentTok{\# Histograma}
\NormalTok{filmes }\SpecialCharTok{\%\textgreater{}\%} \FunctionTok{ggplot}\NormalTok{(}\FunctionTok{aes}\NormalTok{(}\AttributeTok{x =}\NormalTok{ nota)) }\SpecialCharTok{+} \FunctionTok{geom\_histogram}\NormalTok{(}\AttributeTok{bins =} \DecValTok{30}\NormalTok{, }\AttributeTok{fill =} \StringTok{"\#18BC9C"}\NormalTok{)}
\end{Highlighting}
\end{Shaded}

\pandocbounded{\includegraphics[keepaspectratio]{02-dias/d1-tarde_files/figure-pdf/unnamed-chunk-6-1.pdf}}

\begin{Shaded}
\begin{Highlighting}[]
\CommentTok{\# Dispersão}
\NormalTok{filmes }\SpecialCharTok{\%\textgreater{}\%} \FunctionTok{ggplot}\NormalTok{(}\FunctionTok{aes}\NormalTok{(}\AttributeTok{x =}\NormalTok{ orcamento, }\AttributeTok{y =}\NormalTok{ faturamento)) }\SpecialCharTok{+} \FunctionTok{geom\_point}\NormalTok{(}\AttributeTok{alpha =} \FloatTok{0.3}\NormalTok{)}
\end{Highlighting}
\end{Shaded}

\pandocbounded{\includegraphics[keepaspectratio]{02-dias/d1-tarde_files/figure-pdf/unnamed-chunk-6-2.pdf}}

\begin{Shaded}
\begin{Highlighting}[]
\CommentTok{\# Barras — top 10 gêneros}
\NormalTok{filmes }\SpecialCharTok{\%\textgreater{}\%} \FunctionTok{count}\NormalTok{(genero, }\AttributeTok{sort =} \ConstantTok{TRUE}\NormalTok{) }\SpecialCharTok{\%\textgreater{}\%} \FunctionTok{head}\NormalTok{(}\DecValTok{10}\NormalTok{) }\SpecialCharTok{\%\textgreater{}\%} 
  \FunctionTok{ggplot}\NormalTok{(}\FunctionTok{aes}\NormalTok{(}\AttributeTok{x =} \FunctionTok{reorder}\NormalTok{(genero, n), }\AttributeTok{y =}\NormalTok{ n)) }\SpecialCharTok{+} \FunctionTok{geom\_col}\NormalTok{(}\AttributeTok{fill =} \StringTok{"\#18BC9C"}\NormalTok{) }\SpecialCharTok{+} \FunctionTok{coord\_flip}\NormalTok{()}
\end{Highlighting}
\end{Shaded}

\pandocbounded{\includegraphics[keepaspectratio]{02-dias/d1-tarde_files/figure-pdf/unnamed-chunk-6-3.pdf}}

\begin{Shaded}
\begin{Highlighting}[]
\CommentTok{\# Boxplot}
\NormalTok{filmes }\SpecialCharTok{\%\textgreater{}\%} \FunctionTok{filter}\NormalTok{(genero }\SpecialCharTok{\%in\%} \FunctionTok{c}\NormalTok{(}\StringTok{"Ação"}\NormalTok{, }\StringTok{"Drama"}\NormalTok{, }\StringTok{"Comédia"}\NormalTok{)) }\SpecialCharTok{\%\textgreater{}\%} 
  \FunctionTok{ggplot}\NormalTok{(}\FunctionTok{aes}\NormalTok{(}\AttributeTok{x =}\NormalTok{ genero, }\AttributeTok{y =}\NormalTok{ nota, }\AttributeTok{fill =}\NormalTok{ genero)) }\SpecialCharTok{+} \FunctionTok{geom\_boxplot}\NormalTok{()}
\end{Highlighting}
\end{Shaded}

\pandocbounded{\includegraphics[keepaspectratio]{02-dias/d1-tarde_files/figure-pdf/unnamed-chunk-6-4.pdf}}

\section*{Exercícios}\label{exercuxedcios-4}
\addcontentsline{toc}{section}{Exercícios}

\markright{Exercícios}

\subsection{Exercício 1}\label{exercuxedcio-1-1}

Filtre filmes com nota ≥ 8.5. Quantos são?\\
\textbf{Resultado esperado:} 103 filmes.

\subsection{Exercício 2}\label{exercuxedcio-2-1}

Crie \texttt{lucro\ =\ faturamento\ -\ orcamento}. Qual o filme mais
lucrativo?\\
\textbf{Resultado esperado:} \texttt{"The\ story\ of\ cap\ and\ trade"}
com lucro de 48.7 milhões.

\subsection{Exercício 3}\label{exercuxedcio-3}

Nota média por gênero, ordenada da maior para a menor.\\
\textbf{Resultado esperado:} Documentário (6.84), Guerra (6.75),
História (6.73)\ldots{}

\subsection{Exercício 4}\label{exercuxedcio-4}

Crie um histograma de \texttt{duracao} com 40 bins.\\
\textbf{Resultado esperado:} Pico em 90-100 min, assimétrica à direita.

\chapter{Dia 2 --- Manhã}\label{dia-2-manhuxe3}

\textbf{Dia 2 --- Manhã \textbar{} 4 horas}

\section*{Joins: cruzando tabelas}\label{joins-cruzando-tabelas}
\addcontentsline{toc}{section}{Joins: cruzando tabelas}

\markright{Joins: cruzando tabelas}

\begin{Shaded}
\begin{Highlighting}[]
\CommentTok{\# left\_join: adiciona colunas da tabela direita à esquerda}
\NormalTok{filmes\_com\_elenco }\OtherTok{\textless{}{-}}\NormalTok{ filmes }\SpecialCharTok{\%\textgreater{}\%} \FunctionTok{left\_join}\NormalTok{(elenco, }\AttributeTok{by =} \FunctionTok{c}\NormalTok{(}\StringTok{"codigo\_filme"} \OtherTok{=} \StringTok{"codigo"}\NormalTok{))}

\CommentTok{\# Anti\_join: o que NÃO tem correspondência}
\NormalTok{inv\_florestal }\SpecialCharTok{\%\textgreater{}\%} \FunctionTok{anti\_join}\NormalTok{(parcelas, }\AttributeTok{by =} \StringTok{"parcela"}\NormalTok{)}
\end{Highlighting}
\end{Shaded}

\begin{verbatim}
# A tibble: 0 x 6
# i 6 variables: parcela <dbl>, arvore <dbl>, especie <chr>, dap <dbl>,
#   altura <dbl>, volume <dbl>
\end{verbatim}

\section*{Visualização com
ggplot2}\label{visualizauxe7uxe3o-com-ggplot2}
\addcontentsline{toc}{section}{Visualização com ggplot2}

\markright{Visualização com ggplot2}

\begin{Shaded}
\begin{Highlighting}[]
\CommentTok{\# Dispersão + tendência}
\NormalTok{inv\_florestal }\SpecialCharTok{\%\textgreater{}\%} 
  \FunctionTok{ggplot}\NormalTok{(}\FunctionTok{aes}\NormalTok{(}\AttributeTok{x =}\NormalTok{ dap, }\AttributeTok{y =}\NormalTok{ altura, }\AttributeTok{color =}\NormalTok{ especie)) }\SpecialCharTok{+}
  \FunctionTok{geom\_point}\NormalTok{(}\AttributeTok{alpha =} \FloatTok{0.5}\NormalTok{) }\SpecialCharTok{+}
  \FunctionTok{geom\_smooth}\NormalTok{(}\AttributeTok{method =} \StringTok{"lm"}\NormalTok{, }\AttributeTok{se =} \ConstantTok{FALSE}\NormalTok{) }\SpecialCharTok{+}
  \FunctionTok{scale\_color\_brewer}\NormalTok{(}\AttributeTok{palette =} \StringTok{"Set1"}\NormalTok{) }\SpecialCharTok{+}
  \FunctionTok{labs}\NormalTok{(}\AttributeTok{title =} \StringTok{"Relação hipsométrica"}\NormalTok{, }\AttributeTok{x =} \StringTok{"DAP (cm)"}\NormalTok{, }\AttributeTok{y =} \StringTok{"Altura (m)"}\NormalTok{) }\SpecialCharTok{+}
  \FunctionTok{theme\_minimal}\NormalTok{()}
\end{Highlighting}
\end{Shaded}

\pandocbounded{\includegraphics[keepaspectratio]{02-dias/d2-manha_files/figure-pdf/unnamed-chunk-2-1.pdf}}

\begin{Shaded}
\begin{Highlighting}[]
\CommentTok{\# Facetas}
\NormalTok{inv\_florestal }\SpecialCharTok{\%\textgreater{}\%} 
  \FunctionTok{ggplot}\NormalTok{(}\FunctionTok{aes}\NormalTok{(}\AttributeTok{x =}\NormalTok{ dap, }\AttributeTok{y =}\NormalTok{ altura)) }\SpecialCharTok{+}
  \FunctionTok{geom\_point}\NormalTok{(}\AttributeTok{alpha =} \FloatTok{0.5}\NormalTok{) }\SpecialCharTok{+}
  \FunctionTok{geom\_smooth}\NormalTok{(}\AttributeTok{method =} \StringTok{"lm"}\NormalTok{, }\AttributeTok{formula =}\NormalTok{ y }\SpecialCharTok{\textasciitilde{}} \FunctionTok{log}\NormalTok{(x), }\AttributeTok{se =} \ConstantTok{FALSE}\NormalTok{) }\SpecialCharTok{+}
  \FunctionTok{facet\_wrap}\NormalTok{(}\SpecialCharTok{\textasciitilde{}}\NormalTok{especie)}
\end{Highlighting}
\end{Shaded}

\pandocbounded{\includegraphics[keepaspectratio]{02-dias/d2-manha_files/figure-pdf/unnamed-chunk-2-2.pdf}}

\begin{Shaded}
\begin{Highlighting}[]
\CommentTok{\# Boxplot}
\NormalTok{inv\_florestal }\SpecialCharTok{\%\textgreater{}\%} 
  \FunctionTok{ggplot}\NormalTok{(}\FunctionTok{aes}\NormalTok{(}\AttributeTok{x =}\NormalTok{ especie, }\AttributeTok{y =}\NormalTok{ volume, }\AttributeTok{fill =}\NormalTok{ especie)) }\SpecialCharTok{+}
  \FunctionTok{geom\_boxplot}\NormalTok{() }\SpecialCharTok{+} \FunctionTok{theme}\NormalTok{(}\AttributeTok{legend.position =} \StringTok{"none"}\NormalTok{)}
\end{Highlighting}
\end{Shaded}

\pandocbounded{\includegraphics[keepaspectratio]{02-dias/d2-manha_files/figure-pdf/unnamed-chunk-2-3.pdf}}

\begin{Shaded}
\begin{Highlighting}[]
\CommentTok{\# Barras — volume por parcela}
\NormalTok{inv\_florestal }\SpecialCharTok{\%\textgreater{}\%} \FunctionTok{group\_by}\NormalTok{(parcela) }\SpecialCharTok{\%\textgreater{}\%} \FunctionTok{summarise}\NormalTok{(}\AttributeTok{v =} \FunctionTok{sum}\NormalTok{(volume)) }\SpecialCharTok{\%\textgreater{}\%} 
  \FunctionTok{ggplot}\NormalTok{(}\FunctionTok{aes}\NormalTok{(}\AttributeTok{x =} \FunctionTok{reorder}\NormalTok{(}\FunctionTok{factor}\NormalTok{(parcela), v), }\AttributeTok{y =}\NormalTok{ v)) }\SpecialCharTok{+}
  \FunctionTok{geom\_col}\NormalTok{(}\AttributeTok{fill =} \StringTok{"\#18BC9C"}\NormalTok{) }\SpecialCharTok{+} \FunctionTok{coord\_flip}\NormalTok{() }\SpecialCharTok{+} \FunctionTok{labs}\NormalTok{(}\AttributeTok{x =} \StringTok{"Parcela"}\NormalTok{, }\AttributeTok{y =} \StringTok{"Volume (m³)"}\NormalTok{)}
\end{Highlighting}
\end{Shaded}

\pandocbounded{\includegraphics[keepaspectratio]{02-dias/d2-manha_files/figure-pdf/unnamed-chunk-2-4.pdf}}

\section*{Exercícios}\label{exercuxedcios-5}
\addcontentsline{toc}{section}{Exercícios}

\markright{Exercícios}

\subsection{Exercício 1}\label{exercuxedcio-1-2}

Faça \texttt{left\_join} entre \texttt{inv\_florestal} e
\texttt{parcelas}. Calcule volume/ha e encontre a parcela mais
produtiva.\\
\textbf{Resultado esperado:} Parcela 7 com 420 m³/ha.

\subsection{Exercício 2}\label{exercuxedcio-2-2}

Crie boxplot de \texttt{dap} por \texttt{especie} com cores.\\
\textbf{Resultado esperado:} Eucalipto com DAP mediano \textasciitilde20
cm, Pinus \textasciitilde17 cm.

\subsection{Exercício 3}\label{exercuxedcio-3-1}

Plote \texttt{dap} × \texttt{volume} com \texttt{geom\_point()} +
\texttt{geom\_smooth()}, cor por espécie.\\
\textbf{Resultado esperado:} Curva exponencial --- volume cresce com
DAP².

\chapter{Dia 2 --- Tarde}\label{dia-2-tarde}

\textbf{Dia 2 --- Tarde (parte 1) \textbar{} 2 horas}

\section*{Customização e
exportação}\label{customizauxe7uxe3o-e-exportauxe7uxe3o}
\addcontentsline{toc}{section}{Customização e exportação}

\markright{Customização e exportação}

\begin{Shaded}
\begin{Highlighting}[]
\CommentTok{\# Gráfico completo para relatório}
\NormalTok{inv\_florestal }\SpecialCharTok{\%\textgreater{}\%} 
  \FunctionTok{ggplot}\NormalTok{(}\FunctionTok{aes}\NormalTok{(}\AttributeTok{x =}\NormalTok{ dap, }\AttributeTok{y =}\NormalTok{ altura, }\AttributeTok{color =}\NormalTok{ especie)) }\SpecialCharTok{+}
  \FunctionTok{geom\_point}\NormalTok{(}\AttributeTok{alpha =} \FloatTok{0.5}\NormalTok{, }\AttributeTok{size =} \DecValTok{2}\NormalTok{) }\SpecialCharTok{+}
  \FunctionTok{geom\_smooth}\NormalTok{(}\AttributeTok{method =} \StringTok{"lm"}\NormalTok{, }\AttributeTok{formula =}\NormalTok{ y }\SpecialCharTok{\textasciitilde{}} \FunctionTok{log}\NormalTok{(x), }\AttributeTok{se =} \ConstantTok{TRUE}\NormalTok{, }\AttributeTok{alpha =} \FloatTok{0.2}\NormalTok{) }\SpecialCharTok{+}
  \FunctionTok{scale\_color\_brewer}\NormalTok{(}\AttributeTok{palette =} \StringTok{"Set1"}\NormalTok{) }\SpecialCharTok{+}
  \FunctionTok{labs}\NormalTok{(}\AttributeTok{title =} \StringTok{"Relação hipsométrica — Inventário Florestal"}\NormalTok{,}
       \AttributeTok{subtitle =} \StringTok{"Curvas ajustadas por espécie (IC 95\%)"}\NormalTok{,}
       \AttributeTok{x =} \StringTok{"DAP (cm)"}\NormalTok{, }\AttributeTok{y =} \StringTok{"Altura total (m)"}\NormalTok{, }\AttributeTok{color =} \StringTok{"Espécie"}\NormalTok{,}
       \AttributeTok{caption =} \StringTok{"Fonte: Inventário florestal"}\NormalTok{) }\SpecialCharTok{+}
  \FunctionTok{theme\_minimal}\NormalTok{(}\AttributeTok{base\_size =} \DecValTok{13}\NormalTok{) }\SpecialCharTok{+}
  \FunctionTok{theme}\NormalTok{(}\AttributeTok{plot.title =} \FunctionTok{element\_text}\NormalTok{(}\AttributeTok{face =} \StringTok{"bold"}\NormalTok{), }\AttributeTok{legend.position =} \StringTok{"bottom"}\NormalTok{)}
\end{Highlighting}
\end{Shaded}

\pandocbounded{\includegraphics[keepaspectratio]{02-dias/d2-tarde_files/figure-pdf/unnamed-chunk-1-1.pdf}}

\begin{Shaded}
\begin{Highlighting}[]
\CommentTok{\# Dashboard com patchwork}
\NormalTok{p1 }\OtherTok{\textless{}{-}}\NormalTok{ inv\_florestal }\SpecialCharTok{\%\textgreater{}\%} \FunctionTok{ggplot}\NormalTok{(}\FunctionTok{aes}\NormalTok{(}\AttributeTok{x =}\NormalTok{ dap)) }\SpecialCharTok{+} \FunctionTok{geom\_histogram}\NormalTok{(}\AttributeTok{bins =} \DecValTok{30}\NormalTok{, }\AttributeTok{fill =} \StringTok{"\#18BC9C"}\NormalTok{) }\SpecialCharTok{+} \FunctionTok{labs}\NormalTok{(}\AttributeTok{title =} \StringTok{"Dist. diamétrica"}\NormalTok{)}
\NormalTok{p2 }\OtherTok{\textless{}{-}}\NormalTok{ inv\_florestal }\SpecialCharTok{\%\textgreater{}\%} \FunctionTok{ggplot}\NormalTok{(}\FunctionTok{aes}\NormalTok{(}\AttributeTok{x =}\NormalTok{ dap, }\AttributeTok{y =}\NormalTok{ altura, }\AttributeTok{color =}\NormalTok{ especie)) }\SpecialCharTok{+} \FunctionTok{geom\_point}\NormalTok{(}\AttributeTok{alpha =} \FloatTok{0.5}\NormalTok{) }\SpecialCharTok{+} \FunctionTok{labs}\NormalTok{(}\AttributeTok{title =} \StringTok{"Hipsométrica"}\NormalTok{)}
\NormalTok{p3 }\OtherTok{\textless{}{-}}\NormalTok{ inv\_florestal }\SpecialCharTok{\%\textgreater{}\%} \FunctionTok{ggplot}\NormalTok{(}\FunctionTok{aes}\NormalTok{(}\AttributeTok{x =}\NormalTok{ especie, }\AttributeTok{y =}\NormalTok{ volume, }\AttributeTok{fill =}\NormalTok{ especie)) }\SpecialCharTok{+} \FunctionTok{geom\_boxplot}\NormalTok{() }\SpecialCharTok{+} \FunctionTok{theme}\NormalTok{(}\AttributeTok{legend.position =} \StringTok{"none"}\NormalTok{)}
\NormalTok{p4 }\OtherTok{\textless{}{-}}\NormalTok{ inv\_florestal }\SpecialCharTok{\%\textgreater{}\%} \FunctionTok{group\_by}\NormalTok{(parcela) }\SpecialCharTok{\%\textgreater{}\%} \FunctionTok{summarise}\NormalTok{(}\AttributeTok{v =} \FunctionTok{sum}\NormalTok{(volume)) }\SpecialCharTok{\%\textgreater{}\%} 
  \FunctionTok{ggplot}\NormalTok{(}\FunctionTok{aes}\NormalTok{(}\AttributeTok{x =} \FunctionTok{factor}\NormalTok{(parcela), }\AttributeTok{y =}\NormalTok{ v)) }\SpecialCharTok{+} \FunctionTok{geom\_col}\NormalTok{(}\AttributeTok{fill =} \StringTok{"\#2C3E50"}\NormalTok{) }\SpecialCharTok{+} \FunctionTok{labs}\NormalTok{(}\AttributeTok{title =} \StringTok{"Volume/parcela"}\NormalTok{, }\AttributeTok{x =} \StringTok{"Parcela"}\NormalTok{)}

\NormalTok{(p1 }\SpecialCharTok{+}\NormalTok{ p2) }\SpecialCharTok{/}\NormalTok{ (p3 }\SpecialCharTok{+}\NormalTok{ p4)}
\end{Highlighting}
\end{Shaded}

\pandocbounded{\includegraphics[keepaspectratio]{02-dias/d2-tarde_files/figure-pdf/unnamed-chunk-2-1.pdf}}

\begin{Shaded}
\begin{Highlighting}[]
\CommentTok{\# Exportar}
\CommentTok{\# ggsave("dashboard.png", width = 10, height = 7, dpi = 300)}
\end{Highlighting}
\end{Shaded}

\section*{Exercícios}\label{exercuxedcios-6}
\addcontentsline{toc}{section}{Exercícios}

\markright{Exercícios}

\begin{enumerate}
\def\labelenumi{\arabic{enumi}.}
\tightlist
\item
  Crie histograma do DAP com 30 bins, título e rótulos.
\item
  Boxplot de volume por espécie. Qual tem maior variabilidade?
\item
  Gráfico de barras do volume total por parcela (top 10), com
  \texttt{coord\_flip()}.
\item
  Exporte seu melhor gráfico como PNG (8×5 pol, 300 DPI).
\end{enumerate}

\chapter{Dia 2 --- Tarde
(continuação)}\label{dia-2-tarde-continuauxe7uxe3o}

\textbf{Dia 2 --- Tarde (parte 2) \textbar{} 2 horas}

Este capítulo consolida o aprendizado do curso de 2 dias com exercícios
práticos que integram todas as técnicas.

\section*{Instruções}\label{instruuxe7uxf5es}
\addcontentsline{toc}{section}{Instruções}

\markright{Instruções}

Cada exercício mostra o \textbf{resultado esperado}. Escreva o código e
confira se sua solução produz o mesmo resultado.

\begin{center}\rule{0.5\linewidth}{0.5pt}\end{center}

\section{Exercício 1 --- Explorando o
dataset}\label{exercuxedcio-1-explorando-o-dataset}

Carregue \texttt{inv\_florestal.csv} e \texttt{parcelas.csv}. Faça o
join entre eles e responda:

\textbf{a)} Quantas árvores no total? E quantas parcelas?\\
\textbf{Resultado esperado:} 992 árvores, 30 parcelas.

\textbf{b)} Qual o DAP médio geral? E por espécie?\\
\textbf{Resultado esperado:} Eucalipto: 20.6 cm, Pinus: 17.8 cm.

\textbf{c)} Qual parcela tem o maior volume por hectare?\\
\textbf{Resultado esperado:} Parcela 7 com 420 m³/ha.

\begin{center}\rule{0.5\linewidth}{0.5pt}\end{center}

\section{Exercício 2 ---
Visualização}\label{exercuxedcio-2-visualizauxe7uxe3o}

\textbf{a)} Crie um histograma da distribuição diamétrica.\\
\textbf{Resultado esperado:} Pico em 15-20 cm, assimétrica à direita.

\textbf{b)} Crie um gráfico de dispersão \texttt{dap} × \texttt{altura}
com cor por espécie e curva de tendência.\\
\textbf{Resultado esperado:} Duas curvas quase paralelas, Eucalipto
ligeiramente acima.

\textbf{c)} Crie um dashboard com 4 gráficos: histograma do DAP, relação
hipsométrica, boxplot de volume por espécie e barras de volume por
parcela.\\
\textbf{Resultado esperado:} Painel 2×2 com tema consistente.

\begin{center}\rule{0.5\linewidth}{0.5pt}\end{center}

\section{Exercício 3 --- Análise por gênero
(filmes)}\label{exercuxedcio-3-anuxe1lise-por-guxeanero-filmes}

Usando \texttt{filmes\_imdb.csv}:

\textbf{a)} Gênero com maior nota média (mínimo 50 filmes)?\\
\textbf{Resultado esperado:} Documentário com 6.84.

\textbf{b)} Gráfico de dispersão \texttt{orcamento} ×
\texttt{faturamento}, com cor por gênero.\\
\textbf{Resultado esperado:} Correlação positiva, Ação e Aventura com
maiores valores.

\textbf{c)} Qual diretor tem mais filmes no dataset?\\
\textbf{Resultado esperado:} Woody Allen.

\begin{center}\rule{0.5\linewidth}{0.5pt}\end{center}

\section{Exercício 4 --- Desafio
final}\label{exercuxedcio-4-desafio-final}

Usando \texttt{inv\_florestal}, faça uma análise completa:

\begin{enumerate}
\def\labelenumi{\arabic{enumi}.}
\tightlist
\item
  Calcule área basal por árvore: \texttt{ab\ =\ pi\ *\ (dap/200)\^{}2}
\item
  Calcule volume: \texttt{vol\ =\ ab\ *\ altura\ *\ 0.5}
\item
  Agrupe por parcela e calcule volume/ha (área da parcela = 600 m²)
\item
  Crie um gráfico de barras do volume/ha por parcela (top 10)
\item
  Ajuste
  \texttt{lm(log(volume)\ \textasciitilde{}\ log(dap)\ +\ log(altura))}
  e reporte o R²
\end{enumerate}

\textbf{Resultado esperado do R²:} 0.96

\part{4 dias}

\chapter{R Completo --- 4 dias}\label{r-completo-4-dias}

\chapter*{R Completo}\label{r-completo}
\addcontentsline{toc}{chapter}{R Completo}

\markboth{R Completo}{R Completo}

\textbf{32 horas} · \textbf{8 capítulos}

Domine R para análise de dados --- da importação à modelagem
estatística, produzindo relatórios técnicos completos.

\section{Estrutura}\label{estrutura-2}

{\def\LTcaptype{none} % do not increment counter
\begin{longtable}[]{@{}
  >{\centering\arraybackslash}p{(\linewidth - 4\tabcolsep) * \real{0.1724}}
  >{\raggedright\arraybackslash}p{(\linewidth - 4\tabcolsep) * \real{0.4138}}
  >{\raggedright\arraybackslash}p{(\linewidth - 4\tabcolsep) * \real{0.4138}}@{}}
\toprule\noalign{}
\begin{minipage}[b]{\linewidth}\centering
\end{minipage} & \begin{minipage}[b]{\linewidth}\raggedright
Manhã (4h)
\end{minipage} & \begin{minipage}[b]{\linewidth}\raggedright
Tarde (4h)
\end{minipage} \\
\midrule\noalign{}
\endhead
\bottomrule\noalign{}
\endlastfoot
\textbf{Dia 1} & \href{d1-manha.qmd}{Dia 1 --- Manhã} &
\href{d1-tarde.qmd}{Dia 1 --- Tarde} \\
\textbf{Dia 2} & \href{d2-manha.qmd}{Dia 2 --- Manhã} &
\href{d2-tarde.qmd}{Dia 2 --- Tarde} \\
\textbf{Dia 3} & \href{d3-manha.qmd}{Dia 3 --- Manhã} &
\href{d3-tarde.qmd}{Dia 3 --- Tarde} \\
\textbf{Dia 4} & \href{d4-manha.qmd}{Dia 4 --- Manhã} &
\href{d4-tarde.qmd}{Dia 4 --- Tarde} \\
\end{longtable}
}

\section{Capítulos}\label{capuxedtulos-2}

\begin{itemize}
\tightlist
\item
  \href{d1-manha.qmd}{Dia 1 --- Manhã} --- Introdução ao R e RStudio
\item
  \href{d1-tarde.qmd}{Dia 1 --- Tarde} --- Manipulação de dados com
  dplyr
\item
  \href{d2-manha.qmd}{Dia 2 --- Manhã} --- Manipulação avançada e Joins
\item
  \href{d2-tarde.qmd}{Dia 2 --- Tarde} --- Visualização de dados com
  ggplot2
\item
  \href{d3-manha.qmd}{Dia 3 --- Manhã} --- Strings, Fatores e Datas
\item
  \href{d3-tarde.qmd}{Dia 3 --- Tarde} --- Tidy Data e Importação
  Avançada
\item
  \href{d4-manha.qmd}{Dia 4 --- Manhã} --- Iteração e Funções
\item
  \href{d4-tarde.qmd}{Dia 4 --- Tarde} --- Modelagem Estatística e
  Projeto Final
\end{itemize}

\section{Datasets}\label{datasets-2}

\begin{itemize}
\tightlist
\item
  🎬 \texttt{filmes\_imdb.csv} --- filmes (nota, orçamento, faturamento,
  gênero)
\item
  🌲 \texttt{inv\_florestal.csv}, \texttt{parcelas.csv},
  \texttt{clima.xlsx}, \texttt{talhoes.xlsx}, \texttt{experimento.xlsx}
\end{itemize}

\section{Pré-requisitos}\label{pruxe9-requisitos-3}

\begin{itemize}
\tightlist
\item
  \href{https://cran.r-project.org/}{R} (4.0+)
\item
  \href{https://posit.co/download/rstudio-desktop/}{RStudio}
\end{itemize}

\section{Repositório}\label{reposituxf3rio-3}

👉
\href{https://github.com/italocegatta/curso_r}{github.com/italocegatta/curso\_r}

\chapter{Dia 1 --- Manhã}\label{dia-1-manhuxe3-2}

\textbf{Dia 1 --- Manhã \textbar{} 4 horas}

\section*{Objetivos}\label{objetivos}
\addcontentsline{toc}{section}{Objetivos}

\markright{Objetivos}

Ao final deste capítulo, você será capaz de:

\begin{itemize}
\tightlist
\item
  Entender o que é o R e por que usá-lo
\item
  Instalar o R e o RStudio
\item
  Navegar pelo ambiente do RStudio
\item
  Criar e manipular objetos (vetores, fatores, data frames)
\item
  Executar seus primeiros scripts em R
\end{itemize}

\begin{center}\rule{0.5\linewidth}{0.5pt}\end{center}

\section{O que é o R?}\label{o-que-uxe9-o-r-1}

O R é uma linguagem de programação e um ambiente para computação
estatística e gráficos. É mantido por um grupo de voluntários de vários
países, o \textbf{R-core team}, e distribuído gratuitamente pelo site do
\href{https://cran.r-project.org/}{CRAN}.

O R nasceu nos anos 90 a partir de uma implementação da linguagem S,
criada por John Chambers no Bell Laboratories. Desde então, cresceu em
ritmo acelerado e hoje é um dos principais softwares de análise de dados
do mundo.

\begin{tcolorbox}[enhanced jigsaw, arc=.35mm, bottomrule=.15mm, bottomtitle=1mm, breakable, colback=white, colbacktitle=quarto-callout-tip-color!10!white, colframe=quarto-callout-tip-color-frame, coltitle=black, left=2mm, leftrule=.75mm, opacityback=0, opacitybacktitle=0.6, rightrule=.15mm, title=\textcolor{quarto-callout-tip-color}{\faLightbulb}\hspace{0.5em}{Por que aprender R?}, titlerule=0mm, toprule=.15mm, toptitle=1mm]

\begin{itemize}
\tightlist
\item
  \textbf{Gratuito} --- sem licenças, sem custos
\item
  \textbf{Comunidade enorme} --- mais de 20.000 pacotes disponíveis
\item
  \textbf{Reprodutibilidade} --- scripts documentam toda a análise
\item
  \textbf{Gráficos profissionais} --- ggplot2 é referência em
  visualização
\item
  \textbf{Integração} --- RMarkdown/Quarto geram relatórios, sites,
  dashboards e slides
\end{itemize}

\end{tcolorbox}

\begin{center}\rule{0.5\linewidth}{0.5pt}\end{center}

\section{Instalação}\label{instalauxe7uxe3o-1}

\subsection{R}\label{r}

\begin{enumerate}
\def\labelenumi{\arabic{enumi}.}
\tightlist
\item
  Acesse \href{https://cran.r-project.org/}{cran.r-project.org}
\item
  Escolha seu sistema operacional (Windows, Mac ou Linux)
\item
  Faça o download e execute o instalador
\item
  Siga as instruções padrão
\end{enumerate}

\subsection{RStudio}\label{rstudio}

O RStudio é um ambiente de desenvolvimento integrado (IDE) que facilita
muito o trabalho com R. Ele oferece editor de código, console,
visualizador de gráficos, gerenciador de arquivos e muito mais.

\begin{enumerate}
\def\labelenumi{\arabic{enumi}.}
\tightlist
\item
  Acesse
  \href{https://posit.co/download/rstudio-desktop/}{posit.co/download/rstudio-desktop}
\item
  Faça o download da versão gratuita (RStudio Desktop)
\item
  Instale normalmente
\end{enumerate}

\begin{tcolorbox}[enhanced jigsaw, arc=.35mm, bottomrule=.15mm, bottomtitle=1mm, breakable, colback=white, colbacktitle=quarto-callout-warning-color!10!white, colframe=quarto-callout-warning-color-frame, coltitle=black, left=2mm, leftrule=.75mm, opacityback=0, opacitybacktitle=0.6, rightrule=.15mm, title=\textcolor{quarto-callout-warning-color}{\faExclamationTriangle}\hspace{0.5em}{Atenção!}, titlerule=0mm, toprule=.15mm, toptitle=1mm]

Instale o \textbf{R primeiro}, depois o RStudio. O RStudio precisa
encontrar o R instalado no sistema para funcionar.

\end{tcolorbox}

\begin{center}\rule{0.5\linewidth}{0.5pt}\end{center}

\section{Conhecendo o RStudio}\label{conhecendo-o-rstudio}

Ao abrir o RStudio, você verá 4 painéis principais:

{\def\LTcaptype{none} % do not increment counter
\begin{longtable}[]{@{}lcl@{}}
\toprule\noalign{}
Painel & Localização & Função \\
\midrule\noalign{}
\endhead
\bottomrule\noalign{}
\endlastfoot
\textbf{Editor} & Superior esquerdo & Escrever scripts \texttt{.R} \\
\textbf{Console} & Inferior esquerdo & Executar comandos e ver
resultados \\
\textbf{Environment} & Superior direito & Ver objetos criados \\
\textbf{Files/Plots/Help} & Inferior direito & Arquivos, gráficos e
ajuda \\
\end{longtable}
}

\begin{tcolorbox}[enhanced jigsaw, arc=.35mm, bottomrule=.15mm, bottomtitle=1mm, breakable, colback=white, colbacktitle=quarto-callout-tip-color!10!white, colframe=quarto-callout-tip-color-frame, coltitle=black, left=2mm, leftrule=.75mm, opacityback=0, opacitybacktitle=0.6, rightrule=.15mm, title=\textcolor{quarto-callout-tip-color}{\faLightbulb}\hspace{0.5em}{Atalho essencial}, titlerule=0mm, toprule=.15mm, toptitle=1mm]

\texttt{Ctrl\ +\ Enter} (ou \texttt{Cmd\ +\ Enter} no Mac) executa a
linha atual do editor no console. É o atalho que você mais vai usar!

\end{tcolorbox}

\begin{center}\rule{0.5\linewidth}{0.5pt}\end{center}

\section{Primeiros passos no console}\label{primeiros-passos-no-console}

O console é onde o R recebe comandos e mostra resultados. Vamos começar
usando o R como uma calculadora:

\begin{Shaded}
\begin{Highlighting}[]
\CommentTok{\# Soma}
\DecValTok{2} \SpecialCharTok{+} \DecValTok{3}
\end{Highlighting}
\end{Shaded}

\begin{verbatim}
[1] 5
\end{verbatim}

\begin{Shaded}
\begin{Highlighting}[]
\CommentTok{\# Multiplicação}
\DecValTok{4} \SpecialCharTok{*} \DecValTok{5}
\end{Highlighting}
\end{Shaded}

\begin{verbatim}
[1] 20
\end{verbatim}

\begin{Shaded}
\begin{Highlighting}[]
\CommentTok{\# Potência}
\DecValTok{2}\SpecialCharTok{\^{}}\DecValTok{3}
\end{Highlighting}
\end{Shaded}

\begin{verbatim}
[1] 8
\end{verbatim}

\begin{Shaded}
\begin{Highlighting}[]
\CommentTok{\# Raiz quadrada}
\FunctionTok{sqrt}\NormalTok{(}\DecValTok{81}\NormalTok{)}
\end{Highlighting}
\end{Shaded}

\begin{verbatim}
[1] 9
\end{verbatim}

\begin{Shaded}
\begin{Highlighting}[]
\CommentTok{\# Expressão completa}
\NormalTok{(}\DecValTok{10} \SpecialCharTok{+} \DecValTok{5}\NormalTok{) }\SpecialCharTok{*} \DecValTok{3} \SpecialCharTok{/} \DecValTok{2}
\end{Highlighting}
\end{Shaded}

\begin{verbatim}
[1] 22.5
\end{verbatim}

\subsection{Atribuição: guardando valores em
objetos}\label{atribuiuxe7uxe3o-guardando-valores-em-objetos}

Para guardar resultados e reutilizá-los depois, usamos o operador de
atribuição \texttt{\textless{}-}:

\begin{Shaded}
\begin{Highlighting}[]
\CommentTok{\# Guardando um número}
\NormalTok{nota\_filme }\OtherTok{\textless{}{-}} \FloatTok{8.5}
\NormalTok{nota\_filme}
\end{Highlighting}
\end{Shaded}

\begin{verbatim}
[1] 8.5
\end{verbatim}

\begin{Shaded}
\begin{Highlighting}[]
\CommentTok{\# Guardando texto}
\NormalTok{titulo }\OtherTok{\textless{}{-}} \StringTok{"O Poderoso Chefão"}
\NormalTok{titulo}
\end{Highlighting}
\end{Shaded}

\begin{verbatim}
[1] "O Poderoso Chefão"
\end{verbatim}

\begin{Shaded}
\begin{Highlighting}[]
\CommentTok{\# Usando o objeto em operações}
\NormalTok{nota\_filme }\SpecialCharTok{*} \DecValTok{2}
\end{Highlighting}
\end{Shaded}

\begin{verbatim}
[1] 17
\end{verbatim}

\begin{Shaded}
\begin{Highlighting}[]
\NormalTok{nota\_filme }\SpecialCharTok{+} \FloatTok{1.5}
\end{Highlighting}
\end{Shaded}

\begin{verbatim}
[1] 10
\end{verbatim}

\begin{tcolorbox}[enhanced jigsaw, arc=.35mm, bottomrule=.15mm, bottomtitle=1mm, breakable, colback=white, colbacktitle=quarto-callout-important-color!10!white, colframe=quarto-callout-important-color-frame, coltitle=black, left=2mm, leftrule=.75mm, opacityback=0, opacitybacktitle=0.6, rightrule=.15mm, title=\textcolor{quarto-callout-important-color}{\faExclamation}\hspace{0.5em}{\texttt{\textless{}-} ou \texttt{=}?}, titlerule=0mm, toprule=.15mm, toptitle=1mm]

Em R, o operador tradicional de atribuição é \texttt{\textless{}-}
(atalho: \texttt{Alt\ +\ -}). Você também pode usar \texttt{=}, mas
\texttt{\textless{}-} é o padrão da comunidade R. Use o que preferir,
mas seja consistente.

\end{tcolorbox}

\subsection{Comentários}\label{comentuxe1rios}

Tudo que vem depois de \texttt{\#} é ignorado pelo R. Use comentários
para explicar seu código:

\begin{Shaded}
\begin{Highlighting}[]
\CommentTok{\# Isto é um comentário — o R ignora esta linha}
\NormalTok{x }\OtherTok{\textless{}{-}} \DecValTok{10}  \CommentTok{\# comentário ao lado do código também funciona}
\end{Highlighting}
\end{Shaded}

\subsection{Pedindo ajuda}\label{pedindo-ajuda}

\begin{Shaded}
\begin{Highlighting}[]
\CommentTok{\# Ajuda de uma função específica}
\CommentTok{\# ?mean}
\CommentTok{\# ?ggplot}

\CommentTok{\# Busca por termo}
\CommentTok{\# ??histograma}
\end{Highlighting}
\end{Shaded}

\begin{center}\rule{0.5\linewidth}{0.5pt}\end{center}

\section{Estruturas de dados no R}\label{estruturas-de-dados-no-r}

No R, tudo é um \textbf{objeto}. Os objetos podem ser de diferentes
tipos e estruturas. Entender essas estruturas é fundamental --- é com
elas que você vai trabalhar todos os dias.

\subsection{Vetores atômicos}\label{vetores-atuxf4micos}

O vetor é a estrutura mais básica: uma sequência de valores do mesmo
tipo.

\begin{Shaded}
\begin{Highlighting}[]
\CommentTok{\# Vetor numérico}
\NormalTok{notas }\OtherTok{\textless{}{-}} \FunctionTok{c}\NormalTok{(}\FloatTok{7.5}\NormalTok{, }\FloatTok{8.0}\NormalTok{, }\FloatTok{9.2}\NormalTok{, }\FloatTok{6.8}\NormalTok{, }\FloatTok{8.7}\NormalTok{)}
\NormalTok{notas}
\end{Highlighting}
\end{Shaded}

\begin{verbatim}
[1] 7.5 8.0 9.2 6.8 8.7
\end{verbatim}

\begin{Shaded}
\begin{Highlighting}[]
\CommentTok{\# Vetor de texto (character)}
\NormalTok{filmes }\OtherTok{\textless{}{-}} \FunctionTok{c}\NormalTok{(}\StringTok{"Matrix"}\NormalTok{, }\StringTok{"Interestelar"}\NormalTok{, }\StringTok{"Parasita"}\NormalTok{)}
\NormalTok{filmes}
\end{Highlighting}
\end{Shaded}

\begin{verbatim}
[1] "Matrix"       "Interestelar" "Parasita"    
\end{verbatim}

\begin{Shaded}
\begin{Highlighting}[]
\CommentTok{\# Vetor lógico (TRUE/FALSE)}
\NormalTok{filmes\_brasileiros }\OtherTok{\textless{}{-}} \FunctionTok{c}\NormalTok{(}\ConstantTok{TRUE}\NormalTok{, }\ConstantTok{FALSE}\NormalTok{, }\ConstantTok{FALSE}\NormalTok{)}
\NormalTok{filmes\_brasileiros}
\end{Highlighting}
\end{Shaded}

\begin{verbatim}
[1]  TRUE FALSE FALSE
\end{verbatim}

\begin{Shaded}
\begin{Highlighting}[]
\CommentTok{\# Operações com vetores}
\FunctionTok{mean}\NormalTok{(notas)       }\CommentTok{\# média}
\end{Highlighting}
\end{Shaded}

\begin{verbatim}
[1] 8.04
\end{verbatim}

\begin{Shaded}
\begin{Highlighting}[]
\FunctionTok{sum}\NormalTok{(notas)        }\CommentTok{\# soma}
\end{Highlighting}
\end{Shaded}

\begin{verbatim}
[1] 40.2
\end{verbatim}

\begin{Shaded}
\begin{Highlighting}[]
\FunctionTok{length}\NormalTok{(notas)     }\CommentTok{\# comprimento}
\end{Highlighting}
\end{Shaded}

\begin{verbatim}
[1] 5
\end{verbatim}

\begin{Shaded}
\begin{Highlighting}[]
\FunctionTok{max}\NormalTok{(notas)        }\CommentTok{\# máximo}
\end{Highlighting}
\end{Shaded}

\begin{verbatim}
[1] 9.2
\end{verbatim}

\begin{Shaded}
\begin{Highlighting}[]
\FunctionTok{min}\NormalTok{(notas)        }\CommentTok{\# mínimo}
\end{Highlighting}
\end{Shaded}

\begin{verbatim}
[1] 6.8
\end{verbatim}

\begin{tcolorbox}[enhanced jigsaw, arc=.35mm, bottomrule=.15mm, bottomtitle=1mm, breakable, colback=white, colbacktitle=quarto-callout-tip-color!10!white, colframe=quarto-callout-tip-color-frame, coltitle=black, left=2mm, leftrule=.75mm, opacityback=0, opacitybacktitle=0.6, rightrule=.15mm, title=\textcolor{quarto-callout-tip-color}{\faLightbulb}\hspace{0.5em}{Use \texttt{c()} para criar vetores}, titlerule=0mm, toprule=.15mm, toptitle=1mm]

A função \texttt{c()} significa \emph{combine} (combinar). Use-a sempre
que quiser juntar valores em um vetor.

\end{tcolorbox}

\subsection{Valores especiais}\label{valores-especiais}

\begin{Shaded}
\begin{Highlighting}[]
\CommentTok{\# NA = Not Available (dado ausente)}
\NormalTok{dados }\OtherTok{\textless{}{-}} \FunctionTok{c}\NormalTok{(}\DecValTok{23}\NormalTok{, }\DecValTok{45}\NormalTok{, }\ConstantTok{NA}\NormalTok{, }\DecValTok{32}\NormalTok{, }\ConstantTok{NA}\NormalTok{)}
\NormalTok{dados}
\end{Highlighting}
\end{Shaded}

\begin{verbatim}
[1] 23 45 NA 32 NA
\end{verbatim}

\begin{Shaded}
\begin{Highlighting}[]
\FunctionTok{mean}\NormalTok{(dados)             }\CommentTok{\# NA — precisa lidar com os ausentes}
\end{Highlighting}
\end{Shaded}

\begin{verbatim}
[1] NA
\end{verbatim}

\begin{Shaded}
\begin{Highlighting}[]
\FunctionTok{mean}\NormalTok{(dados, }\AttributeTok{na.rm =} \ConstantTok{TRUE}\NormalTok{)  }\CommentTok{\# remove os NA antes de calcular}
\end{Highlighting}
\end{Shaded}

\begin{verbatim}
[1] 33.33333
\end{verbatim}

\begin{Shaded}
\begin{Highlighting}[]
\CommentTok{\# NaN = Not a Number}
\DecValTok{0} \SpecialCharTok{/} \DecValTok{0}
\end{Highlighting}
\end{Shaded}

\begin{verbatim}
[1] NaN
\end{verbatim}

\begin{Shaded}
\begin{Highlighting}[]
\CommentTok{\# NULL = ausência de valor (objeto vazio)}
\ConstantTok{NULL}
\end{Highlighting}
\end{Shaded}

\begin{verbatim}
NULL
\end{verbatim}

\begin{Shaded}
\begin{Highlighting}[]
\CommentTok{\# Inf = infinito}
\DecValTok{1} \SpecialCharTok{/} \DecValTok{0}
\end{Highlighting}
\end{Shaded}

\begin{verbatim}
[1] Inf
\end{verbatim}

\subsection{Conversão entre tipos
(coerção)}\label{conversuxe3o-entre-tipos-coeruxe7uxe3o}

Às vezes você precisa converter um tipo de dado em outro:

\begin{Shaded}
\begin{Highlighting}[]
\CommentTok{\# Texto para número}
\FunctionTok{as.numeric}\NormalTok{(}\StringTok{"42"}\NormalTok{)}
\end{Highlighting}
\end{Shaded}

\begin{verbatim}
[1] 42
\end{verbatim}

\begin{Shaded}
\begin{Highlighting}[]
\CommentTok{\# Número para texto}
\FunctionTok{as.character}\NormalTok{(}\DecValTok{42}\NormalTok{)}
\end{Highlighting}
\end{Shaded}

\begin{verbatim}
[1] "42"
\end{verbatim}

\begin{Shaded}
\begin{Highlighting}[]
\CommentTok{\# Número para lógico}
\FunctionTok{as.logical}\NormalTok{(}\FunctionTok{c}\NormalTok{(}\DecValTok{1}\NormalTok{, }\DecValTok{0}\NormalTok{, }\DecValTok{1}\NormalTok{, }\DecValTok{0}\NormalTok{))}
\end{Highlighting}
\end{Shaded}

\begin{verbatim}
[1]  TRUE FALSE  TRUE FALSE
\end{verbatim}

\begin{Shaded}
\begin{Highlighting}[]
\CommentTok{\# Cuidado: conversão impossível gera NA}
\FunctionTok{as.numeric}\NormalTok{(}\StringTok{"quarenta e dois"}\NormalTok{)}
\end{Highlighting}
\end{Shaded}

\begin{verbatim}
[1] NA
\end{verbatim}

\subsection{Fatores}\label{fatores}

Fatores são usados para representar \textbf{variáveis categóricas} ---
como gênero de filme, espécie de árvore ou tipo de tratamento.

\begin{Shaded}
\begin{Highlighting}[]
\CommentTok{\# Criando um fator}
\NormalTok{genero }\OtherTok{\textless{}{-}} \FunctionTok{factor}\NormalTok{(}\FunctionTok{c}\NormalTok{(}\StringTok{"Ação"}\NormalTok{, }\StringTok{"Drama"}\NormalTok{, }\StringTok{"Comédia"}\NormalTok{, }\StringTok{"Drama"}\NormalTok{, }\StringTok{"Ação"}\NormalTok{))}
\NormalTok{genero}
\end{Highlighting}
\end{Shaded}

\begin{verbatim}
[1] Ação    Drama   Comédia Drama   Ação   
Levels: Ação Comédia Drama
\end{verbatim}

\begin{Shaded}
\begin{Highlighting}[]
\CommentTok{\# Níveis do fator}
\FunctionTok{levels}\NormalTok{(genero)}
\end{Highlighting}
\end{Shaded}

\begin{verbatim}
[1] "Ação"    "Comédia" "Drama"  
\end{verbatim}

\begin{Shaded}
\begin{Highlighting}[]
\CommentTok{\# Contagem por nível}
\FunctionTok{table}\NormalTok{(genero)}
\end{Highlighting}
\end{Shaded}

\begin{verbatim}
genero
   Ação Comédia   Drama 
      2       1       2 
\end{verbatim}

\subsection{Listas}\label{listas}

Listas podem conter elementos de \textbf{tipos diferentes} --- diferente
dos vetores, onde tudo precisa ser do mesmo tipo.

\begin{Shaded}
\begin{Highlighting}[]
\CommentTok{\# Uma lista com número, texto e vetor}
\NormalTok{filme\_info }\OtherTok{\textless{}{-}} \FunctionTok{list}\NormalTok{(}
  \AttributeTok{titulo =} \StringTok{"Matrix"}\NormalTok{,}
  \AttributeTok{ano =} \DecValTok{1999}\NormalTok{,}
  \AttributeTok{generos =} \FunctionTok{c}\NormalTok{(}\StringTok{"Ação"}\NormalTok{, }\StringTok{"Ficção científica"}\NormalTok{),}
  \AttributeTok{nota =} \FloatTok{8.7}
\NormalTok{)}
\NormalTok{filme\_info}
\end{Highlighting}
\end{Shaded}

\begin{verbatim}
$titulo
[1] "Matrix"

$ano
[1] 1999

$generos
[1] "Ação"              "Ficção científica"

$nota
[1] 8.7
\end{verbatim}

\begin{Shaded}
\begin{Highlighting}[]
\CommentTok{\# Acessando elementos da lista}
\NormalTok{filme\_info}\SpecialCharTok{$}\NormalTok{titulo}
\end{Highlighting}
\end{Shaded}

\begin{verbatim}
[1] "Matrix"
\end{verbatim}

\begin{Shaded}
\begin{Highlighting}[]
\NormalTok{filme\_info}\SpecialCharTok{$}\NormalTok{ano}
\end{Highlighting}
\end{Shaded}

\begin{verbatim}
[1] 1999
\end{verbatim}

\subsection{Data frames e tibbles}\label{data-frames-e-tibbles}

O data frame é a estrutura mais importante para análise de dados: uma
\textbf{tabela} com linhas (observações) e colunas (variáveis).

\begin{Shaded}
\begin{Highlighting}[]
\CommentTok{\# Criando um data frame manualmente}
\NormalTok{df }\OtherTok{\textless{}{-}} \FunctionTok{data.frame}\NormalTok{(}
  \AttributeTok{titulo =} \FunctionTok{c}\NormalTok{(}\StringTok{"Matrix"}\NormalTok{, }\StringTok{"Interestelar"}\NormalTok{, }\StringTok{"Parasita"}\NormalTok{),}
  \AttributeTok{ano =} \FunctionTok{c}\NormalTok{(}\DecValTok{1999}\NormalTok{, }\DecValTok{2014}\NormalTok{, }\DecValTok{2019}\NormalTok{),}
  \AttributeTok{nota =} \FunctionTok{c}\NormalTok{(}\FloatTok{8.7}\NormalTok{, }\FloatTok{8.6}\NormalTok{, }\FloatTok{8.5}\NormalTok{)}
\NormalTok{)}
\NormalTok{df}
\end{Highlighting}
\end{Shaded}

\begin{verbatim}
        titulo  ano nota
1       Matrix 1999  8.7
2 Interestelar 2014  8.6
3     Parasita 2019  8.5
\end{verbatim}

No \texttt{tidyverse}, usamos uma versão moderna do data frame chamada
\textbf{tibble}:

\begin{Shaded}
\begin{Highlighting}[]
\FunctionTok{library}\NormalTok{(dplyr)}
\FunctionTok{library}\NormalTok{(ggplot2)}
\FunctionTok{library}\NormalTok{(readr)}
\FunctionTok{library}\NormalTok{(tidyr)}
\FunctionTok{library}\NormalTok{(stringr)}
\FunctionTok{library}\NormalTok{(forcats)}
\FunctionTok{library}\NormalTok{(purrr)}
\FunctionTok{library}\NormalTok{(tibble)}

\FunctionTok{tibble}\NormalTok{(}
  \AttributeTok{titulo =} \FunctionTok{c}\NormalTok{(}\StringTok{"Matrix"}\NormalTok{, }\StringTok{"Interestelar"}\NormalTok{, }\StringTok{"Parasita"}\NormalTok{),}
  \AttributeTok{ano =} \FunctionTok{c}\NormalTok{(}\DecValTok{1999}\NormalTok{, }\DecValTok{2014}\NormalTok{, }\DecValTok{2019}\NormalTok{),}
  \AttributeTok{nota =} \FunctionTok{c}\NormalTok{(}\FloatTok{8.7}\NormalTok{, }\FloatTok{8.6}\NormalTok{, }\FloatTok{8.5}\NormalTok{)}
\NormalTok{)}
\end{Highlighting}
\end{Shaded}

\begin{verbatim}
# A tibble: 3 x 3
  titulo         ano  nota
  <chr>        <dbl> <dbl>
1 Matrix        1999   8.7
2 Interestelar  2014   8.6
3 Parasita      2019   8.5
\end{verbatim}

\subsection{Conhecendo seus dados}\label{conhecendo-seus-dados}

Três funções essenciais para inspecionar qualquer dataset:

\begin{Shaded}
\begin{Highlighting}[]
\CommentTok{\# Carregando o dataset de filmes}
\FunctionTok{library}\NormalTok{(readr)}
\NormalTok{filmes }\OtherTok{\textless{}{-}} \FunctionTok{read\_csv2}\NormalTok{(}\StringTok{"../data/filmes\_imdb.csv"}\NormalTok{)}

\CommentTok{\# Estrutura do dataset}
\FunctionTok{glimpse}\NormalTok{(filmes)}
\end{Highlighting}
\end{Shaded}

\begin{verbatim}
Rows: 3,875
Columns: 10
$ codigo_filme     <dbl> 1, 2, 3, 4, 5, 6, 7, 8, 9, 10, 11, 12, 13, 14, 15, 16~
$ titulo           <chr> "Avatar ", "Pirates of the Caribbean: At World's End ~
$ diretor          <chr> "James Cameron", "Gore Verbinski", "Sam Mendes", "Chr~
$ ano              <dbl> 2009, 2007, 2015, 2012, 2012, 2007, 2010, 2015, 2009,~
$ genero           <chr> "Ação", "Aventura", "Comédia", "Mistério", "Fantasia"~
$ duracao          <dbl> 81, 126, 88, 99, 92, 95, 39, 111, 82, 89, 77, 102, 98~
$ orcamento        <dbl> 237000000, 300000000, 245000000, 250000000, 263700000~
$ faturamento      <dbl> 760505847, 309404152, 200074175, 448130642, 73058679,~
$ nota             <dbl> 7.9, 7.1, 6.8, 8.5, 6.6, 6.2, 7.8, 7.5, 7.5, 6.9, 6.1~
$ elenco_principal <chr> "Viola Davis", "Cate Blanchett", "Natalie Portman", "~
\end{verbatim}

\begin{Shaded}
\begin{Highlighting}[]
\CommentTok{\# Primeiras linhas}
\FunctionTok{head}\NormalTok{(filmes, }\DecValTok{5}\NormalTok{)}
\end{Highlighting}
\end{Shaded}

\begin{verbatim}
# A tibble: 5 x 10
  codigo_filme titulo   diretor   ano genero duracao orcamento faturamento  nota
         <dbl> <chr>    <chr>   <dbl> <chr>    <dbl>     <dbl>       <dbl> <dbl>
1            1 Avatar   James ~  2009 Ação        81 237000000   760505847   7.9
2            2 Pirates~ Gore V~  2007 Avent~     126 300000000   309404152   7.1
3            3 Spectre  Sam Me~  2015 Coméd~      88 245000000   200074175   6.8
4            4 The Dar~ Christ~  2012 Misté~      99 250000000   448130642   8.5
5            5 John Ca~ Andrew~  2012 Fanta~      92 263700000    73058679   6.6
# i 1 more variable: elenco_principal <chr>
\end{verbatim}

\begin{Shaded}
\begin{Highlighting}[]
\CommentTok{\# Resumo estatístico das colunas}
\FunctionTok{summary}\NormalTok{(filmes)}
\end{Highlighting}
\end{Shaded}

\begin{verbatim}
  codigo_filme       titulo            diretor               ano      
 Min.   :   1.0   Length:3875        Length:3875        Min.   :1920  
 1st Qu.: 971.5   Class :character   Class :character   1st Qu.:1999  
 Median :1942.0   Mode  :character   Mode  :character   Median :2005  
 Mean   :1941.6                                         Mean   :2003  
 3rd Qu.:2911.5                                         3rd Qu.:2010  
 Max.   :3881.0                                         Max.   :2016  
    genero             duracao        orcamento          faturamento       
 Length:3875        Min.   : 31.0   Min.   :      218   Min.   :      162  
 Class :character   1st Qu.: 89.0   1st Qu.: 10000000   1st Qu.:  6856116  
 Mode  :character   Median :105.0   Median : 24000000   Median : 28064226  
                    Mean   :104.9   Mean   : 38165464   Mean   : 51085555  
                    3rd Qu.:121.0   3rd Qu.: 50000000   3rd Qu.: 65517534  
                    Max.   :201.0   Max.   :553632000   Max.   :760505847  
      nota       elenco_principal  
 Min.   :1.600   Length:3875       
 1st Qu.:5.900   Class :character  
 Median :6.600   Mode  :character  
 Mean   :6.462                     
 3rd Qu.:7.200                     
 Max.   :9.300                     
\end{verbatim}

\begin{Shaded}
\begin{Highlighting}[]
\CommentTok{\# Dimensões}
\FunctionTok{dim}\NormalTok{(filmes)     }\CommentTok{\# linhas × colunas}
\end{Highlighting}
\end{Shaded}

\begin{verbatim}
[1] 3875   10
\end{verbatim}

\begin{Shaded}
\begin{Highlighting}[]
\FunctionTok{nrow}\NormalTok{(filmes)    }\CommentTok{\# número de linhas}
\end{Highlighting}
\end{Shaded}

\begin{verbatim}
[1] 3875
\end{verbatim}

\begin{Shaded}
\begin{Highlighting}[]
\FunctionTok{ncol}\NormalTok{(filmes)    }\CommentTok{\# número de colunas}
\end{Highlighting}
\end{Shaded}

\begin{verbatim}
[1] 10
\end{verbatim}

\begin{Shaded}
\begin{Highlighting}[]
\FunctionTok{names}\NormalTok{(filmes)   }\CommentTok{\# nomes das colunas}
\end{Highlighting}
\end{Shaded}

\begin{verbatim}
 [1] "codigo_filme"     "titulo"           "diretor"          "ano"             
 [5] "genero"           "duracao"          "orcamento"        "faturamento"     
 [9] "nota"             "elenco_principal"
\end{verbatim}

\begin{center}\rule{0.5\linewidth}{0.5pt}\end{center}

\section{Exercícios}\label{exercuxedcios-7}

\begin{tcolorbox}[enhanced jigsaw, arc=.35mm, bottomrule=.15mm, bottomtitle=1mm, breakable, colback=white, colbacktitle=quarto-callout-note-color!10!white, colframe=quarto-callout-note-color-frame, coltitle=black, left=2mm, leftrule=.75mm, opacityback=0, opacitybacktitle=0.6, rightrule=.15mm, title=\textcolor{quarto-callout-note-color}{\faInfo}\hspace{0.5em}{Como fazer os exercícios}, titlerule=0mm, toprule=.15mm, toptitle=1mm]

Cada exercício mostra o \textbf{enunciado} e o \textbf{resultado
esperado}. O código é com você! Escreva sua solução, execute e confira
se o resultado bate com o que está no livro.

\end{tcolorbox}

\subsection{Exercício 1}\label{exercuxedcio-1-3}

Crie um vetor com as notas de 5 filmes e calcule a média.

\textbf{Resultado esperado:}

\begin{verbatim}
[1] 8.4
\end{verbatim}

\subsection{Exercício 2}\label{exercuxedcio-2-3}

Crie um vetor \texttt{c(8.5,\ 9.0,\ NA,\ 7.5,\ NA)} e calcule a média
sem os valores ausentes.

\textbf{Resultado esperado:}

\begin{verbatim}
[1] 8.333333
\end{verbatim}

\subsection{Exercício 3}\label{exercuxedcio-3-2}

Crie um fator com os gêneros \texttt{"Ação"}, \texttt{"Drama"},
\texttt{"Comédia"}, \texttt{"Drama"}, \texttt{"Ação"}. Mostre a tabela
de frequência.

\textbf{Resultado esperado:}

\begin{verbatim}
Ação   Comédia   Drama 
   2         1       2 
\end{verbatim}

\subsection{Exercício 4}\label{exercuxedcio-4-1}

Carregue o dataset \texttt{filmes\_imdb.csv} e responda:

\textbf{a)} Quantas linhas e quantas colunas o dataset possui?

\textbf{Resultado esperado:}

\begin{verbatim}
[1] 4826   11
\end{verbatim}

\textbf{b)} Qual é a nota média dos filmes?

\textbf{Resultado esperado:}

\begin{verbatim}
[1] 6.37
\end{verbatim}

\textbf{c)} Qual é o filme mais recente e qual o mais antigo?

\textbf{Resultado esperado:}

\begin{verbatim}
# A tibble: 1 × 2
   ano    n
  <dbl> <int>
1  2016  4826
\end{verbatim}

\chapter{Dia 1 --- Tarde}\label{dia-1-tarde-2}

\textbf{Dia 1 --- Tarde \textbar{} 4 horas}

\section*{Objetivos}\label{objetivos-1}
\addcontentsline{toc}{section}{Objetivos}

\markright{Objetivos}

Ao final deste capítulo, você será capaz de:

\begin{itemize}
\tightlist
\item
  Importar dados de arquivos CSV
\item
  Filtrar, selecionar, transformar e resumir dados com \texttt{dplyr}
\item
  Usar o pipe (\texttt{\%\textgreater{}\%}) para encadear operações
\item
  Criar seus primeiros gráficos com \texttt{ggplot2}
\end{itemize}

\begin{center}\rule{0.5\linewidth}{0.5pt}\end{center}

\section{Importando dados}\label{importando-dados}

Tudo começa com a importação. O pacote \texttt{readr} (parte do
\texttt{tidyverse}) oferece funções rápidas e consistentes para ler
arquivos:

\begin{Shaded}
\begin{Highlighting}[]
\FunctionTok{library}\NormalTok{(dplyr)}
\FunctionTok{library}\NormalTok{(ggplot2)}
\FunctionTok{library}\NormalTok{(readr)}
\FunctionTok{library}\NormalTok{(tidyr)}
\FunctionTok{library}\NormalTok{(stringr)}
\FunctionTok{library}\NormalTok{(forcats)}
\FunctionTok{library}\NormalTok{(purrr)}
\FunctionTok{library}\NormalTok{(tibble)}

\CommentTok{\# CSV com vírgula como separador (padrão americano)}
\NormalTok{filmes }\OtherTok{\textless{}{-}} \FunctionTok{read\_csv2}\NormalTok{(}\StringTok{"../data/filmes\_imdb.csv"}\NormalTok{)}

\CommentTok{\# Se o arquivo usar ponto{-}e{-}vírgula (comum no Brasil)}
\CommentTok{\# filmes \textless{}{-} read\_csv2("../data/filmes\_imdb.csv")}
\end{Highlighting}
\end{Shaded}

\begin{Shaded}
\begin{Highlighting}[]
\FunctionTok{glimpse}\NormalTok{(filmes)}
\end{Highlighting}
\end{Shaded}

\begin{verbatim}
Rows: 3,875
Columns: 10
$ codigo_filme     <dbl> 1, 2, 3, 4, 5, 6, 7, 8, 9, 10, 11, 12, 13, 14, 15, 16~
$ titulo           <chr> "Avatar ", "Pirates of the Caribbean: At World's End ~
$ diretor          <chr> "James Cameron", "Gore Verbinski", "Sam Mendes", "Chr~
$ ano              <dbl> 2009, 2007, 2015, 2012, 2012, 2007, 2010, 2015, 2009,~
$ genero           <chr> "Ação", "Aventura", "Comédia", "Mistério", "Fantasia"~
$ duracao          <dbl> 81, 126, 88, 99, 92, 95, 39, 111, 82, 89, 77, 102, 98~
$ orcamento        <dbl> 237000000, 300000000, 245000000, 250000000, 263700000~
$ faturamento      <dbl> 760505847, 309404152, 200074175, 448130642, 73058679,~
$ nota             <dbl> 7.9, 7.1, 6.8, 8.5, 6.6, 6.2, 7.8, 7.5, 7.5, 6.9, 6.1~
$ elenco_principal <chr> "Viola Davis", "Cate Blanchett", "Natalie Portman", "~
\end{verbatim}

Temos 4826 filmes com 11 colunas: título, ano, nota, orçamento,
faturamento, gênero, diretor, duração, etc.

\begin{center}\rule{0.5\linewidth}{0.5pt}\end{center}

\section{\texorpdfstring{O pipe:
\texttt{\%\textgreater{}\%}}{O pipe: \%\textgreater\%}}\label{o-pipe-1}

O pipe (\texttt{\%\textgreater{}\%}) é o encadeamento de operações. Em
vez de aninhar funções ou criar objetos intermediários, você conecta os
passos com \texttt{\%\textgreater{}\%}:

\begin{Shaded}
\begin{Highlighting}[]
\CommentTok{\# Sem pipe (difícil de ler):}
\FunctionTok{head}\NormalTok{(}\FunctionTok{arrange}\NormalTok{(}\FunctionTok{summarise}\NormalTok{(}\FunctionTok{group\_by}\NormalTok{(}\FunctionTok{filter}\NormalTok{(filmes, ano }\SpecialCharTok{\textgreater{}} \DecValTok{2000}\NormalTok{), genero), }\AttributeTok{nota\_media =} \FunctionTok{mean}\NormalTok{(nota)), }\FunctionTok{desc}\NormalTok{(nota\_media)), }\DecValTok{5}\NormalTok{)}
\end{Highlighting}
\end{Shaded}

\begin{verbatim}
# A tibble: 5 x 2
  genero    nota_media
  <chr>          <dbl>
1 Musical         6.70
2 Esporte         6.59
3 Aventura        6.55
4 Guerra          6.52
5 Biografia       6.51
\end{verbatim}

\begin{Shaded}
\begin{Highlighting}[]
\CommentTok{\# Com pipe (leitura natural, da esquerda para a direita):}
\NormalTok{filmes }\SpecialCharTok{\%\textgreater{}\%} 
  \FunctionTok{filter}\NormalTok{(ano }\SpecialCharTok{\textgreater{}} \DecValTok{2000}\NormalTok{) }\SpecialCharTok{\%\textgreater{}\%} 
  \FunctionTok{group\_by}\NormalTok{(genero) }\SpecialCharTok{\%\textgreater{}\%} 
  \FunctionTok{summarise}\NormalTok{(}\AttributeTok{nota\_media =} \FunctionTok{mean}\NormalTok{(nota)) }\SpecialCharTok{\%\textgreater{}\%} 
  \FunctionTok{arrange}\NormalTok{(}\FunctionTok{desc}\NormalTok{(nota\_media)) }\SpecialCharTok{\%\textgreater{}\%} 
  \FunctionTok{head}\NormalTok{(}\DecValTok{5}\NormalTok{)}
\end{Highlighting}
\end{Shaded}

\begin{verbatim}
# A tibble: 5 x 2
  genero    nota_media
  <chr>          <dbl>
1 Musical         6.70
2 Esporte         6.59
3 Aventura        6.55
4 Guerra          6.52
5 Biografia       6.51
\end{verbatim}

\begin{tcolorbox}[enhanced jigsaw, arc=.35mm, bottomrule=.15mm, bottomtitle=1mm, breakable, colback=white, colbacktitle=quarto-callout-tip-color!10!white, colframe=quarto-callout-tip-color-frame, coltitle=black, left=2mm, leftrule=.75mm, opacityback=0, opacitybacktitle=0.6, rightrule=.15mm, title=\textcolor{quarto-callout-tip-color}{\faLightbulb}\hspace{0.5em}{Atalho do pipe}, titlerule=0mm, toprule=.15mm, toptitle=1mm]

\texttt{Ctrl\ +\ Shift\ +\ M} (Windows/Linux) ou
\texttt{Cmd\ +\ Shift\ +\ M} (Mac). Grave esse atalho --- você vai
usá-lo centenas de vezes por dia!

\end{tcolorbox}

\begin{center}\rule{0.5\linewidth}{0.5pt}\end{center}

\section{Os 6 verbos do dplyr}\label{os-6-verbos-do-dplyr-1}

O \texttt{dplyr} é baseado em seis funções principais (verbos) que
cobrem a maioria das operações com tabelas:

{\def\LTcaptype{none} % do not increment counter
\begin{longtable}[]{@{}lll@{}}
\toprule\noalign{}
Verbo & O que faz & Equivalente mental \\
\midrule\noalign{}
\endhead
\bottomrule\noalign{}
\endlastfoot
\texttt{filter()} & Filtra \textbf{linhas} & ``Quero só
esses\ldots{}'' \\
\texttt{select()} & Seleciona \textbf{colunas} & ``Só me
interessa\ldots{}'' \\
\texttt{mutate()} & Cria/modifica \textbf{colunas} & ``Calcula pra
mim\ldots{}'' \\
\texttt{arrange()} & \textbf{Ordena} linhas & ``Organiza
por\ldots{}'' \\
\texttt{summarise()} & \textbf{Resume} em um número & ``Resume
isso\ldots{}'' \\
\texttt{group\_by()} & \textbf{Agrupa} para operações & ``Para
cada\ldots{}'' \\
\end{longtable}
}

\begin{center}\rule{0.5\linewidth}{0.5pt}\end{center}

\subsection{\texorpdfstring{\texttt{filter()} --- Filtrar
linhas}{filter() --- Filtrar linhas}}\label{filter-filtrar-linhas}

Filtra linhas com base em condições lógicas:

\begin{Shaded}
\begin{Highlighting}[]
\CommentTok{\# Filmes com nota maior ou igual a 9.0}
\NormalTok{filmes }\SpecialCharTok{\%\textgreater{}\%} 
  \FunctionTok{filter}\NormalTok{(nota }\SpecialCharTok{\textgreater{}=} \FloatTok{9.0}\NormalTok{)}
\end{Highlighting}
\end{Shaded}

\begin{verbatim}
# A tibble: 4 x 10
  codigo_filme titulo   diretor   ano genero duracao orcamento faturamento  nota
         <dbl> <chr>    <chr>   <dbl> <chr>    <dbl>     <dbl>       <dbl> <dbl>
1           65 The Dar~ Christ~  2008 Avent~     151 185000000   533316061   9  
2         1807 The Sha~ Frank ~  1994 Fanta~     104  25000000    28341469   9.3
3         2563 The God~ Franci~  1974 Coméd~      45  13000000    57300000   9  
4         3035 The God~ Franci~  1972 Crime      131   6000000   134821952   9.2
# i 1 more variable: elenco_principal <chr>
\end{verbatim}

\begin{Shaded}
\begin{Highlighting}[]
\CommentTok{\# Filmes de ação com nota acima de 8.0}
\NormalTok{filmes }\SpecialCharTok{\%\textgreater{}\%} 
  \FunctionTok{filter}\NormalTok{(genero }\SpecialCharTok{==} \StringTok{"Ação"}\NormalTok{, nota }\SpecialCharTok{\textgreater{}} \FloatTok{8.0}\NormalTok{)}
\end{Highlighting}
\end{Shaded}

\begin{verbatim}
# A tibble: 21 x 10
   codigo_filme titulo  diretor   ano genero duracao orcamento faturamento  nota
          <dbl> <chr>   <chr>   <dbl> <chr>    <dbl>     <dbl>       <dbl> <dbl>
 1           17 The Av~ Joss W~  2012 Ação        53 220000000   623279547   8.1
 2           95 Incept~ Christ~  2010 Ação       103 160000000   292568851   8.8
 3          564 Inglou~ Quenti~  2009 Ação       108  75000000   120523073   8.3
 4          776 Deadpo~ Tim Mi~  2016 Ação       123  58000000   363024263   8.1
 5         1026 Prison~ Denis ~  2013 Ação       101  46000000    60962878   8.1
 6         1232 Amélie  Jean-P~  2001 Ação        55  77000000    33201661   8.4
 7         1341 L.A. C~ Curtis~  1997 Ação        89  35000000    64604977   8.3
 8         1775 Goodfe~ Martin~  1990 Ação       118  25000000    46836394   8.7
 9         1778 Scarfa~ Brian ~  1983 Ação        96  25000000    44700000   8.3
10         2034 Spotli~ Tom Mc~  2015 Ação        96  20000000    44988180   8.1
# i 11 more rows
# i 1 more variable: elenco_principal <chr>
\end{verbatim}

\begin{Shaded}
\begin{Highlighting}[]
\CommentTok{\# Filmes de 2010 ou mais recentes, com orçamento acima de 100 milhões}
\NormalTok{filmes }\SpecialCharTok{\%\textgreater{}\%} 
  \FunctionTok{filter}\NormalTok{(ano }\SpecialCharTok{\textgreater{}=} \DecValTok{2010}\NormalTok{, orcamento }\SpecialCharTok{\textgreater{}} \DecValTok{100000000}\NormalTok{)}
\end{Highlighting}
\end{Shaded}

\begin{verbatim}
# A tibble: 162 x 10
   codigo_filme titulo  diretor   ano genero duracao orcamento faturamento  nota
          <dbl> <chr>   <chr>   <dbl> <chr>    <dbl>     <dbl>       <dbl> <dbl>
 1            3 Spectr~ Sam Me~  2015 Coméd~      88 245000000   200074175   6.8
 2            4 The Da~ Christ~  2012 Misté~      99 250000000   448130642   8.5
 3            5 John C~ Andrew~  2012 Fanta~      92 263700000    73058679   6.6
 4            7 Tangle~ Nathan~  2010 Terror      39 260000000   200807262   7.8
 5            8 Avenge~ Joss W~  2015 Misté~     111 250000000   458991599   7.5
 6           10 Batman~ Zack S~  2016 Terror      89 250000000   330249062   6.9
 7           14 The Lo~ Gore V~  2013 Ficçã~     129 215000000    89289910   6.5
 8           15 Man of~ Zack S~  2013 Drama      123 225000000   291021565   7.2
 9           17 The Av~ Joss W~  2012 Ação        53 220000000   623279547   8.1
10           18 Pirate~ Rob Ma~  2011 Drama      113 250000000   241063875   6.7
# i 152 more rows
# i 1 more variable: elenco_principal <chr>
\end{verbatim}

\textbf{Operadores úteis no \texttt{filter()}:}

\begin{Shaded}
\begin{Highlighting}[]
\CommentTok{\# Igual a          ==}
\CommentTok{\# Diferente de     !=}
\CommentTok{\# Maior / Menor    \textgreater{}   \textless{}   \textgreater{}=   \textless{}=}
\CommentTok{\# Está na lista    \%in\% c("A", "B")}
\CommentTok{\# E                \&}
\CommentTok{\# OU               |}
\CommentTok{\# Negação          !}
\end{Highlighting}
\end{Shaded}

\begin{Shaded}
\begin{Highlighting}[]
\CommentTok{\# Filmes de Drama OU Comédia}
\NormalTok{filmes }\SpecialCharTok{\%\textgreater{}\%} 
  \FunctionTok{filter}\NormalTok{(genero }\SpecialCharTok{\%in\%} \FunctionTok{c}\NormalTok{(}\StringTok{"Drama"}\NormalTok{, }\StringTok{"Comédia"}\NormalTok{)) }\SpecialCharTok{\%\textgreater{}\%} 
  \FunctionTok{count}\NormalTok{(genero, }\AttributeTok{sort =} \ConstantTok{TRUE}\NormalTok{)}
\end{Highlighting}
\end{Shaded}

\begin{verbatim}
# A tibble: 2 x 2
  genero      n
  <chr>   <int>
1 Drama     719
2 Comédia   452
\end{verbatim}

\begin{Shaded}
\begin{Highlighting}[]
\CommentTok{\# Filmes com nota \textgreater{}= 8.0 E ano \textgreater{} 2010}
\NormalTok{filmes }\SpecialCharTok{\%\textgreater{}\%} 
  \FunctionTok{filter}\NormalTok{(nota }\SpecialCharTok{\textgreater{}=} \FloatTok{8.0}\NormalTok{, ano }\SpecialCharTok{\textgreater{}} \DecValTok{2010}\NormalTok{)}
\end{Highlighting}
\end{Shaded}

\begin{verbatim}
# A tibble: 36 x 10
   codigo_filme titulo  diretor   ano genero duracao orcamento faturamento  nota
          <dbl> <chr>   <chr>   <dbl> <chr>    <dbl>     <dbl>       <dbl> <dbl>
 1            4 The Da~ Christ~  2012 Misté~      99 250000000   448130642   8.5
 2           17 The Av~ Joss W~  2012 Ação        53 220000000   623279547   8.1
 3           47 X-Men:~ Bryan ~  2014 Ação       126 200000000   233914986   8  
 4           77 Inside~ Pete D~  2015 Coméd~     113 175000000   356454367   8.3
 5           93 Guardi~ James ~  2014 Coméd~     118 170000000   333130696   8.1
 6           94 Inters~ Christ~  2014 Roman~      90 165000000   187991439   8.6
 7          125 Mad Ma~ George~  2015 Coméd~      61 150000000   153629485   8.1
 8          175 The Re~ Alejan~  2015 Roman~      91 135000000   183635922   8.1
 9          216 Life o~ Ang Lee  2012 Thril~     123 120000000   124976634   8  
10          269 The Ma~ Ridley~  2015 Drama       78 108000000   228430993   8.1
# i 26 more rows
# i 1 more variable: elenco_principal <chr>
\end{verbatim}

\begin{center}\rule{0.5\linewidth}{0.5pt}\end{center}

\subsection{\texorpdfstring{\texttt{select()} --- Selecionar
colunas}{select() --- Selecionar colunas}}\label{select-selecionar-colunas}

\begin{Shaded}
\begin{Highlighting}[]
\CommentTok{\# Selecionar colunas por nome}
\NormalTok{filmes }\SpecialCharTok{\%\textgreater{}\%} 
  \FunctionTok{select}\NormalTok{(titulo, ano, nota)}
\end{Highlighting}
\end{Shaded}

\begin{verbatim}
# A tibble: 3,875 x 3
   titulo                                      ano  nota
   <chr>                                     <dbl> <dbl>
 1 Avatar                                     2009   7.9
 2 Pirates of the Caribbean: At World's End   2007   7.1
 3 Spectre                                    2015   6.8
 4 The Dark Knight Rises                      2012   8.5
 5 John Carter                                2012   6.6
 6 Spider-Man 3                               2007   6.2
 7 Tangled                                    2010   7.8
 8 Avengers: Age of Ultron                    2015   7.5
 9 Harry Potter and the Half-Blood Prince     2009   7.5
10 Batman v Superman: Dawn of Justice         2016   6.9
# i 3,865 more rows
\end{verbatim}

\begin{Shaded}
\begin{Highlighting}[]
\CommentTok{\# Selecionar um intervalo de colunas}
\NormalTok{filmes }\SpecialCharTok{\%\textgreater{}\%} 
  \FunctionTok{select}\NormalTok{(titulo}\SpecialCharTok{:}\NormalTok{nota)}
\end{Highlighting}
\end{Shaded}

\begin{verbatim}
# A tibble: 3,875 x 8
   titulo               diretor   ano genero duracao orcamento faturamento  nota
   <chr>                <chr>   <dbl> <chr>    <dbl>     <dbl>       <dbl> <dbl>
 1 Avatar               James ~  2009 Ação        81 237000000   760505847   7.9
 2 Pirates of the Cari~ Gore V~  2007 Avent~     126 300000000   309404152   7.1
 3 Spectre              Sam Me~  2015 Coméd~      88 245000000   200074175   6.8
 4 The Dark Knight Ris~ Christ~  2012 Misté~      99 250000000   448130642   8.5
 5 John Carter          Andrew~  2012 Fanta~      92 263700000    73058679   6.6
 6 Spider-Man 3         Sam Ra~  2007 Drama       95 258000000   336530303   6.2
 7 Tangled              Nathan~  2010 Terror      39 260000000   200807262   7.8
 8 Avengers: Age of Ul~ Joss W~  2015 Misté~     111 250000000   458991599   7.5
 9 Harry Potter and th~ David ~  2009 Ficçã~      82 250000000   301956980   7.5
10 Batman v Superman: ~ Zack S~  2016 Terror      89 250000000   330249062   6.9
# i 3,865 more rows
\end{verbatim}

\begin{Shaded}
\begin{Highlighting}[]
\CommentTok{\# Remover colunas com {-}}
\NormalTok{filmes }\SpecialCharTok{\%\textgreater{}\%} 
  \FunctionTok{select}\NormalTok{(}\SpecialCharTok{{-}}\NormalTok{codigo\_filme, }\SpecialCharTok{{-}}\NormalTok{diretor)}
\end{Highlighting}
\end{Shaded}

\begin{verbatim}
# A tibble: 3,875 x 8
   titulo        ano genero duracao orcamento faturamento  nota elenco_principal
   <chr>       <dbl> <chr>    <dbl>     <dbl>       <dbl> <dbl> <chr>           
 1 Avatar       2009 Ação        81 237000000   760505847   7.9 Viola Davis     
 2 Pirates of~  2007 Avent~     126 300000000   309404152   7.1 Cate Blanchett  
 3 Spectre      2015 Coméd~      88 245000000   200074175   6.8 Natalie Portman 
 4 The Dark K~  2012 Misté~      99 250000000   448130642   8.5 Cate Blanchett  
 5 John Carte~  2012 Fanta~      92 263700000    73058679   6.6 Tom Hanks       
 6 Spider-Man~  2007 Drama       95 258000000   336530303   6.2 Viola Davis     
 7 Tangled      2010 Terror      39 260000000   200807262   7.8 Viola Davis     
 8 Avengers: ~  2015 Misté~     111 250000000   458991599   7.5 Robert De Niro  
 9 Harry Pott~  2009 Ficçã~      82 250000000   301956980   7.5 Natalie Portman 
10 Batman v S~  2016 Terror      89 250000000   330249062   6.9 Tom Hanks       
# i 3,865 more rows
\end{verbatim}

\begin{Shaded}
\begin{Highlighting}[]
\CommentTok{\# Selecionar por padrão no nome}
\NormalTok{filmes }\SpecialCharTok{\%\textgreater{}\%} 
  \FunctionTok{select}\NormalTok{(}\FunctionTok{starts\_with}\NormalTok{(}\StringTok{"cod"}\NormalTok{))}
\end{Highlighting}
\end{Shaded}

\begin{verbatim}
# A tibble: 3,875 x 1
   codigo_filme
          <dbl>
 1            1
 2            2
 3            3
 4            4
 5            5
 6            6
 7            7
 8            8
 9            9
10           10
# i 3,865 more rows
\end{verbatim}

\begin{Shaded}
\begin{Highlighting}[]
\NormalTok{filmes }\SpecialCharTok{\%\textgreater{}\%} 
  \FunctionTok{select}\NormalTok{(}\FunctionTok{contains}\NormalTok{(}\StringTok{"mento"}\NormalTok{))}
\end{Highlighting}
\end{Shaded}

\begin{verbatim}
# A tibble: 3,875 x 2
   orcamento faturamento
       <dbl>       <dbl>
 1 237000000   760505847
 2 300000000   309404152
 3 245000000   200074175
 4 250000000   448130642
 5 263700000    73058679
 6 258000000   336530303
 7 260000000   200807262
 8 250000000   458991599
 9 250000000   301956980
10 250000000   330249062
# i 3,865 more rows
\end{verbatim}

\begin{center}\rule{0.5\linewidth}{0.5pt}\end{center}

\subsection{\texorpdfstring{\texttt{mutate()} --- Criar ou modificar
colunas}{mutate() --- Criar ou modificar colunas}}\label{mutate-criar-ou-modificar-colunas}

A função \texttt{mutate()} calcula novas colunas a partir das
existentes:

\begin{Shaded}
\begin{Highlighting}[]
\NormalTok{filmes }\SpecialCharTok{\%\textgreater{}\%} 
  \FunctionTok{mutate}\NormalTok{(}
    \AttributeTok{lucro =}\NormalTok{ faturamento }\SpecialCharTok{{-}}\NormalTok{ orcamento,}
    \AttributeTok{lucro\_milhoes =}\NormalTok{ lucro }\SpecialCharTok{/} \DecValTok{1000000}\NormalTok{,}
    \AttributeTok{duracao\_horas =}\NormalTok{ duracao }\SpecialCharTok{/} \DecValTok{60}
\NormalTok{  ) }\SpecialCharTok{\%\textgreater{}\%} 
  \FunctionTok{select}\NormalTok{(titulo, orcamento, faturamento, lucro\_milhoes, duracao\_horas)}
\end{Highlighting}
\end{Shaded}

\begin{verbatim}
# A tibble: 3,875 x 5
   titulo                      orcamento faturamento lucro_milhoes duracao_horas
   <chr>                           <dbl>       <dbl>         <dbl>         <dbl>
 1 Avatar                      237000000   760505847        524.            1.35
 2 Pirates of the Caribbean: ~ 300000000   309404152          9.40          2.1 
 3 Spectre                     245000000   200074175        -44.9           1.47
 4 The Dark Knight Rises       250000000   448130642        198.            1.65
 5 John Carter                 263700000    73058679       -191.            1.53
 6 Spider-Man 3                258000000   336530303         78.5           1.58
 7 Tangled                     260000000   200807262        -59.2           0.65
 8 Avengers: Age of Ultron     250000000   458991599        209.            1.85
 9 Harry Potter and the Half-~ 250000000   301956980         52.0           1.37
10 Batman v Superman: Dawn of~ 250000000   330249062         80.2           1.48
# i 3,865 more rows
\end{verbatim}

\begin{Shaded}
\begin{Highlighting}[]
\CommentTok{\# Qual o filme mais lucrativo?}
\NormalTok{filmes }\SpecialCharTok{\%\textgreater{}\%} 
  \FunctionTok{mutate}\NormalTok{(}\AttributeTok{lucro =}\NormalTok{ faturamento }\SpecialCharTok{{-}}\NormalTok{ orcamento) }\SpecialCharTok{\%\textgreater{}\%} 
  \FunctionTok{arrange}\NormalTok{(}\FunctionTok{desc}\NormalTok{(lucro)) }\SpecialCharTok{\%\textgreater{}\%} 
  \FunctionTok{select}\NormalTok{(titulo, lucro) }\SpecialCharTok{\%\textgreater{}\%} 
  \FunctionTok{head}\NormalTok{(}\DecValTok{5}\NormalTok{)}
\end{Highlighting}
\end{Shaded}

\begin{verbatim}
# A tibble: 5 x 2
  titulo                                  lucro
  <chr>                                   <dbl>
1 Avatar                              523505847
2 Jurassic World                      502177271
3 Titanic                             458672302
4 Star Wars: Episode IV - A New Hope  449935665
5 E.T. the Extra-Terrestrial          424449459
\end{verbatim}

\begin{Shaded}
\begin{Highlighting}[]
\CommentTok{\# Nota padronizada (z{-}score) por filme}
\NormalTok{filmes }\SpecialCharTok{\%\textgreater{}\%} 
  \FunctionTok{mutate}\NormalTok{(}\AttributeTok{nota\_z =}\NormalTok{ (nota }\SpecialCharTok{{-}} \FunctionTok{mean}\NormalTok{(nota)) }\SpecialCharTok{/} \FunctionTok{sd}\NormalTok{(nota)) }\SpecialCharTok{\%\textgreater{}\%} 
  \FunctionTok{select}\NormalTok{(titulo, nota, nota\_z) }\SpecialCharTok{\%\textgreater{}\%} 
  \FunctionTok{arrange}\NormalTok{(}\FunctionTok{desc}\NormalTok{(nota\_z)) }\SpecialCharTok{\%\textgreater{}\%} 
  \FunctionTok{head}\NormalTok{(}\DecValTok{5}\NormalTok{)}
\end{Highlighting}
\end{Shaded}

\begin{verbatim}
# A tibble: 5 x 3
  titulo                                          nota nota_z
  <chr>                                          <dbl>  <dbl>
1 The Shawshank Redemption                         9.3   2.69
2 The Godfather                                    9.2   2.60
3 The Dark Knight                                  9     2.41
4 The Godfather: Part II                           9     2.41
5 The Lord of the Rings: The Return of the King    8.9   2.31
\end{verbatim}

\begin{center}\rule{0.5\linewidth}{0.5pt}\end{center}

\subsection{\texorpdfstring{\texttt{arrange()} ---
Ordenar}{arrange() --- Ordenar}}\label{arrange-ordenar}

\begin{Shaded}
\begin{Highlighting}[]
\CommentTok{\# Ordenar por nota (crescente)}
\NormalTok{filmes }\SpecialCharTok{\%\textgreater{}\%} 
  \FunctionTok{arrange}\NormalTok{(nota) }\SpecialCharTok{\%\textgreater{}\%} 
  \FunctionTok{select}\NormalTok{(titulo, nota) }\SpecialCharTok{\%\textgreater{}\%} 
  \FunctionTok{head}\NormalTok{(}\DecValTok{5}\NormalTok{)}
\end{Highlighting}
\end{Shaded}

\begin{verbatim}
# A tibble: 5 x 2
  titulo                           nota
  <chr>                           <dbl>
1 Justin Bieber: Never Say Never    1.6
2 Disaster Movie                    1.9
3 Superbabies: Baby Geniuses 2      1.9
4 Who's Your Caddy?                 2  
5 From Justin to Kelly              2.1
\end{verbatim}

\begin{Shaded}
\begin{Highlighting}[]
\CommentTok{\# Ordenar por nota (decrescente)}
\NormalTok{filmes }\SpecialCharTok{\%\textgreater{}\%} 
  \FunctionTok{arrange}\NormalTok{(}\FunctionTok{desc}\NormalTok{(nota)) }\SpecialCharTok{\%\textgreater{}\%} 
  \FunctionTok{select}\NormalTok{(titulo, nota) }\SpecialCharTok{\%\textgreater{}\%} 
  \FunctionTok{head}\NormalTok{(}\DecValTok{5}\NormalTok{)}
\end{Highlighting}
\end{Shaded}

\begin{verbatim}
# A tibble: 5 x 2
  titulo                                          nota
  <chr>                                          <dbl>
1 The Shawshank Redemption                         9.3
2 The Godfather                                    9.2
3 The Dark Knight                                  9  
4 The Godfather: Part II                           9  
5 The Lord of the Rings: The Return of the King    8.9
\end{verbatim}

\begin{Shaded}
\begin{Highlighting}[]
\CommentTok{\# Ordenar por gênero e depois por nota (decrescente)}
\NormalTok{filmes }\SpecialCharTok{\%\textgreater{}\%} 
  \FunctionTok{arrange}\NormalTok{(genero, }\FunctionTok{desc}\NormalTok{(nota)) }\SpecialCharTok{\%\textgreater{}\%} 
  \FunctionTok{select}\NormalTok{(genero, titulo, nota) }\SpecialCharTok{\%\textgreater{}\%} 
  \FunctionTok{head}\NormalTok{(}\DecValTok{10}\NormalTok{)}
\end{Highlighting}
\end{Shaded}

\begin{verbatim}
# A tibble: 10 x 3
   genero   titulo                                              nota
   <chr>    <chr>                                              <dbl>
 1 Animação Schindler's List                                     8.9
 2 Animação The Lord of the Rings: The Fellowship of the Ring    8.8
 3 Animação The Lord of the Rings: The Two Towers                8.7
 4 Animação One Flew Over the Cuckoo's Nest                      8.7
 5 Animação Jurassic Park                                        8.1
 6 Animação The Sixth Sense                                      8.1
 7 Animação Platoon                                              8.1
 8 Animação JFK                                                  8  
 9 Animação Winged Migration                                     8  
10 Animação The Hangover                                         7.8
\end{verbatim}

\begin{center}\rule{0.5\linewidth}{0.5pt}\end{center}

\subsection{\texorpdfstring{\texttt{summarise()} e \texttt{group\_by()}
--- Resumir e
agrupar}{summarise() e group\_by() --- Resumir e agrupar}}\label{summarise-e-group_by-resumir-e-agrupar}

Juntos, \texttt{group\_by()} e \texttt{summarise()} são a dupla mais
poderosa do dplyr:

\begin{Shaded}
\begin{Highlighting}[]
\CommentTok{\# Nota média geral}
\NormalTok{filmes }\SpecialCharTok{\%\textgreater{}\%} 
  \FunctionTok{summarise}\NormalTok{(}
    \AttributeTok{nota\_media =} \FunctionTok{mean}\NormalTok{(nota),}
    \AttributeTok{nota\_desvio =} \FunctionTok{sd}\NormalTok{(nota),}
    \AttributeTok{nota\_min =} \FunctionTok{min}\NormalTok{(nota),}
    \AttributeTok{nota\_max =} \FunctionTok{max}\NormalTok{(nota),}
    \AttributeTok{n =} \FunctionTok{n}\NormalTok{()}
\NormalTok{  )}
\end{Highlighting}
\end{Shaded}

\begin{verbatim}
# A tibble: 1 x 5
  nota_media nota_desvio nota_min nota_max     n
       <dbl>       <dbl>    <dbl>    <dbl> <int>
1       6.46        1.05      1.6      9.3  3875
\end{verbatim}

\begin{Shaded}
\begin{Highlighting}[]
\CommentTok{\# Nota média por gênero}
\NormalTok{filmes }\SpecialCharTok{\%\textgreater{}\%} 
  \FunctionTok{group\_by}\NormalTok{(genero) }\SpecialCharTok{\%\textgreater{}\%} 
  \FunctionTok{summarise}\NormalTok{(}
    \AttributeTok{nota\_media =} \FunctionTok{mean}\NormalTok{(nota),}
    \AttributeTok{n =} \FunctionTok{n}\NormalTok{()}
\NormalTok{  ) }\SpecialCharTok{\%\textgreater{}\%} 
  \FunctionTok{arrange}\NormalTok{(}\FunctionTok{desc}\NormalTok{(nota\_media))}
\end{Highlighting}
\end{Shaded}

\begin{verbatim}
# A tibble: 19 x 3
   genero            nota_media     n
   <chr>                  <dbl> <int>
 1 Musical                 6.75    39
 2 Esporte                 6.73    16
 3 História                6.56    31
 4 Guerra                  6.55    33
 5 Biografia               6.50   121
 6 Comédia                 6.50   452
 7 Ação                    6.49   566
 8 Drama                   6.48   719
 9 Romance                 6.48   245
10 Documentário            6.47   117
11 Ficção científica       6.47   306
12 Fantasia                6.47    82
13 Mistério                6.46   113
14 Aventura                6.45   204
15 Thriller                6.41   185
16 Terror                  6.40   311
17 Crime                   6.39   193
18 Animação                6.29   114
19 Família                 6.1     28
\end{verbatim}

\begin{Shaded}
\begin{Highlighting}[]
\CommentTok{\# Orçamento médio por ano (apenas anos com 5+ filmes)}
\NormalTok{filmes }\SpecialCharTok{\%\textgreater{}\%} 
  \FunctionTok{group\_by}\NormalTok{(ano) }\SpecialCharTok{\%\textgreater{}\%} 
  \FunctionTok{summarise}\NormalTok{(}
    \AttributeTok{orcamento\_medio =} \FunctionTok{mean}\NormalTok{(orcamento, }\AttributeTok{na.rm =} \ConstantTok{TRUE}\NormalTok{),}
    \AttributeTok{n =} \FunctionTok{n}\NormalTok{()}
\NormalTok{  ) }\SpecialCharTok{\%\textgreater{}\%} 
  \FunctionTok{filter}\NormalTok{(n }\SpecialCharTok{\textgreater{}=} \DecValTok{5}\NormalTok{) }\SpecialCharTok{\%\textgreater{}\%} 
  \FunctionTok{arrange}\NormalTok{(}\FunctionTok{desc}\NormalTok{(orcamento\_medio)) }\SpecialCharTok{\%\textgreater{}\%} 
  \FunctionTok{head}\NormalTok{(}\DecValTok{10}\NormalTok{)}
\end{Highlighting}
\end{Shaded}

\begin{verbatim}
# A tibble: 10 x 3
     ano orcamento_medio     n
   <dbl>           <dbl> <int>
 1  2016       66007647.    68
 2  2015       53006627.   139
 3  2013       49385475.   171
 4  2014       48291950.   159
 5  2012       47744328.   166
 6  2010       47258132.   174
 7  2008       44004546.   183
 8  2011       42996023.   175
 9  2009       41423690.   187
10  2005       40487560.   186
\end{verbatim}

\begin{center}\rule{0.5\linewidth}{0.5pt}\end{center}

\subsection{\texorpdfstring{\texttt{count()} --- Atalho para
frequência}{count() --- Atalho para frequência}}\label{count-atalho-para-frequuxeancia}

\begin{Shaded}
\begin{Highlighting}[]
\CommentTok{\# Gêneros mais frequentes}
\NormalTok{filmes }\SpecialCharTok{\%\textgreater{}\%} 
  \FunctionTok{count}\NormalTok{(genero, }\AttributeTok{sort =} \ConstantTok{TRUE}\NormalTok{)}
\end{Highlighting}
\end{Shaded}

\begin{verbatim}
# A tibble: 19 x 2
   genero                n
   <chr>             <int>
 1 Drama               719
 2 Ação                566
 3 Comédia             452
 4 Terror              311
 5 Ficção científica   306
 6 Romance             245
 7 Aventura            204
 8 Crime               193
 9 Thriller            185
10 Biografia           121
11 Documentário        117
12 Animação            114
13 Mistério            113
14 Fantasia             82
15 Musical              39
16 Guerra               33
17 História             31
18 Família              28
19 Esporte              16
\end{verbatim}

\begin{Shaded}
\begin{Highlighting}[]
\CommentTok{\# Diretores mais prolíficos}
\NormalTok{filmes }\SpecialCharTok{\%\textgreater{}\%} 
  \FunctionTok{count}\NormalTok{(diretor, }\AttributeTok{sort =} \ConstantTok{TRUE}\NormalTok{) }\SpecialCharTok{\%\textgreater{}\%} 
  \FunctionTok{head}\NormalTok{(}\DecValTok{10}\NormalTok{)}
\end{Highlighting}
\end{Shaded}

\begin{verbatim}
# A tibble: 10 x 2
   diretor               n
   <chr>             <int>
 1 Steven Spielberg     25
 2 Clint Eastwood       19
 3 Woody Allen          19
 4 Ridley Scott         17
 5 Martin Scorsese      16
 6 Steven Soderbergh    16
 7 Tim Burton           16
 8 Renny Harlin         15
 9 Spike Lee            15
10 Barry Levinson       13
\end{verbatim}

\begin{center}\rule{0.5\linewidth}{0.5pt}\end{center}

\section{🎬 Seus primeiros gráficos com
ggplot2}\label{seus-primeiros-gruxe1ficos-com-ggplot2}

E se eu te disser que você já pode criar gráficos profissionais
\textbf{hoje}? Não precisa de teoria complicada --- é só seguir uma
receita simples.

A estrutura básica de qualquer gráfico no ggplot2 é:

\begin{verbatim}
ggplot(dados, aes(x = coluna_x, y = coluna_y)) +
  geom_tipo()
\end{verbatim}

Vamos criar 4 gráficos essenciais com os dados de filmes:

\subsection{Histograma --- distribuição das
notas}\label{histograma-distribuiuxe7uxe3o-das-notas}

\begin{Shaded}
\begin{Highlighting}[]
\NormalTok{filmes }\SpecialCharTok{\%\textgreater{}\%} 
  \FunctionTok{ggplot}\NormalTok{(}\FunctionTok{aes}\NormalTok{(}\AttributeTok{x =}\NormalTok{ nota)) }\SpecialCharTok{+}
  \FunctionTok{geom\_histogram}\NormalTok{(}\AttributeTok{bins =} \DecValTok{30}\NormalTok{, }\AttributeTok{fill =} \StringTok{"\#18BC9C"}\NormalTok{, }\AttributeTok{color =} \StringTok{"white"}\NormalTok{)}
\end{Highlighting}
\end{Shaded}

\pandocbounded{\includegraphics[keepaspectratio]{04-dias/d1-tarde_files/figure-pdf/unnamed-chunk-17-1.pdf}}

Cada barra mostra quantos filmes têm nota naquele intervalo. A
distribuição é aproximadamente simétrica, centrada em torno de 6.

\subsection{Dispersão --- orçamento ×
faturamento}\label{dispersuxe3o-oruxe7amento-faturamento}

\begin{Shaded}
\begin{Highlighting}[]
\NormalTok{filmes }\SpecialCharTok{\%\textgreater{}\%} 
  \FunctionTok{ggplot}\NormalTok{(}\FunctionTok{aes}\NormalTok{(}\AttributeTok{x =}\NormalTok{ orcamento, }\AttributeTok{y =}\NormalTok{ faturamento)) }\SpecialCharTok{+}
  \FunctionTok{geom\_point}\NormalTok{(}\AttributeTok{alpha =} \FloatTok{0.3}\NormalTok{, }\AttributeTok{color =} \StringTok{"\#2C3E50"}\NormalTok{) }\SpecialCharTok{+}
  \FunctionTok{geom\_smooth}\NormalTok{(}\AttributeTok{method =} \StringTok{"lm"}\NormalTok{, }\AttributeTok{se =} \ConstantTok{FALSE}\NormalTok{, }\AttributeTok{color =} \StringTok{"\#E74C3C"}\NormalTok{)}
\end{Highlighting}
\end{Shaded}

\pandocbounded{\includegraphics[keepaspectratio]{04-dias/d1-tarde_files/figure-pdf/unnamed-chunk-18-1.pdf}}

Cada ponto é um filme. A linha vermelha mostra a tendência: quanto maior
o orçamento, maior o faturamento --- em média.

\subsection{Barras --- top 10 gêneros}\label{barras-top-10-guxeaneros}

\begin{Shaded}
\begin{Highlighting}[]
\NormalTok{filmes }\SpecialCharTok{\%\textgreater{}\%} 
  \FunctionTok{count}\NormalTok{(genero, }\AttributeTok{sort =} \ConstantTok{TRUE}\NormalTok{) }\SpecialCharTok{\%\textgreater{}\%} 
  \FunctionTok{head}\NormalTok{(}\DecValTok{10}\NormalTok{) }\SpecialCharTok{\%\textgreater{}\%} 
  \FunctionTok{ggplot}\NormalTok{(}\FunctionTok{aes}\NormalTok{(}\AttributeTok{x =} \FunctionTok{reorder}\NormalTok{(genero, n), }\AttributeTok{y =}\NormalTok{ n)) }\SpecialCharTok{+}
  \FunctionTok{geom\_col}\NormalTok{(}\AttributeTok{fill =} \StringTok{"\#18BC9C"}\NormalTok{) }\SpecialCharTok{+}
  \FunctionTok{coord\_flip}\NormalTok{() }\SpecialCharTok{+}
  \FunctionTok{labs}\NormalTok{(}\AttributeTok{x =} \ConstantTok{NULL}\NormalTok{, }\AttributeTok{y =} \StringTok{"Número de filmes"}\NormalTok{)}
\end{Highlighting}
\end{Shaded}

\pandocbounded{\includegraphics[keepaspectratio]{04-dias/d1-tarde_files/figure-pdf/unnamed-chunk-19-1.pdf}}

\subsection{Boxplot --- nota por
gênero}\label{boxplot-nota-por-guxeanero}

\begin{Shaded}
\begin{Highlighting}[]
\NormalTok{filmes }\SpecialCharTok{\%\textgreater{}\%} 
  \FunctionTok{filter}\NormalTok{(genero }\SpecialCharTok{\%in\%} \FunctionTok{c}\NormalTok{(}\StringTok{"Ação"}\NormalTok{, }\StringTok{"Comédia"}\NormalTok{, }\StringTok{"Drama"}\NormalTok{, }\StringTok{"Ficção científica"}\NormalTok{, }\StringTok{"Terror"}\NormalTok{)) }\SpecialCharTok{\%\textgreater{}\%} 
  \FunctionTok{ggplot}\NormalTok{(}\FunctionTok{aes}\NormalTok{(}\AttributeTok{x =}\NormalTok{ genero, }\AttributeTok{y =}\NormalTok{ nota, }\AttributeTok{fill =}\NormalTok{ genero)) }\SpecialCharTok{+}
  \FunctionTok{geom\_boxplot}\NormalTok{() }\SpecialCharTok{+}
  \FunctionTok{labs}\NormalTok{(}\AttributeTok{x =} \ConstantTok{NULL}\NormalTok{, }\AttributeTok{y =} \StringTok{"Nota IMDB"}\NormalTok{) }\SpecialCharTok{+}
  \FunctionTok{theme}\NormalTok{(}\AttributeTok{legend.position =} \StringTok{"none"}\NormalTok{)}
\end{Highlighting}
\end{Shaded}

\pandocbounded{\includegraphics[keepaspectratio]{04-dias/d1-tarde_files/figure-pdf/unnamed-chunk-20-1.pdf}}

\begin{tcolorbox}[enhanced jigsaw, arc=.35mm, bottomrule=.15mm, bottomtitle=1mm, breakable, colback=white, colbacktitle=quarto-callout-tip-color!10!white, colframe=quarto-callout-tip-color-frame, coltitle=black, left=2mm, leftrule=.75mm, opacityback=0, opacitybacktitle=0.6, rightrule=.15mm, title=\textcolor{quarto-callout-tip-color}{\faLightbulb}\hspace{0.5em}{Você acaba de fazer análise exploratória!}, titlerule=0mm, toprule=.15mm, toptitle=1mm]

Em menos de 10 linhas de código você filtrou, agrupou, resumiu e
visualizou dados. Isso é o poder do R + tidyverse. O resto do curso é
aprofundamento dessas mesmas ferramentas.

\end{tcolorbox}

\begin{center}\rule{0.5\linewidth}{0.5pt}\end{center}

\section{Exercícios}\label{exercuxedcios-8}

\begin{tcolorbox}[enhanced jigsaw, arc=.35mm, bottomrule=.15mm, bottomtitle=1mm, breakable, colback=white, colbacktitle=quarto-callout-note-color!10!white, colframe=quarto-callout-note-color-frame, coltitle=black, left=2mm, leftrule=.75mm, opacityback=0, opacitybacktitle=0.6, rightrule=.15mm, title=\textcolor{quarto-callout-note-color}{\faInfo}\hspace{0.5em}{Importante}, titlerule=0mm, toprule=.15mm, toptitle=1mm]

Em todos os exercícios, carregue o dataset com:\\
\texttt{filmes\ \textless{}-\ read\_csv2("../data/filmes\_imdb.csv")}

\end{tcolorbox}

\subsection{Exercício 1}\label{exercuxedcio-1-4}

Filtre os filmes com nota maior ou igual a 8.5. Quantos filmes atendem a
esse critério?

\textbf{Resultado esperado:}

\begin{verbatim}
# A tibble: 1 × 1
      n
  <int>
1   103
\end{verbatim}

\subsection{Exercício 2}\label{exercuxedcio-2-4}

Crie uma coluna \texttt{lucro\ =\ faturamento\ -\ orcamento} e encontre
os 5 filmes mais lucrativos.

\textbf{Resultado esperado:}

\begin{verbatim}
# A tibble: 5 × 2
  titulo                        lucro
  <chr>                         <dbl>
1 The story of cap and trade 48699922
2 Sholay                     18599922
3 Ram Teri Ganga Maili       12599881
4 ET                          8999901
5 Black Friday                1999989
\end{verbatim}

\subsection{Exercício 3}\label{exercuxedcio-3-3}

Calcule a nota média e o número de filmes por gênero. Ordene da maior
para a menor nota média.

\textbf{Resultado esperado:}

\begin{verbatim}
# A tibble: 19 × 3
   genero          nota_media     n
   <chr>                <dbl> <int>
 1 Documentário          6.84    70
 2 Guerra                6.75    19
 3 História              6.73    21
 4 Drama                 6.56   964
 5 Musical               6.56    13
 6 Biografia             6.54    90
 7 Crime                 6.49   308
 8 Mistério              6.43    90
 9 Esporte               6.42     5
10 Aventura              6.42    80
# … with 9 more rows
\end{verbatim}

\subsection{Exercício 4}\label{exercuxedcio-4-2}

Crie um histograma da coluna \texttt{duracao} (duração dos filmes em
minutos). Use 40 bins.

\textbf{Resultado esperado:}

\emph{Um histograma com pico em torno de 90-100 minutos, distribuição
assimétrica à direita.}

\subsection{Exercício 5}\label{exercuxedcio-5}

Crie um gráfico de dispersão com \texttt{orcamento} no eixo x,
\texttt{nota} no eixo y. Existe relação visível?

\textbf{Resultado esperado:}

\emph{Nuvem de pontos sem padrão claro --- orçamento não determina a
qualidade do filme.}

\subsection{Exercício 6}\label{exercuxedcio-6}

Crie um gráfico de barras com os 10 diretores que mais aparecem no
dataset.

\textbf{Resultado esperado:}

\emph{Gráfico de barras horizontal com Woody Allen no topo, seguido por
Steven Spielberg, Clint Eastwood, etc.}

\subsection{Exercício 7 (Desafio)}\label{exercuxedcio-7-desafio}

Combine \texttt{filter()}, \texttt{mutate()}, \texttt{group\_by()} e
\texttt{summarise()} para responder: qual gênero tem o maior orçamento
médio entre filmes com nota ≥ 7.0?

\textbf{Resultado esperado:}

\begin{verbatim}
# A tibble: 10 × 3
   genero          orcamento_medio     n
   <chr>                     <dbl> <int>
 1 Aventura              22980338.     9
 2 Animação              19091918.     9
 3 Fantasia              16110251.     4
 4 Ação                  15915052.    37
 5 Ficção científica     14880336.    21
 6 Guerra                11377778.     2
 7 Drama                 11075548.    91
 8 Comédia                9832855.    46
 9 Romance                9465063.    10
10 Thriller               9416381.    12
\end{verbatim}

\chapter{Dia 2 --- Manhã}\label{dia-2-manhuxe3-1}

\textbf{Dia 2 --- Manhã \textbar{} 4 horas}

\section*{Objetivos}\label{objetivos-2}
\addcontentsline{toc}{section}{Objetivos}

\markright{Objetivos}

Ao final deste capítulo, você será capaz de:

\begin{itemize}
\item
  Criar colunas condicionais com \texttt{case\_when()} e
  \texttt{if\_else()}
\item
  Trabalhar com operações de janela (top N, valores defasados)
\item
  Cruzar tabelas com os diferentes tipos de join
\item
  Transitar entre dados de filmes e dados florestais com as mesmas
  técnicas
\end{itemize}

\begin{center}\rule{0.5\linewidth}{0.5pt}\end{center}

\section{Manipulação condicional}\label{manipulauxe7uxe3o-condicional}

Muitas vezes você precisa criar colunas cujo valor \textbf{depende de
uma condição}. Para isso usamos \texttt{case\_when()} (múltiplas
condições) e \texttt{if\_else()} (condição binária).

\subsection{\texorpdfstring{\texttt{case\_when()} --- múltiplas
condições}{case\_when() --- múltiplas condições}}\label{case_when-muxfaltiplas-condiuxe7uxf5es}

\begin{Shaded}
\begin{Highlighting}[]
\CommentTok{\# Classificar filmes por faixa de nota}
\NormalTok{filmes }\SpecialCharTok{\%\textgreater{}\%} 
  \FunctionTok{mutate}\NormalTok{(}
    \AttributeTok{categoria =} \FunctionTok{case\_when}\NormalTok{(}
\NormalTok{      nota }\SpecialCharTok{\textgreater{}=} \FloatTok{8.5} \SpecialCharTok{\textasciitilde{}} \StringTok{"Excelente"}\NormalTok{,}
\NormalTok{      nota }\SpecialCharTok{\textgreater{}=} \FloatTok{7.0} \SpecialCharTok{\textasciitilde{}} \StringTok{"Bom"}\NormalTok{,}
\NormalTok{      nota }\SpecialCharTok{\textgreater{}=} \FloatTok{5.0} \SpecialCharTok{\textasciitilde{}} \StringTok{"Regular"}\NormalTok{,}
      \ConstantTok{TRUE} \SpecialCharTok{\textasciitilde{}} \StringTok{"Ruim"}
\NormalTok{    )}
\NormalTok{  ) }\SpecialCharTok{\%\textgreater{}\%} 
  \FunctionTok{count}\NormalTok{(categoria)}
\end{Highlighting}
\end{Shaded}

\begin{verbatim}
# A tibble: 4 x 2
  categoria     n
  <chr>     <int>
1 Bom        1277
2 Excelente    47
3 Regular    2242
4 Ruim        309
\end{verbatim}

\begin{Shaded}
\begin{Highlighting}[]
\CommentTok{\# Mesma lógica aplicada a dados florestais}
\NormalTok{inv\_florestal }\SpecialCharTok{\%\textgreater{}\%} 
  \FunctionTok{mutate}\NormalTok{(}
    \AttributeTok{classe\_dap =} \FunctionTok{case\_when}\NormalTok{(}
\NormalTok{      dap }\SpecialCharTok{\textgreater{}=} \DecValTok{35} \SpecialCharTok{\textasciitilde{}} \StringTok{"Grossa"}\NormalTok{,}
\NormalTok{      dap }\SpecialCharTok{\textgreater{}=} \DecValTok{15} \SpecialCharTok{\textasciitilde{}} \StringTok{"Média"}\NormalTok{,}
      \ConstantTok{TRUE} \SpecialCharTok{\textasciitilde{}} \StringTok{"Fina"}
\NormalTok{    )}
\NormalTok{  ) }\SpecialCharTok{\%\textgreater{}\%} 
  \FunctionTok{count}\NormalTok{(classe\_dap)}
\end{Highlighting}
\end{Shaded}

\begin{verbatim}
# A tibble: 3 x 2
  classe_dap     n
  <chr>      <int>
1 Fina         232
2 Grossa        35
3 Média        864
\end{verbatim}

\subsection{\texorpdfstring{\texttt{if\_else()} --- condição
binária}{if\_else() --- condição binária}}\label{if_else-condiuxe7uxe3o-binuxe1ria}

\begin{Shaded}
\begin{Highlighting}[]
\CommentTok{\# Classificar como "Sucesso" se faturamento \textgreater{} orçamento}
\NormalTok{filmes }\SpecialCharTok{\%\textgreater{}\%} 
  \FunctionTok{mutate}\NormalTok{(}
    \AttributeTok{sucesso =} \FunctionTok{if\_else}\NormalTok{(faturamento }\SpecialCharTok{\textgreater{}}\NormalTok{ orcamento, }\StringTok{"Sim"}\NormalTok{, }\StringTok{"Não"}\NormalTok{)}
\NormalTok{  ) }\SpecialCharTok{\%\textgreater{}\%} 
  \FunctionTok{count}\NormalTok{(sucesso)}
\end{Highlighting}
\end{Shaded}

\begin{verbatim}
# A tibble: 2 x 2
  sucesso     n
  <chr>   <int>
1 Não      1830
2 Sim      2045
\end{verbatim}

\begin{center}\rule{0.5\linewidth}{0.5pt}\end{center}

\section{Operações com janelas}\label{operauxe7uxf5es-com-janelas}

Operações de janela permitem comparar cada linha com outras do mesmo
grupo, sem perder os dados originais.

\subsection{Top N por grupo}\label{top-n-por-grupo}

\begin{Shaded}
\begin{Highlighting}[]
\CommentTok{\# Top 3 filmes por gênero (por nota)}
\NormalTok{filmes }\SpecialCharTok{\%\textgreater{}\%} 
  \FunctionTok{group\_by}\NormalTok{(genero) }\SpecialCharTok{\%\textgreater{}\%} 
  \FunctionTok{slice\_max}\NormalTok{(nota, }\AttributeTok{n =} \DecValTok{3}\NormalTok{) }\SpecialCharTok{\%\textgreater{}\%} 
  \FunctionTok{select}\NormalTok{(genero, titulo, nota) }\SpecialCharTok{\%\textgreater{}\%} 
  \FunctionTok{arrange}\NormalTok{(genero, }\FunctionTok{desc}\NormalTok{(nota))}
\end{Highlighting}
\end{Shaded}

\begin{verbatim}
# A tibble: 64 x 3
# Groups:   genero [19]
   genero   titulo                                              nota
   <chr>    <chr>                                              <dbl>
 1 Animação Schindler's List                                     8.9
 2 Animação The Lord of the Rings: The Fellowship of the Ring    8.8
 3 Animação The Lord of the Rings: The Two Towers                8.7
 4 Animação One Flew Over the Cuckoo's Nest                      8.7
 5 Aventura The Dark Knight                                      9  
 6 Aventura Braveheart                                           8.4
 7 Aventura A Separation                                         8.4
 8 Ação     Inception                                            8.8
 9 Ação     Goodfellas                                           8.7
10 Ação     The Usual Suspects                                   8.6
# i 54 more rows
\end{verbatim}

\begin{Shaded}
\begin{Highlighting}[]
\CommentTok{\# 5 árvores com maior DAP em cada parcela (dado florestal)}
\NormalTok{inv\_florestal }\SpecialCharTok{\%\textgreater{}\%} 
  \FunctionTok{group\_by}\NormalTok{(parcela) }\SpecialCharTok{\%\textgreater{}\%} 
  \FunctionTok{slice\_max}\NormalTok{(dap, }\AttributeTok{n =} \DecValTok{5}\NormalTok{) }\SpecialCharTok{\%\textgreater{}\%} 
  \FunctionTok{select}\NormalTok{(parcela, especie, dap, altura)}
\end{Highlighting}
\end{Shaded}

\begin{verbatim}
# A tibble: 150 x 4
# Groups:   parcela [30]
   parcela especie                dap altura
     <dbl> <chr>                <dbl>  <dbl>
 1       1 Eucalyptus Grandis    37.2   54.2
 2       1 Pinus Taeda           33.8   54.4
 3       1 Eucalyptus Urophylla  32.6   52.1
 4       1 Eucalyptus Grandis    31     52.4
 5       1 Eucalyptus Grandis    30.9   50.7
 6       2 Eucalyptus Grandis    36.9   56.5
 7       2 Eucalyptus Grandis    34.9   56.5
 8       2 Eucalyptus Grandis    34     54.3
 9       2 Eucalyptus Grandis    32.7   51.2
10       2 Eucalyptus Grandis    30.4   51.5
# i 140 more rows
\end{verbatim}

\subsection{\texorpdfstring{Valores defasados: \texttt{lag()} e
\texttt{lead()}}{Valores defasados: lag() e lead()}}\label{valores-defasados-lag-e-lead}

\begin{Shaded}
\begin{Highlighting}[]
\CommentTok{\# Comparar faturamento entre filmes consecutivos do mesmo diretor}
\NormalTok{filmes }\SpecialCharTok{\%\textgreater{}\%} 
  \FunctionTok{group\_by}\NormalTok{(diretor) }\SpecialCharTok{\%\textgreater{}\%} 
  \FunctionTok{arrange}\NormalTok{(ano) }\SpecialCharTok{\%\textgreater{}\%} 
  \FunctionTok{mutate}\NormalTok{(}
    \AttributeTok{faturamento\_anterior =} \FunctionTok{lag}\NormalTok{(faturamento),}
    \AttributeTok{variacao =}\NormalTok{ faturamento }\SpecialCharTok{{-}}\NormalTok{ faturamento\_anterior}
\NormalTok{  ) }\SpecialCharTok{\%\textgreater{}\%} 
  \FunctionTok{select}\NormalTok{(diretor, ano, titulo, faturamento, faturamento\_anterior, variacao) }\SpecialCharTok{\%\textgreater{}\%} 
  \FunctionTok{filter}\NormalTok{(}\SpecialCharTok{!}\FunctionTok{is.na}\NormalTok{(variacao))}
\end{Highlighting}
\end{Shaded}

\begin{verbatim}
# A tibble: 2,130 x 6
# Groups:   diretor [735]
   diretor                ano titulo   faturamento faturamento_anterior variacao
   <chr>                <dbl> <chr>          <dbl>                <dbl>    <dbl>
 1 Victor Fleming        1939 The Wiz~    22202612            198655278  -1.76e8
 2 David Lean            1962 Lawrenc~     6000000             27200000  -2.12e7
 3 Terence Young         1963 From Ru~    24800000             16067035   8.73e6
 4 David Lean            1965 Doctor ~   111722000              6000000   1.06e8
 5 Terence Young         1965 Thunder~    63600000             24800000   3.88e7
 6 Sergio Leone          1966 The Goo~     6100000              3500000   2.6 e6
 7 Guy Hamilton          1971 Diamond~    43800000             51100000  -7.3 e6
 8 Guy Hamilton          1973 Live an~    35400000             43800000  -8.40e6
 9 George Roy Hill       1973 The Sti~   159600000            102308900   5.73e7
10 Francis Ford Coppola  1974 The God~    57300000            134821952  -7.75e7
# i 2,120 more rows
\end{verbatim}

\subsection{Operações cumulativas}\label{operauxe7uxf5es-cumulativas}

\begin{Shaded}
\begin{Highlighting}[]
\CommentTok{\# Soma cumulativa de faturamento por ano}
\NormalTok{filmes }\SpecialCharTok{\%\textgreater{}\%} 
  \FunctionTok{group\_by}\NormalTok{(ano) }\SpecialCharTok{\%\textgreater{}\%} 
  \FunctionTok{arrange}\NormalTok{(faturamento) }\SpecialCharTok{\%\textgreater{}\%} 
  \FunctionTok{mutate}\NormalTok{(}
    \AttributeTok{faturamento\_cumulativo =} \FunctionTok{cumsum}\NormalTok{(faturamento)}
\NormalTok{  ) }\SpecialCharTok{\%\textgreater{}\%} 
  \FunctionTok{select}\NormalTok{(ano, titulo, faturamento, faturamento\_cumulativo)}
\end{Highlighting}
\end{Shaded}

\begin{verbatim}
# A tibble: 3,875 x 4
# Groups:   ano [75]
     ano titulo                            faturamento faturamento_cumulativo
   <dbl> <chr>                                   <dbl>                  <dbl>
 1  2014 Skin Trade                                162                    162
 2  2001 The Jimmy Show                            703                    703
 3  2006 In Her Line of Fire                       721                    721
 4  2006 The Trials of Darryl Hunt                1111                   1832
 5  2012 Detention of the Dead                    1332                   1332
 6  2012 Eddie: The Sleepwalking Cannibal         1521                   2853
 7  2011 The Harvest/La Cosecha                   2245                   2245
 8  2012 The Ghastly Love of Johnny X             2436                   5289
 9  2014 Locker 13                                2468                   2630
10  2003 I Love Your Work                         2580                   2580
# i 3,865 more rows
\end{verbatim}

\begin{Shaded}
\begin{Highlighting}[]
\CommentTok{\# Volume acumulado por parcela (dado florestal)}
\NormalTok{inv\_florestal }\SpecialCharTok{\%\textgreater{}\%} 
  \FunctionTok{group\_by}\NormalTok{(parcela) }\SpecialCharTok{\%\textgreater{}\%} 
  \FunctionTok{mutate}\NormalTok{(}\AttributeTok{volume\_cumulativo =} \FunctionTok{cumsum}\NormalTok{(volume)) }\SpecialCharTok{\%\textgreater{}\%} 
  \FunctionTok{select}\NormalTok{(parcela, especie, volume, volume\_cumulativo)}
\end{Highlighting}
\end{Shaded}

\begin{verbatim}
# A tibble: 1,131 x 4
# Groups:   parcela [30]
   parcela especie            volume volume_cumulativo
     <dbl> <chr>               <dbl>             <dbl>
 1       1 Eucalyptus Grandis  0.876             0.876
 2       1 Pinus Taeda         0.699             1.57 
 3       1 Eucalyptus Grandis  1.24              2.81 
 4       1 Pinus Taeda         0.217             3.03 
 5       1 Eucalyptus Grandis  0.746             3.77 
 6       1 Eucalyptus Grandis  0.596             4.37 
 7       1 Pinus Taeda         1.02              5.39 
 8       1 Eucalyptus Grandis  1.56              6.95 
 9       1 Eucalyptus Grandis  1.11              8.06 
10       1 Pinus Taeda         1.57              9.62 
# i 1,121 more rows
\end{verbatim}

\begin{center}\rule{0.5\linewidth}{0.5pt}\end{center}

\section{Trabalhando com múltiplas tabelas:
Joins}\label{trabalhando-com-muxfaltiplas-tabelas-joins}

Na vida real, os dados raramente estão em uma única tabela. Informações
complementares ficam em tabelas separadas que precisam ser
\textbf{cruzadas}.

\subsection{Os tipos de join}\label{os-tipos-de-join}

\begin{verbatim}
left_join()   → Mantém TODAS as linhas da tabela da esquerda
inner_join()  → Mantém APENAS linhas com correspondência nas duas
full_join()   → Mantém TODAS as linhas de ambas
anti_join()   → Mostra o que NÃO tem correspondência (diagnóstico)
\end{verbatim}

\subsection{Exemplo prático: filmes +
elenco}\label{exemplo-pruxe1tico-filmes-elenco}

\begin{Shaded}
\begin{Highlighting}[]
\CommentTok{\# Carregar tabelas separadas}
\NormalTok{elenco }\OtherTok{\textless{}{-}} \FunctionTok{read\_csv}\NormalTok{(}\StringTok{"../data/elenco.csv"}\NormalTok{)}

\CommentTok{\# left\_join: adicionar informações do elenco aos filmes}
\NormalTok{filmes\_com\_elenco }\OtherTok{\textless{}{-}}\NormalTok{ filmes }\SpecialCharTok{\%\textgreater{}\%} 
  \FunctionTok{left\_join}\NormalTok{(elenco, }\AttributeTok{by =} \FunctionTok{c}\NormalTok{(}\StringTok{"codigo\_filme"} \OtherTok{=} \StringTok{"codigo"}\NormalTok{))}

\FunctionTok{glimpse}\NormalTok{(filmes\_com\_elenco)}
\end{Highlighting}
\end{Shaded}

\begin{verbatim}
Rows: 3,875
Columns: 12
$ codigo_filme     <dbl> 1, 2, 3, 4, 5, 6, 7, 8, 9, 10, 11, 12, 13, 14, 15, 16~
$ titulo           <chr> "Avatar ", "Pirates of the Caribbean: At World's End ~
$ diretor          <chr> "James Cameron", "Gore Verbinski", "Sam Mendes", "Chr~
$ ano              <dbl> 2009, 2007, 2015, 2012, 2012, 2007, 2010, 2015, 2009,~
$ genero           <chr> "Ação", "Aventura", "Comédia", "Mistério", "Fantasia"~
$ duracao          <dbl> 81, 126, 88, 99, 92, 95, 39, 111, 82, 89, 77, 102, 98~
$ orcamento        <dbl> 237000000, 300000000, 245000000, 250000000, 263700000~
$ faturamento      <dbl> 760505847, 309404152, 200074175, 448130642, 73058679,~
$ nota             <dbl> 7.9, 7.1, 6.8, 8.5, 6.6, 6.2, 7.8, 7.5, 7.5, 6.9, 6.1~
$ elenco_principal <chr> "Viola Davis", "Cate Blanchett", "Natalie Portman", "~
$ ator_principal   <chr> "Viola Davis", "Viola Davis", "Tom Hanks", "Leonardo ~
$ orcamento_elenco <dbl> 11783666, 30312077, 47499828, 27114606, 38641928, 143~
\end{verbatim}

\subsection{Exemplo florestal: inventário +
cadastro}\label{exemplo-florestal-inventuxe1rio-cadastro}

\begin{Shaded}
\begin{Highlighting}[]
\CommentTok{\# Dados de árvores (inv\_florestal) + informações das parcelas (parcelas)}
\NormalTok{parcelas }\OtherTok{\textless{}{-}} \FunctionTok{read\_csv}\NormalTok{(}\StringTok{"../data/parcelas.csv"}\NormalTok{)}

\NormalTok{inv\_completo }\OtherTok{\textless{}{-}}\NormalTok{ inv\_florestal }\SpecialCharTok{\%\textgreater{}\%} 
  \FunctionTok{left\_join}\NormalTok{(parcelas, }\AttributeTok{by =} \StringTok{"parcela"}\NormalTok{)}

\FunctionTok{glimpse}\NormalTok{(inv\_completo)}
\end{Highlighting}
\end{Shaded}

\begin{verbatim}
Rows: 1,131
Columns: 10
$ parcela      <dbl> 1, 1, 1, 1, 1, 1, 1, 1, 1, 1, 1, 1, 1, 1, 1, 1, 1, 1, 1, ~
$ arvore       <dbl> 1, 2, 3, 4, 5, 6, 7, 8, 9, 10, 11, 12, 13, 14, 15, 16, 17~
$ especie      <chr> "Eucalyptus Grandis", "Pinus Taeda", "Eucalyptus Grandis"~
$ dap          <dbl> 21.7, 20.5, 26.2, 12.6, 19.5, 17.9, 23.5, 27.0, 24.5, 27.~
$ altura       <dbl> 46.0, 44.0, 51.1, 38.7, 49.8, 44.3, 49.8, 52.9, 48.9, 50.~
$ volume       <dbl> 0.8759, 0.6989, 1.2358, 0.2173, 0.7461, 0.5956, 1.0173, 1~
$ talhao       <dbl> 9, 9, 9, 9, 9, 9, 9, 9, 9, 9, 9, 9, 9, 9, 9, 9, 9, 9, 9, ~
$ area_parcela <dbl> 600, 600, 600, 600, 600, 600, 600, 600, 600, 600, 600, 60~
$ idade        <dbl> 11, 11, 11, 11, 11, 11, 11, 11, 11, 11, 11, 11, 11, 11, 1~
$ sitio        <chr> "Médio", "Médio", "Médio", "Médio", "Médio", "Médio", "Mé~
\end{verbatim}

\begin{Shaded}
\begin{Highlighting}[]
\CommentTok{\# Agora podemos calcular volume por hectare!}
\NormalTok{inv\_completo }\SpecialCharTok{\%\textgreater{}\%} 
  \FunctionTok{group\_by}\NormalTok{(parcela, idade, sitio) }\SpecialCharTok{\%\textgreater{}\%} 
  \FunctionTok{summarise}\NormalTok{(}
    \AttributeTok{volume\_total =} \FunctionTok{sum}\NormalTok{(volume),}
    \AttributeTok{area\_ha =} \FunctionTok{first}\NormalTok{(area\_parcela) }\SpecialCharTok{/} \DecValTok{10000}\NormalTok{,}
    \AttributeTok{volume\_ha =}\NormalTok{ volume\_total }\SpecialCharTok{/}\NormalTok{ area\_ha}
\NormalTok{  ) }\SpecialCharTok{\%\textgreater{}\%} 
  \FunctionTok{arrange}\NormalTok{(}\FunctionTok{desc}\NormalTok{(volume\_ha))}
\end{Highlighting}
\end{Shaded}

\begin{verbatim}
# A tibble: 30 x 6
# Groups:   parcela, idade [30]
   parcela idade sitio volume_total area_ha volume_ha
     <dbl> <dbl> <chr>        <dbl>   <dbl>     <dbl>
 1      22     5 Ruim          50.4    0.06      839.
 2      18     5 Médio         49.7    0.06      828.
 3      13    11 Médio         44.4    0.06      740.
 4       5    11 Médio         43.0    0.06      717.
 5       6     8 Médio         42.2    0.06      704.
 6      11    12 Médio         41.9    0.06      698.
 7      17     5 Médio         41.7    0.06      695.
 8       2     9 Bom           40.1    0.06      668.
 9      25     9 Médio         39.5    0.06      658.
10       4     6 Bom           39.1    0.06      652.
# i 20 more rows
\end{verbatim}

\subsection{\texorpdfstring{Diagnóstico com
\texttt{anti\_join()}}{Diagnóstico com anti\_join()}}\label{diagnuxf3stico-com-anti_join}

\begin{Shaded}
\begin{Highlighting}[]
\CommentTok{\# Parcelas que estão no inventário mas NÃO no cadastro}
\NormalTok{inv\_florestal }\SpecialCharTok{\%\textgreater{}\%} 
  \FunctionTok{anti\_join}\NormalTok{(parcelas, }\AttributeTok{by =} \StringTok{"parcela"}\NormalTok{) }\SpecialCharTok{\%\textgreater{}\%} 
  \FunctionTok{distinct}\NormalTok{(parcela)}
\end{Highlighting}
\end{Shaded}

\begin{verbatim}
# A tibble: 0 x 1
# i 1 variable: parcela <dbl>
\end{verbatim}

\begin{Shaded}
\begin{Highlighting}[]
\CommentTok{\# Parcelas no cadastro mas SEM árvores no inventário}
\NormalTok{parcelas }\SpecialCharTok{\%\textgreater{}\%} 
  \FunctionTok{anti\_join}\NormalTok{(inv\_florestal, }\AttributeTok{by =} \StringTok{"parcela"}\NormalTok{)}
\end{Highlighting}
\end{Shaded}

\begin{verbatim}
# A tibble: 0 x 5
# i 5 variables: parcela <dbl>, talhao <dbl>, area_parcela <dbl>, idade <dbl>,
#   sitio <chr>
\end{verbatim}

\subsection{\texorpdfstring{\texttt{bind\_rows()} --- Empilhando
tabelas}{bind\_rows() --- Empilhando tabelas}}\label{bind_rows-empilhando-tabelas}

\begin{Shaded}
\begin{Highlighting}[]
\CommentTok{\# Simulando inventários de anos diferentes}
\NormalTok{inv\_2022 }\OtherTok{\textless{}{-}}\NormalTok{ inv\_florestal }\SpecialCharTok{\%\textgreater{}\%} \FunctionTok{mutate}\NormalTok{(}\AttributeTok{ano\_medicao =} \DecValTok{2022}\NormalTok{)}
\NormalTok{inv\_2023 }\OtherTok{\textless{}{-}}\NormalTok{ inv\_florestal }\SpecialCharTok{\%\textgreater{}\%} \FunctionTok{mutate}\NormalTok{(}\AttributeTok{ano\_medicao =} \DecValTok{2023}\NormalTok{, }\AttributeTok{dap =}\NormalTok{ dap }\SpecialCharTok{+} \FunctionTok{rnorm}\NormalTok{(}\FunctionTok{n}\NormalTok{(), }\DecValTok{2}\NormalTok{, }\DecValTok{1}\NormalTok{))}

\NormalTok{inv\_historico }\OtherTok{\textless{}{-}} \FunctionTok{bind\_rows}\NormalTok{(inv\_2022, inv\_2023)}

\NormalTok{inv\_historico }\SpecialCharTok{\%\textgreater{}\%} 
  \FunctionTok{group\_by}\NormalTok{(ano\_medicao) }\SpecialCharTok{\%\textgreater{}\%} 
  \FunctionTok{summarise}\NormalTok{(}\AttributeTok{dap\_medio =} \FunctionTok{mean}\NormalTok{(dap))}
\end{Highlighting}
\end{Shaded}

\begin{verbatim}
# A tibble: 2 x 2
  ano_medicao dap_medio
        <dbl>     <dbl>
1        2022      20.7
2        2023      22.7
\end{verbatim}

\begin{center}\rule{0.5\linewidth}{0.5pt}\end{center}

\section{Transição: de filmes para
floresta}\label{transiuxe7uxe3o-de-filmes-para-floresta}

A beleza do que você aprendeu é que as \textbf{mesmas técnicas} se
aplicam a qualquer domínio. Compare:

\begin{Shaded}
\begin{Highlighting}[]
\CommentTok{\# Análise com filmes}
\NormalTok{filmes }\SpecialCharTok{\%\textgreater{}\%} 
  \FunctionTok{group\_by}\NormalTok{(genero) }\SpecialCharTok{\%\textgreater{}\%} 
  \FunctionTok{summarise}\NormalTok{(}\AttributeTok{media =} \FunctionTok{mean}\NormalTok{(nota), }\AttributeTok{n =} \FunctionTok{n}\NormalTok{()) }\SpecialCharTok{\%\textgreater{}\%} 
  \FunctionTok{arrange}\NormalTok{(}\FunctionTok{desc}\NormalTok{(media))}
\end{Highlighting}
\end{Shaded}

\begin{verbatim}
# A tibble: 19 x 3
   genero            media     n
   <chr>             <dbl> <int>
 1 Musical            6.75    39
 2 Esporte            6.73    16
 3 História           6.56    31
 4 Guerra             6.55    33
 5 Biografia          6.50   121
 6 Comédia            6.50   452
 7 Ação               6.49   566
 8 Drama              6.48   719
 9 Romance            6.48   245
10 Documentário       6.47   117
11 Ficção científica  6.47   306
12 Fantasia           6.47    82
13 Mistério           6.46   113
14 Aventura           6.45   204
15 Thriller           6.41   185
16 Terror             6.40   311
17 Crime              6.39   193
18 Animação           6.29   114
19 Família            6.1     28
\end{verbatim}

\begin{Shaded}
\begin{Highlighting}[]
\CommentTok{\# Mesma análise com dados florestais}
\NormalTok{inv\_florestal }\SpecialCharTok{\%\textgreater{}\%} 
  \FunctionTok{group\_by}\NormalTok{(especie) }\SpecialCharTok{\%\textgreater{}\%} 
  \FunctionTok{summarise}\NormalTok{(}\AttributeTok{media\_dap =} \FunctionTok{mean}\NormalTok{(dap), }\AttributeTok{n =} \FunctionTok{n}\NormalTok{()) }\SpecialCharTok{\%\textgreater{}\%} 
  \FunctionTok{arrange}\NormalTok{(}\FunctionTok{desc}\NormalTok{(media\_dap))}
\end{Highlighting}
\end{Shaded}

\begin{verbatim}
# A tibble: 3 x 3
  especie              media_dap     n
  <chr>                    <dbl> <int>
1 Eucalyptus Grandis        22.1   682
2 Eucalyptus Urophylla      21.7    99
3 Pinus Taeda               17.6   350
\end{verbatim}

A estrutura é idêntica. Só mudam os nomes das colunas e o contexto!

\begin{center}\rule{0.5\linewidth}{0.5pt}\end{center}

\section{Exercícios}\label{exercuxedcios-9}

\subsection{Exercício 1}\label{exercuxedcio-1-5}

Usando \texttt{case\_when()}, classifique as árvores do
\texttt{inv\_florestal} em classes de DAP: ``Fina'' (\textless{} 15 cm),
``Média'' (15-25 cm), ``Grossa'' (\textgreater{} 25 cm). Conte quantas
árvores em cada classe.

\textbf{Resultado esperado:}

\begin{verbatim}
# A tibble: 3 × 2
  classe_dap     n
  <chr>      <int>
1 Grossa        96
2 Média        401
3 Fina         495
\end{verbatim}

\subsection{Exercício 2}\label{exercuxedcio-2-5}

Usando \texttt{slice\_max()}, encontre as 3 maiores árvores (por DAP) de
cada espécie.

\textbf{Resultado esperado:}

\begin{verbatim}
# A tibble: 9 × 3
  especie    dap altura
  <chr>    <dbl>  <dbl>
1 Eucalipto  42.3   32.1
2 Eucalipto  41.8   31.5
3 Eucalipto  40.2   30.8
4 Pinus      38.7   28.4
5 Pinus      37.2   27.8
6 Pinus      36.5   27.1
# … with 3 more rows
\end{verbatim}

\subsection{Exercício 3}\label{exercuxedcio-3-4}

Faça um \texttt{left\_join} entre \texttt{inv\_florestal} e
\texttt{parcelas}. Agrupe por \texttt{parcela} e calcule: volume total,
área em hectares, volume por hectare.

\textbf{Resultado esperado (top 3 parcelas):}

\begin{verbatim}
# A tibble: 30 × 5
   parcela volume_total area_ha volume_ha
     <dbl>        <dbl>   <dbl>     <dbl>
 1       7        25.2    0.06      420. 
 2       2        24.3    0.06      405. 
 3      30        23.8    0.06      397. 
\end{verbatim}

\subsection{Exercício 4}\label{exercuxedcio-4-3}

Calcule o IMA (Incremento Médio Anual) = volume\_ha / idade para cada
parcela. Ordene do maior para o menor.

\textbf{Resultado esperado (top 3):}

\begin{verbatim}
# A tibble: 30 × 4
  parcela volume_ha idade   IMA
    <dbl>     <dbl> <dbl> <dbl>
1       7      420.     5  84.1
2      12      395.     5  79.1
3      15      389.     7  55.6
\end{verbatim}

\subsection{Exercício 5}\label{exercuxedcio-5-1}

Use \texttt{anti\_join()} para verificar se há alguma parcela no
\texttt{inv\_florestal} que não existe no cadastro de \texttt{parcelas}.

\textbf{Resultado esperado:}

\begin{verbatim}
# A tibble: 0 × 1
# … with 1 variable: parcela
\end{verbatim}

\emph{(Zero linhas = tudo certo, todas as parcelas têm cadastro.)}

\chapter{Dia 2 --- Tarde}\label{dia-2-tarde-1}

\textbf{Dia 2 --- Tarde \textbar{} 4 horas}

\section*{Objetivos}\label{objetivos-3}
\addcontentsline{toc}{section}{Objetivos}

\markright{Objetivos}

Ao final deste capítulo, você será capaz de:

\begin{itemize}
\item
  Entender a gramática de gráficos do ggplot2
\item
  Criar visualizações profissionais com customização completa
\item
  Usar facetas para comparar grupos
\item
  Exportar gráficos em alta resolução
\end{itemize}

\begin{center}\rule{0.5\linewidth}{0.5pt}\end{center}

\section{A gramática de gráficos}\label{a-gramuxe1tica-de-gruxe1ficos}

O ggplot2 é baseado na \textbf{Gramática de Gráficos} (\emph{Grammar of
Graphics}). Todo gráfico é construído por camadas:

\begin{verbatim}
ggplot(dados, aes(x, y)) +   # dados + mapeamento estético
  geom_tipo() +               # geometria (o que desenhar)
  scale_tipo() +              # escalas (cores, eixos)
  theme_tipo() +              # tema (aparência)
  labs(...)                   # rótulos (títulos)
\end{verbatim}

Vamos construir um gráfico passo a passo, adicionando camadas:

\begin{Shaded}
\begin{Highlighting}[]
\CommentTok{\# Passo 1: dados + eixos (tela vazia)}
\NormalTok{inv\_florestal }\SpecialCharTok{\%\textgreater{}\%}
  \FunctionTok{ggplot}\NormalTok{(}\FunctionTok{aes}\NormalTok{(}\AttributeTok{x =}\NormalTok{ dap, }\AttributeTok{y =}\NormalTok{ altura))}
\end{Highlighting}
\end{Shaded}

\pandocbounded{\includegraphics[keepaspectratio]{04-dias/d2-tarde_files/figure-pdf/unnamed-chunk-1-1.pdf}}

\begin{Shaded}
\begin{Highlighting}[]
\CommentTok{\# Passo 2: geometria de pontos}
\NormalTok{inv\_florestal }\SpecialCharTok{\%\textgreater{}\%}
  \FunctionTok{ggplot}\NormalTok{(}\FunctionTok{aes}\NormalTok{(}\AttributeTok{x =}\NormalTok{ dap, }\AttributeTok{y =}\NormalTok{ altura)) }\SpecialCharTok{+}
  \FunctionTok{geom\_point}\NormalTok{()}
\end{Highlighting}
\end{Shaded}

\pandocbounded{\includegraphics[keepaspectratio]{04-dias/d2-tarde_files/figure-pdf/unnamed-chunk-2-1.pdf}}

\begin{Shaded}
\begin{Highlighting}[]
\CommentTok{\# Passo 3: cor por espécie}
\NormalTok{inv\_florestal }\SpecialCharTok{\%\textgreater{}\%}
  \FunctionTok{ggplot}\NormalTok{(}\FunctionTok{aes}\NormalTok{(}\AttributeTok{x =}\NormalTok{ dap, }\AttributeTok{y =}\NormalTok{ altura, }\AttributeTok{color =}\NormalTok{ especie)) }\SpecialCharTok{+}
  \FunctionTok{geom\_point}\NormalTok{(}\AttributeTok{alpha =} \FloatTok{0.5}\NormalTok{)}
\end{Highlighting}
\end{Shaded}

\pandocbounded{\includegraphics[keepaspectratio]{04-dias/d2-tarde_files/figure-pdf/unnamed-chunk-3-1.pdf}}

\begin{Shaded}
\begin{Highlighting}[]
\CommentTok{\# Passo 4: linha de tendência + tema + rótulos}
\NormalTok{inv\_florestal }\SpecialCharTok{\%\textgreater{}\%}
  \FunctionTok{ggplot}\NormalTok{(}\FunctionTok{aes}\NormalTok{(}\AttributeTok{x =}\NormalTok{ dap, }\AttributeTok{y =}\NormalTok{ altura, }\AttributeTok{color =}\NormalTok{ especie)) }\SpecialCharTok{+}
  \FunctionTok{geom\_point}\NormalTok{(}\AttributeTok{alpha =} \FloatTok{0.5}\NormalTok{) }\SpecialCharTok{+}
  \FunctionTok{geom\_smooth}\NormalTok{(}\AttributeTok{method =} \StringTok{"lm"}\NormalTok{, }\AttributeTok{se =} \ConstantTok{FALSE}\NormalTok{) }\SpecialCharTok{+}
  \FunctionTok{labs}\NormalTok{(}
    \AttributeTok{title =} \StringTok{"Relação hipsométrica"}\NormalTok{,}
    \AttributeTok{subtitle =} \StringTok{"DAP × Altura por espécie"}\NormalTok{,}
    \AttributeTok{x =} \StringTok{"DAP (cm)"}\NormalTok{,}
    \AttributeTok{y =} \StringTok{"Altura (m)"}\NormalTok{,}
    \AttributeTok{color =} \StringTok{"Espécie"}
\NormalTok{  ) }\SpecialCharTok{+}
  \FunctionTok{theme\_minimal}\NormalTok{()}
\end{Highlighting}
\end{Shaded}

\pandocbounded{\includegraphics[keepaspectratio]{04-dias/d2-tarde_files/figure-pdf/unnamed-chunk-4-1.pdf}}

\begin{center}\rule{0.5\linewidth}{0.5pt}\end{center}

\section{Geometrias essenciais}\label{geometrias-essenciais}

\subsection{Uma variável contínua: histograma e
densidade}\label{uma-variuxe1vel-contuxednua-histograma-e-densidade}

\begin{Shaded}
\begin{Highlighting}[]
\CommentTok{\# Distribuição diamétrica — histograma}
\NormalTok{inv\_florestal }\SpecialCharTok{\%\textgreater{}\%}
  \FunctionTok{ggplot}\NormalTok{(}\FunctionTok{aes}\NormalTok{(}\AttributeTok{x =}\NormalTok{ dap)) }\SpecialCharTok{+}
  \FunctionTok{geom\_histogram}\NormalTok{(}\AttributeTok{bins =} \DecValTok{30}\NormalTok{, }\AttributeTok{fill =} \StringTok{"\#18BC9C"}\NormalTok{, }\AttributeTok{color =} \StringTok{"white"}\NormalTok{) }\SpecialCharTok{+}
  \FunctionTok{labs}\NormalTok{(}\AttributeTok{title =} \StringTok{"Distribuição diamétrica"}\NormalTok{, }\AttributeTok{x =} \StringTok{"DAP (cm)"}\NormalTok{, }\AttributeTok{y =} \StringTok{"Frequência"}\NormalTok{)}
\end{Highlighting}
\end{Shaded}

\pandocbounded{\includegraphics[keepaspectratio]{04-dias/d2-tarde_files/figure-pdf/unnamed-chunk-5-1.pdf}}

\begin{Shaded}
\begin{Highlighting}[]
\CommentTok{\# Densidade — mais suave que histograma, bom para comparar grupos}
\NormalTok{inv\_florestal }\SpecialCharTok{\%\textgreater{}\%}
  \FunctionTok{ggplot}\NormalTok{(}\FunctionTok{aes}\NormalTok{(}\AttributeTok{x =}\NormalTok{ dap, }\AttributeTok{fill =}\NormalTok{ especie)) }\SpecialCharTok{+}
  \FunctionTok{geom\_density}\NormalTok{(}\AttributeTok{alpha =} \FloatTok{0.5}\NormalTok{) }\SpecialCharTok{+}
  \FunctionTok{labs}\NormalTok{(}\AttributeTok{title =} \StringTok{"Distribuição diamétrica por espécie"}\NormalTok{, }\AttributeTok{x =} \StringTok{"DAP (cm)"}\NormalTok{, }\AttributeTok{y =} \StringTok{"Densidade"}\NormalTok{)}
\end{Highlighting}
\end{Shaded}

\pandocbounded{\includegraphics[keepaspectratio]{04-dias/d2-tarde_files/figure-pdf/unnamed-chunk-6-1.pdf}}

\subsection{Duas variáveis: dispersão, linha,
boxplot}\label{duas-variuxe1veis-dispersuxe3o-linha-boxplot}

\begin{Shaded}
\begin{Highlighting}[]
\CommentTok{\# Dispersão com tendência}
\NormalTok{inv\_florestal }\SpecialCharTok{\%\textgreater{}\%}
  \FunctionTok{ggplot}\NormalTok{(}\FunctionTok{aes}\NormalTok{(}\AttributeTok{x =}\NormalTok{ dap, }\AttributeTok{y =}\NormalTok{ volume)) }\SpecialCharTok{+}
  \FunctionTok{geom\_point}\NormalTok{(}\AttributeTok{alpha =} \FloatTok{0.3}\NormalTok{, }\AttributeTok{color =} \StringTok{"\#2C3E50"}\NormalTok{) }\SpecialCharTok{+}
  \FunctionTok{geom\_smooth}\NormalTok{(}\AttributeTok{method =} \StringTok{"lm"}\NormalTok{, }\AttributeTok{color =} \StringTok{"\#E74C3C"}\NormalTok{) }\SpecialCharTok{+}
  \FunctionTok{labs}\NormalTok{(}\AttributeTok{title =} \StringTok{"Relação DAP × Volume"}\NormalTok{, }\AttributeTok{x =} \StringTok{"DAP (cm)"}\NormalTok{, }\AttributeTok{y =} \FunctionTok{expression}\NormalTok{(Volume}\SpecialCharTok{\textasciitilde{}}\NormalTok{(m}\SpecialCharTok{\^{}}\DecValTok{3}\NormalTok{)))}
\end{Highlighting}
\end{Shaded}

\pandocbounded{\includegraphics[keepaspectratio]{04-dias/d2-tarde_files/figure-pdf/unnamed-chunk-7-1.pdf}}

\begin{Shaded}
\begin{Highlighting}[]
\CommentTok{\# Boxplot — comparar distribuições}
\NormalTok{inv\_florestal }\SpecialCharTok{\%\textgreater{}\%}
  \FunctionTok{ggplot}\NormalTok{(}\FunctionTok{aes}\NormalTok{(}\AttributeTok{x =}\NormalTok{ especie, }\AttributeTok{y =}\NormalTok{ volume, }\AttributeTok{fill =}\NormalTok{ especie)) }\SpecialCharTok{+}
  \FunctionTok{geom\_boxplot}\NormalTok{() }\SpecialCharTok{+}
  \FunctionTok{labs}\NormalTok{(}\AttributeTok{title =} \StringTok{"Volume por espécie"}\NormalTok{, }\AttributeTok{x =} \ConstantTok{NULL}\NormalTok{, }\AttributeTok{y =} \FunctionTok{expression}\NormalTok{(Volume}\SpecialCharTok{\textasciitilde{}}\NormalTok{(m}\SpecialCharTok{\^{}}\DecValTok{3}\NormalTok{))) }\SpecialCharTok{+}
  \FunctionTok{theme}\NormalTok{(}\AttributeTok{legend.position =} \StringTok{"none"}\NormalTok{)}
\end{Highlighting}
\end{Shaded}

\pandocbounded{\includegraphics[keepaspectratio]{04-dias/d2-tarde_files/figure-pdf/unnamed-chunk-8-1.pdf}}

\begin{Shaded}
\begin{Highlighting}[]
\CommentTok{\# Volume total por parcela (barras)}
\NormalTok{inv\_florestal }\SpecialCharTok{\%\textgreater{}\%}
  \FunctionTok{group\_by}\NormalTok{(parcela) }\SpecialCharTok{\%\textgreater{}\%}
  \FunctionTok{summarise}\NormalTok{(}\AttributeTok{volume\_total =} \FunctionTok{sum}\NormalTok{(volume)) }\SpecialCharTok{\%\textgreater{}\%}
  \FunctionTok{ggplot}\NormalTok{(}\FunctionTok{aes}\NormalTok{(}\AttributeTok{x =} \FunctionTok{reorder}\NormalTok{(}\FunctionTok{factor}\NormalTok{(parcela), volume\_total), }\AttributeTok{y =}\NormalTok{ volume\_total)) }\SpecialCharTok{+}
  \FunctionTok{geom\_col}\NormalTok{(}\AttributeTok{fill =} \StringTok{"\#18BC9C"}\NormalTok{) }\SpecialCharTok{+}
  \FunctionTok{coord\_flip}\NormalTok{() }\SpecialCharTok{+}
  \FunctionTok{labs}\NormalTok{(}\AttributeTok{title =} \StringTok{"Volume total por parcela"}\NormalTok{, }\AttributeTok{x =} \StringTok{"Parcela"}\NormalTok{, }\AttributeTok{y =} \FunctionTok{expression}\NormalTok{(Volume}\SpecialCharTok{\textasciitilde{}}\NormalTok{(m}\SpecialCharTok{\^{}}\DecValTok{3}\NormalTok{)))}
\end{Highlighting}
\end{Shaded}

\pandocbounded{\includegraphics[keepaspectratio]{04-dias/d2-tarde_files/figure-pdf/unnamed-chunk-9-1.pdf}}

\begin{center}\rule{0.5\linewidth}{0.5pt}\end{center}

\section{Customização de cores}\label{customizauxe7uxe3o-de-cores}

\begin{Shaded}
\begin{Highlighting}[]
\CommentTok{\# Cores pré{-}definidas (ColorBrewer)}
\NormalTok{inv\_florestal }\SpecialCharTok{\%\textgreater{}\%}
  \FunctionTok{ggplot}\NormalTok{(}\FunctionTok{aes}\NormalTok{(}\AttributeTok{x =}\NormalTok{ dap, }\AttributeTok{y =}\NormalTok{ altura, }\AttributeTok{color =}\NormalTok{ especie)) }\SpecialCharTok{+}
  \FunctionTok{geom\_point}\NormalTok{(}\AttributeTok{alpha =} \FloatTok{0.5}\NormalTok{) }\SpecialCharTok{+}
  \FunctionTok{scale\_color\_brewer}\NormalTok{(}\AttributeTok{palette =} \StringTok{"Set1"}\NormalTok{) }\SpecialCharTok{+}
  \FunctionTok{labs}\NormalTok{(}\AttributeTok{color =} \StringTok{"Espécie"}\NormalTok{)}
\end{Highlighting}
\end{Shaded}

\pandocbounded{\includegraphics[keepaspectratio]{04-dias/d2-tarde_files/figure-pdf/unnamed-chunk-10-1.pdf}}

\begin{Shaded}
\begin{Highlighting}[]
\CommentTok{\# Escala viridis (acessível para daltônicos)}
\NormalTok{inv\_florestal }\SpecialCharTok{\%\textgreater{}\%}
  \FunctionTok{ggplot}\NormalTok{(}\FunctionTok{aes}\NormalTok{(}\AttributeTok{x =}\NormalTok{ dap, }\AttributeTok{y =}\NormalTok{ altura, }\AttributeTok{color =}\NormalTok{ volume)) }\SpecialCharTok{+}
  \FunctionTok{geom\_point}\NormalTok{(}\AttributeTok{alpha =} \FloatTok{0.7}\NormalTok{) }\SpecialCharTok{+}
  \FunctionTok{scale\_color\_viridis\_c}\NormalTok{(}\AttributeTok{option =} \StringTok{"plasma"}\NormalTok{) }\SpecialCharTok{+}
  \FunctionTok{labs}\NormalTok{(}\AttributeTok{color =} \FunctionTok{expression}\NormalTok{(Volume}\SpecialCharTok{\textasciitilde{}}\NormalTok{(m}\SpecialCharTok{\^{}}\DecValTok{3}\NormalTok{)))}
\end{Highlighting}
\end{Shaded}

\pandocbounded{\includegraphics[keepaspectratio]{04-dias/d2-tarde_files/figure-pdf/unnamed-chunk-10-2.pdf}}

\begin{center}\rule{0.5\linewidth}{0.5pt}\end{center}

\section{Facetas: subgráficos por
grupo}\label{facetas-subgruxe1ficos-por-grupo}

\begin{Shaded}
\begin{Highlighting}[]
\CommentTok{\# Um gráfico por espécie}
\NormalTok{inv\_florestal }\SpecialCharTok{\%\textgreater{}\%}
  \FunctionTok{ggplot}\NormalTok{(}\FunctionTok{aes}\NormalTok{(}\AttributeTok{x =}\NormalTok{ dap, }\AttributeTok{y =}\NormalTok{ altura)) }\SpecialCharTok{+}
  \FunctionTok{geom\_point}\NormalTok{(}\AttributeTok{alpha =} \FloatTok{0.5}\NormalTok{, }\AttributeTok{color =} \StringTok{"\#2C3E50"}\NormalTok{) }\SpecialCharTok{+}
  \FunctionTok{geom\_smooth}\NormalTok{(}\AttributeTok{method =} \StringTok{"lm"}\NormalTok{, }\AttributeTok{se =} \ConstantTok{FALSE}\NormalTok{, }\AttributeTok{color =} \StringTok{"\#E74C3C"}\NormalTok{) }\SpecialCharTok{+}
  \FunctionTok{facet\_wrap}\NormalTok{(}\SpecialCharTok{\textasciitilde{}}\NormalTok{especie) }\SpecialCharTok{+}
  \FunctionTok{labs}\NormalTok{(}\AttributeTok{title =} \StringTok{"Relação hipsométrica por espécie"}\NormalTok{, }\AttributeTok{x =} \StringTok{"DAP (cm)"}\NormalTok{, }\AttributeTok{y =} \StringTok{"Altura (m)"}\NormalTok{)}
\end{Highlighting}
\end{Shaded}

\pandocbounded{\includegraphics[keepaspectratio]{04-dias/d2-tarde_files/figure-pdf/unnamed-chunk-11-1.pdf}}

\begin{Shaded}
\begin{Highlighting}[]
\CommentTok{\# Grade por espécie × parcela (selecionando 6 parcelas)}
\NormalTok{inv\_florestal }\SpecialCharTok{\%\textgreater{}\%}
  \FunctionTok{filter}\NormalTok{(parcela }\SpecialCharTok{\%in\%} \FunctionTok{c}\NormalTok{(}\DecValTok{1}\NormalTok{, }\DecValTok{2}\NormalTok{, }\DecValTok{3}\NormalTok{, }\DecValTok{7}\NormalTok{, }\DecValTok{12}\NormalTok{, }\DecValTok{15}\NormalTok{)) }\SpecialCharTok{\%\textgreater{}\%}
  \FunctionTok{ggplot}\NormalTok{(}\FunctionTok{aes}\NormalTok{(}\AttributeTok{x =}\NormalTok{ dap, }\AttributeTok{fill =}\NormalTok{ especie)) }\SpecialCharTok{+}
  \FunctionTok{geom\_histogram}\NormalTok{(}\AttributeTok{bins =} \DecValTok{20}\NormalTok{, }\AttributeTok{alpha =} \FloatTok{0.7}\NormalTok{, }\AttributeTok{position =} \StringTok{"identity"}\NormalTok{) }\SpecialCharTok{+}
  \FunctionTok{facet\_grid}\NormalTok{(parcela }\SpecialCharTok{\textasciitilde{}}\NormalTok{ especie) }\SpecialCharTok{+}
  \FunctionTok{labs}\NormalTok{(}\AttributeTok{title =} \StringTok{"Distribuição diamétrica — parcelas selecionadas"}\NormalTok{)}
\end{Highlighting}
\end{Shaded}

\pandocbounded{\includegraphics[keepaspectratio]{04-dias/d2-tarde_files/figure-pdf/unnamed-chunk-12-1.pdf}}

\begin{center}\rule{0.5\linewidth}{0.5pt}\end{center}

\section{Temas e aparência final}\label{temas-e-aparuxeancia-final}

\begin{Shaded}
\begin{Highlighting}[]
\CommentTok{\# Comparação de temas}
\NormalTok{p }\OtherTok{\textless{}{-}}\NormalTok{ inv\_florestal }\SpecialCharTok{\%\textgreater{}\%}
  \FunctionTok{ggplot}\NormalTok{(}\FunctionTok{aes}\NormalTok{(}\AttributeTok{x =}\NormalTok{ dap, }\AttributeTok{y =}\NormalTok{ altura, }\AttributeTok{color =}\NormalTok{ especie)) }\SpecialCharTok{+}
  \FunctionTok{geom\_point}\NormalTok{(}\AttributeTok{alpha =} \FloatTok{0.5}\NormalTok{)}

\CommentTok{\# Cinco temas diferentes}
\NormalTok{p }\SpecialCharTok{+} \FunctionTok{theme\_minimal}\NormalTok{() }\SpecialCharTok{+} \FunctionTok{labs}\NormalTok{(}\AttributeTok{title =} \StringTok{"theme\_minimal()"}\NormalTok{)}
\end{Highlighting}
\end{Shaded}

\pandocbounded{\includegraphics[keepaspectratio]{04-dias/d2-tarde_files/figure-pdf/unnamed-chunk-13-1.pdf}}

\begin{Shaded}
\begin{Highlighting}[]
\NormalTok{p }\SpecialCharTok{+} \FunctionTok{theme\_bw}\NormalTok{() }\SpecialCharTok{+} \FunctionTok{labs}\NormalTok{(}\AttributeTok{title =} \StringTok{"theme\_bw()"}\NormalTok{)}
\end{Highlighting}
\end{Shaded}

\pandocbounded{\includegraphics[keepaspectratio]{04-dias/d2-tarde_files/figure-pdf/unnamed-chunk-13-2.pdf}}

\begin{Shaded}
\begin{Highlighting}[]
\NormalTok{p }\SpecialCharTok{+} \FunctionTok{theme\_light}\NormalTok{() }\SpecialCharTok{+} \FunctionTok{labs}\NormalTok{(}\AttributeTok{title =} \StringTok{"theme\_light()"}\NormalTok{)}
\end{Highlighting}
\end{Shaded}

\pandocbounded{\includegraphics[keepaspectratio]{04-dias/d2-tarde_files/figure-pdf/unnamed-chunk-13-3.pdf}}

\begin{Shaded}
\begin{Highlighting}[]
\NormalTok{p }\SpecialCharTok{+} \FunctionTok{theme\_classic}\NormalTok{() }\SpecialCharTok{+} \FunctionTok{labs}\NormalTok{(}\AttributeTok{title =} \StringTok{"theme\_classic()"}\NormalTok{)}
\end{Highlighting}
\end{Shaded}

\pandocbounded{\includegraphics[keepaspectratio]{04-dias/d2-tarde_files/figure-pdf/unnamed-chunk-13-4.pdf}}

\begin{Shaded}
\begin{Highlighting}[]
\NormalTok{p }\SpecialCharTok{+} \FunctionTok{theme\_dark}\NormalTok{() }\SpecialCharTok{+} \FunctionTok{labs}\NormalTok{(}\AttributeTok{title =} \StringTok{"theme\_dark()"}\NormalTok{)}
\end{Highlighting}
\end{Shaded}

\pandocbounded{\includegraphics[keepaspectratio]{04-dias/d2-tarde_files/figure-pdf/unnamed-chunk-13-5.pdf}}

\begin{Shaded}
\begin{Highlighting}[]
\CommentTok{\# Gráfico pronto para relatório}
\NormalTok{inv\_florestal }\SpecialCharTok{\%\textgreater{}\%}
  \FunctionTok{ggplot}\NormalTok{(}\FunctionTok{aes}\NormalTok{(}\AttributeTok{x =}\NormalTok{ dap, }\AttributeTok{y =}\NormalTok{ altura, }\AttributeTok{color =}\NormalTok{ especie)) }\SpecialCharTok{+}
  \FunctionTok{geom\_point}\NormalTok{(}\AttributeTok{alpha =} \FloatTok{0.5}\NormalTok{, }\AttributeTok{size =} \DecValTok{2}\NormalTok{) }\SpecialCharTok{+}
  \FunctionTok{geom\_smooth}\NormalTok{(}\AttributeTok{method =} \StringTok{"lm"}\NormalTok{, }\AttributeTok{se =} \ConstantTok{TRUE}\NormalTok{, }\AttributeTok{alpha =} \FloatTok{0.2}\NormalTok{) }\SpecialCharTok{+}
  \FunctionTok{scale\_color\_brewer}\NormalTok{(}\AttributeTok{palette =} \StringTok{"Set1"}\NormalTok{) }\SpecialCharTok{+}
  \FunctionTok{labs}\NormalTok{(}
    \AttributeTok{title =} \StringTok{"Relação hipsométrica — Inventário Florestal"}\NormalTok{,}
    \AttributeTok{subtitle =} \StringTok{"Curvas ajustadas por espécie com intervalo de confiança de 95\%"}\NormalTok{,}
    \AttributeTok{x =} \StringTok{"DAP (cm)"}\NormalTok{,}
    \AttributeTok{y =} \StringTok{"Altura total (m)"}\NormalTok{,}
    \AttributeTok{color =} \StringTok{"Espécie"}\NormalTok{,}
    \AttributeTok{caption =} \StringTok{"Fonte: Dados do inventário florestal"}
\NormalTok{  ) }\SpecialCharTok{+}
  \FunctionTok{theme\_minimal}\NormalTok{(}\AttributeTok{base\_size =} \DecValTok{13}\NormalTok{) }\SpecialCharTok{+}
  \FunctionTok{theme}\NormalTok{(}
    \AttributeTok{plot.title =} \FunctionTok{element\_text}\NormalTok{(}\AttributeTok{face =} \StringTok{"bold"}\NormalTok{),}
    \AttributeTok{legend.position =} \StringTok{"bottom"}
\NormalTok{  )}
\end{Highlighting}
\end{Shaded}

\pandocbounded{\includegraphics[keepaspectratio]{04-dias/d2-tarde_files/figure-pdf/unnamed-chunk-14-1.pdf}}

\begin{center}\rule{0.5\linewidth}{0.5pt}\end{center}

\section{Exportando gráficos}\label{exportando-gruxe1ficos}

\begin{Shaded}
\begin{Highlighting}[]
\CommentTok{\# Salvar o último gráfico exibido}
\CommentTok{\# ggsave("relacao\_hipsometrica.png", width = 8, height = 5, dpi = 300)}

\CommentTok{\# Ou guardar o gráfico em um objeto e salvar}
\NormalTok{grafico\_final }\OtherTok{\textless{}{-}}\NormalTok{ inv\_florestal }\SpecialCharTok{\%\textgreater{}\%}
  \FunctionTok{ggplot}\NormalTok{(}\FunctionTok{aes}\NormalTok{(}\AttributeTok{x =}\NormalTok{ dap, }\AttributeTok{y =}\NormalTok{ altura, }\AttributeTok{color =}\NormalTok{ especie)) }\SpecialCharTok{+}
  \FunctionTok{geom\_point}\NormalTok{(}\AttributeTok{alpha =} \FloatTok{0.5}\NormalTok{) }\SpecialCharTok{+}
  \FunctionTok{geom\_smooth}\NormalTok{(}\AttributeTok{method =} \StringTok{"lm"}\NormalTok{, }\AttributeTok{se =} \ConstantTok{FALSE}\NormalTok{) }\SpecialCharTok{+}
  \FunctionTok{scale\_color\_brewer}\NormalTok{(}\AttributeTok{palette =} \StringTok{"Set1"}\NormalTok{) }\SpecialCharTok{+}
  \FunctionTok{labs}\NormalTok{(}\AttributeTok{title =} \StringTok{"Relação hipsométrica"}\NormalTok{, }\AttributeTok{x =} \StringTok{"DAP (cm)"}\NormalTok{, }\AttributeTok{y =} \StringTok{"Altura (m)"}\NormalTok{) }\SpecialCharTok{+}
  \FunctionTok{theme\_minimal}\NormalTok{()}

\CommentTok{\# ggsave("hipsometrica.png", grafico\_final, width = 8, height = 5, dpi = 300)}
\end{Highlighting}
\end{Shaded}

\begin{tcolorbox}[enhanced jigsaw, arc=.35mm, bottomrule=.15mm, bottomtitle=1mm, breakable, colback=white, colbacktitle=quarto-callout-tip-color!10!white, colframe=quarto-callout-tip-color-frame, coltitle=black, left=2mm, leftrule=.75mm, opacityback=0, opacitybacktitle=0.6, rightrule=.15mm, title=\textcolor{quarto-callout-tip-color}{\faLightbulb}\hspace{0.5em}{Formatos e resolução}, titlerule=0mm, toprule=.15mm, toptitle=1mm]

\begin{itemize}
\tightlist
\item
  \textbf{PNG} --- melhor para web e apresentações
\item
  \textbf{PDF} --- melhor para publicações e impressão
\item
  \textbf{SVG} --- gráficos vetoriais redimensionáveis
\item
  \textbf{DPI} --- 150 para web, 300 para impressão, 600 para publicação
\end{itemize}

\end{tcolorbox}

\subsection{Dashboard: combinando gráficos com
patchwork}\label{dashboard-combinando-gruxe1ficos-com-patchwork}

\begin{Shaded}
\begin{Highlighting}[]
\CommentTok{\# Quatro gráficos em um painel}
\NormalTok{p1 }\OtherTok{\textless{}{-}}\NormalTok{ inv\_florestal }\SpecialCharTok{\%\textgreater{}\%}
  \FunctionTok{ggplot}\NormalTok{(}\FunctionTok{aes}\NormalTok{(}\AttributeTok{x =}\NormalTok{ dap)) }\SpecialCharTok{+} \FunctionTok{geom\_histogram}\NormalTok{(}\AttributeTok{bins =} \DecValTok{30}\NormalTok{, }\AttributeTok{fill =} \StringTok{"\#18BC9C"}\NormalTok{) }\SpecialCharTok{+} \FunctionTok{labs}\NormalTok{(}\AttributeTok{title =} \StringTok{"Distribuição diamétrica"}\NormalTok{)}

\NormalTok{p2 }\OtherTok{\textless{}{-}}\NormalTok{ inv\_florestal }\SpecialCharTok{\%\textgreater{}\%}
  \FunctionTok{ggplot}\NormalTok{(}\FunctionTok{aes}\NormalTok{(}\AttributeTok{x =}\NormalTok{ dap, }\AttributeTok{y =}\NormalTok{ altura, }\AttributeTok{color =}\NormalTok{ especie)) }\SpecialCharTok{+} \FunctionTok{geom\_point}\NormalTok{(}\AttributeTok{alpha =} \FloatTok{0.5}\NormalTok{) }\SpecialCharTok{+} \FunctionTok{labs}\NormalTok{(}\AttributeTok{title =} \StringTok{"Relação hipsométrica"}\NormalTok{)}

\NormalTok{p3 }\OtherTok{\textless{}{-}}\NormalTok{ inv\_florestal }\SpecialCharTok{\%\textgreater{}\%}
  \FunctionTok{ggplot}\NormalTok{(}\FunctionTok{aes}\NormalTok{(}\AttributeTok{x =}\NormalTok{ especie, }\AttributeTok{y =}\NormalTok{ volume, }\AttributeTok{fill =}\NormalTok{ especie)) }\SpecialCharTok{+} \FunctionTok{geom\_boxplot}\NormalTok{() }\SpecialCharTok{+} \FunctionTok{labs}\NormalTok{(}\AttributeTok{title =} \StringTok{"Volume por espécie"}\NormalTok{) }\SpecialCharTok{+} \FunctionTok{theme}\NormalTok{(}\AttributeTok{legend.position =} \StringTok{"none"}\NormalTok{)}

\NormalTok{p4 }\OtherTok{\textless{}{-}}\NormalTok{ inv\_florestal }\SpecialCharTok{\%\textgreater{}\%}
  \FunctionTok{group\_by}\NormalTok{(parcela) }\SpecialCharTok{\%\textgreater{}\%} \FunctionTok{summarise}\NormalTok{(}\AttributeTok{v =} \FunctionTok{sum}\NormalTok{(volume)) }\SpecialCharTok{\%\textgreater{}\%}
  \FunctionTok{ggplot}\NormalTok{(}\FunctionTok{aes}\NormalTok{(}\AttributeTok{x =} \FunctionTok{factor}\NormalTok{(parcela), }\AttributeTok{y =}\NormalTok{ v)) }\SpecialCharTok{+} \FunctionTok{geom\_col}\NormalTok{(}\AttributeTok{fill =} \StringTok{"\#2C3E50"}\NormalTok{) }\SpecialCharTok{+} \FunctionTok{labs}\NormalTok{(}\AttributeTok{title =} \StringTok{"Volume por parcela"}\NormalTok{, }\AttributeTok{x =} \StringTok{"Parcela"}\NormalTok{)}

\NormalTok{(p1 }\SpecialCharTok{+}\NormalTok{ p2) }\SpecialCharTok{/}\NormalTok{ (p3 }\SpecialCharTok{+}\NormalTok{ p4) }\SpecialCharTok{+}
  \FunctionTok{plot\_annotation}\NormalTok{(}\AttributeTok{title =} \StringTok{"Dashboard — Inventário Florestal"}\NormalTok{, }\AttributeTok{theme =} \FunctionTok{theme}\NormalTok{(}\AttributeTok{plot.title =} \FunctionTok{element\_text}\NormalTok{(}\AttributeTok{face =} \StringTok{"bold"}\NormalTok{, }\AttributeTok{size =} \DecValTok{16}\NormalTok{)))}
\end{Highlighting}
\end{Shaded}

\pandocbounded{\includegraphics[keepaspectratio]{04-dias/d2-tarde_files/figure-pdf/unnamed-chunk-16-1.pdf}}

\begin{center}\rule{0.5\linewidth}{0.5pt}\end{center}

\section{Exercícios}\label{exercuxedcios-10}

\subsection{Exercício 1}\label{exercuxedcio-1-6}

Crie um histograma da distribuição diamétrica (\texttt{dap}) com 30
bins, colorido em verde (\texttt{\#18BC9C}). Adicione título e rótulos
aos eixos.

\textbf{Resultado esperado:}

\emph{Histograma com pico em torno de 15-20 cm, distribuição assimétrica
à direita.}

\subsection{Exercício 2}\label{exercuxedcio-2-6}

Crie um boxplot de \texttt{volume} por \texttt{especie}. Qual espécie
tem maior variabilidade?

\textbf{Resultado esperado:}

\emph{Eucalipto com caixa mais larga e mediana mais alta que Pinus.}

\subsection{Exercício 3}\label{exercuxedcio-3-5}

Plote \texttt{dap} × \texttt{volume} com pontos (\texttt{alpha\ =\ 0.3})
e uma curva de tendência suave (\texttt{geom\_smooth}). Colora os pontos
por espécie.

\textbf{Resultado esperado:}

\emph{Nuvem de pontos com duas curvas não-lineares --- volume cresce
exponencialmente com DAP.}

\subsection{Exercício 4}\label{exercuxedcio-4-4}

Usando \texttt{facet\_wrap(\textasciitilde{}especie)}, crie um gráfico
de dispersão \texttt{dap} × \texttt{altura} com um painel para cada
espécie.

\textbf{Resultado esperado:}

\emph{Dois painéis lado a lado, cada um com a relação hipsométrica de
uma espécie.}

\subsection{Exercício 5}\label{exercuxedcio-5-2}

Crie um gráfico de barras do volume total por parcela (top 10), ordenado
do maior para o menor. Use \texttt{coord\_flip()} para facilitar a
leitura.

\textbf{Resultado esperado:}

\emph{Gráfico horizontal com 10 barras verdes, parcela 7 com maior
volume.}

\subsection{Exercício 6}\label{exercuxedcio-6-1}

Exporte seu melhor gráfico como PNG com 8 polegadas de largura, 5 de
altura e 300 DPI.

\chapter{Dia 3 --- Manhã}\label{dia-3-manhuxe3}

\textbf{Dia 3 --- Manhã \textbar{} 4 horas}

\section*{Objetivos}\label{objetivos-4}
\addcontentsline{toc}{section}{Objetivos}

\markright{Objetivos}

Ao final deste capítulo, você será capaz de:

\begin{itemize}
\tightlist
\item
  Manipular texto com \texttt{stringr} (buscar, extrair, substituir)
\item
  Trabalhar com fatores usando \texttt{forcats} (reordenar, agrupar)
\item
  Lidar com datas usando \texttt{lubridate} (parse, extrair, calcular
  intervalos)
\end{itemize}

\begin{center}\rule{0.5\linewidth}{0.5pt}\end{center}

\section{\texorpdfstring{Manipulação de texto com
\texttt{stringr}}{Manipulação de texto com stringr}}\label{manipulauxe7uxe3o-de-texto-com-stringr}

Dados reais vêm sujos: espaços extras, maiúsculas inconsistentes,
padrões que precisam ser extraídos. O pacote \texttt{stringr} resolve
isso.

\subsection{Detectar e contar
padrões}\label{detectar-e-contar-padruxf5es}

\begin{Shaded}
\begin{Highlighting}[]
\FunctionTok{library}\NormalTok{(stringr)}

\CommentTok{\# Quantos filmes têm "The" no título?}
\NormalTok{filmes }\SpecialCharTok{\%\textgreater{}\%}
  \FunctionTok{mutate}\NormalTok{(}\AttributeTok{tem\_the =} \FunctionTok{str\_detect}\NormalTok{(titulo, }\StringTok{"\^{}The "}\NormalTok{)) }\SpecialCharTok{\%\textgreater{}\%}
  \FunctionTok{count}\NormalTok{(tem\_the)}
\end{Highlighting}
\end{Shaded}

\begin{verbatim}
# A tibble: 2 x 2
  tem_the     n
  <lgl>   <int>
1 FALSE    3082
2 TRUE      793
\end{verbatim}

\begin{Shaded}
\begin{Highlighting}[]
\CommentTok{\# Quantas espécies contêm "Eucalyptus"?}
\NormalTok{inv\_florestal }\SpecialCharTok{\%\textgreater{}\%}
  \FunctionTok{mutate}\NormalTok{(}\AttributeTok{euca =} \FunctionTok{str\_detect}\NormalTok{(especie, }\StringTok{"Eucalyptus"}\NormalTok{)) }\SpecialCharTok{\%\textgreater{}\%}
  \FunctionTok{count}\NormalTok{(euca)}
\end{Highlighting}
\end{Shaded}

\begin{verbatim}
# A tibble: 2 x 2
  euca      n
  <lgl> <int>
1 FALSE   350
2 TRUE    781
\end{verbatim}

\subsection{Extrair e substituir}\label{extrair-e-substituir}

\begin{Shaded}
\begin{Highlighting}[]
\CommentTok{\# Extrair o ano de uma data em formato texto}
\FunctionTok{str\_extract}\NormalTok{(}\StringTok{"15/03/2023"}\NormalTok{, }\StringTok{"}\SpecialCharTok{\textbackslash{}\textbackslash{}}\StringTok{d\{4\}$"}\NormalTok{)}
\end{Highlighting}
\end{Shaded}

\begin{verbatim}
[1] "2023"
\end{verbatim}

\begin{Shaded}
\begin{Highlighting}[]
\CommentTok{\# Substituir abreviações}
\FunctionTok{str\_replace}\NormalTok{(}\StringTok{"Eucalyptus grandis"}\NormalTok{, }\StringTok{"Eucalyptus"}\NormalTok{, }\StringTok{"E."}\NormalTok{)}
\end{Highlighting}
\end{Shaded}

\begin{verbatim}
[1] "E. grandis"
\end{verbatim}

\begin{Shaded}
\begin{Highlighting}[]
\CommentTok{\# Padronizar nomes de espécies}
\NormalTok{inv\_florestal }\SpecialCharTok{\%\textgreater{}\%}
  \FunctionTok{mutate}\NormalTok{(}\AttributeTok{especie\_curto =} \FunctionTok{str\_replace}\NormalTok{(especie, }\StringTok{"Eucalyptus"}\NormalTok{, }\StringTok{"E."}\NormalTok{)) }\SpecialCharTok{\%\textgreater{}\%}
  \FunctionTok{distinct}\NormalTok{(especie\_curto)}
\end{Highlighting}
\end{Shaded}

\begin{verbatim}
# A tibble: 3 x 1
  especie_curto
  <chr>        
1 E. Grandis   
2 Pinus Taeda  
3 E. Urophylla 
\end{verbatim}

\subsection{Limpeza de texto}\label{limpeza-de-texto}

\begin{Shaded}
\begin{Highlighting}[]
\CommentTok{\# Texto sujo (simulando dado real de campo)}
\NormalTok{nomes\_sujos }\OtherTok{\textless{}{-}} \FunctionTok{c}\NormalTok{(}\StringTok{"  EUCALYPTUS GRANDIS  "}\NormalTok{, }\StringTok{"Pinus taeda"}\NormalTok{, }\StringTok{"Eucalyptus  urophylla"}\NormalTok{)}

\CommentTok{\# Faxina completa}
\NormalTok{nomes\_sujos }\SpecialCharTok{\%\textgreater{}\%}
  \FunctionTok{str\_squish}\NormalTok{() }\SpecialCharTok{\%\textgreater{}\%}        \CommentTok{\# remove espaços extras}
  \FunctionTok{str\_to\_lower}\NormalTok{() }\SpecialCharTok{\%\textgreater{}\%}      \CommentTok{\# tudo minúsculo}
  \FunctionTok{str\_to\_title}\NormalTok{()          }\CommentTok{\# primeira letra maiúscula}
\end{Highlighting}
\end{Shaded}

\begin{verbatim}
[1] "Eucalyptus Grandis"   "Pinus Taeda"          "Eucalyptus Urophylla"
\end{verbatim}

\begin{Shaded}
\begin{Highlighting}[]
\CommentTok{\# Aplicando limpeza ao dataset}
\NormalTok{inv\_florestal }\SpecialCharTok{\%\textgreater{}\%}
  \FunctionTok{mutate}\NormalTok{(}\AttributeTok{especie\_limpa =} \FunctionTok{str\_squish}\NormalTok{(especie) }\SpecialCharTok{\%\textgreater{}\%} \FunctionTok{str\_to\_title}\NormalTok{()) }\SpecialCharTok{\%\textgreater{}\%}
  \FunctionTok{count}\NormalTok{(especie\_limpa)}
\end{Highlighting}
\end{Shaded}

\begin{verbatim}
# A tibble: 3 x 2
  especie_limpa            n
  <chr>                <int>
1 Eucalyptus Grandis     682
2 Eucalyptus Urophylla    99
3 Pinus Taeda            350
\end{verbatim}

\begin{center}\rule{0.5\linewidth}{0.5pt}\end{center}

\section{\texorpdfstring{Fatores com
\texttt{forcats}}{Fatores com forcats}}\label{fatores-com-forcats}

Fatores são variáveis categóricas. O pacote \texttt{forcats} ajuda a
reordená-los e agrupá-los --- essencial para gráficos e modelos.

\subsection{Reordenar por outra
variável}\label{reordenar-por-outra-variuxe1vel}

\begin{Shaded}
\begin{Highlighting}[]
\CommentTok{\# Ordenar espécies por DAP médio (útil para gráficos)}
\NormalTok{inv\_florestal }\SpecialCharTok{\%\textgreater{}\%}
  \FunctionTok{mutate}\NormalTok{(}\AttributeTok{especie =} \FunctionTok{fct\_reorder}\NormalTok{(especie, dap, }\AttributeTok{.fun =}\NormalTok{ median)) }\SpecialCharTok{\%\textgreater{}\%}
  \FunctionTok{ggplot}\NormalTok{(}\FunctionTok{aes}\NormalTok{(}\AttributeTok{x =}\NormalTok{ especie, }\AttributeTok{y =}\NormalTok{ dap, }\AttributeTok{fill =}\NormalTok{ especie)) }\SpecialCharTok{+}
  \FunctionTok{geom\_boxplot}\NormalTok{() }\SpecialCharTok{+}
  \FunctionTok{labs}\NormalTok{(}\AttributeTok{x =} \ConstantTok{NULL}\NormalTok{, }\AttributeTok{y =} \StringTok{"DAP (cm)"}\NormalTok{, }\AttributeTok{title =} \StringTok{"DAP por espécie (ordenado pela mediana)"}\NormalTok{) }\SpecialCharTok{+}
  \FunctionTok{theme}\NormalTok{(}\AttributeTok{legend.position =} \StringTok{"none"}\NormalTok{)}
\end{Highlighting}
\end{Shaded}

\pandocbounded{\includegraphics[keepaspectratio]{04-dias/d3-manha_files/figure-pdf/unnamed-chunk-7-1.pdf}}

\subsection{Agrupar categorias raras}\label{agrupar-categorias-raras}

\begin{Shaded}
\begin{Highlighting}[]
\CommentTok{\# Agrupar gêneros com poucos filmes como "Outros"}
\NormalTok{filmes }\SpecialCharTok{\%\textgreater{}\%}
  \FunctionTok{mutate}\NormalTok{(}\AttributeTok{genero\_agrupado =} \FunctionTok{fct\_lump}\NormalTok{(genero, }\AttributeTok{n =} \DecValTok{5}\NormalTok{)) }\SpecialCharTok{\%\textgreater{}\%}
  \FunctionTok{count}\NormalTok{(genero\_agrupado, }\AttributeTok{sort =} \ConstantTok{TRUE}\NormalTok{)}
\end{Highlighting}
\end{Shaded}

\begin{verbatim}
# A tibble: 6 x 2
  genero_agrupado       n
  <fct>             <int>
1 Other              1521
2 Drama               719
3 Ação                566
4 Comédia             452
5 Terror              311
6 Ficção científica   306
\end{verbatim}

\begin{Shaded}
\begin{Highlighting}[]
\CommentTok{\# Agrupar espécies raras no inventário}
\NormalTok{inv\_florestal }\SpecialCharTok{\%\textgreater{}\%}
  \FunctionTok{mutate}\NormalTok{(}\AttributeTok{especie\_agrupada =} \FunctionTok{fct\_lump}\NormalTok{(especie, }\AttributeTok{n =} \DecValTok{3}\NormalTok{)) }\SpecialCharTok{\%\textgreater{}\%}
  \FunctionTok{count}\NormalTok{(especie\_agrupada, }\AttributeTok{sort =} \ConstantTok{TRUE}\NormalTok{)}
\end{Highlighting}
\end{Shaded}

\begin{verbatim}
# A tibble: 3 x 2
  especie_agrupada         n
  <fct>                <int>
1 Eucalyptus Grandis     682
2 Pinus Taeda            350
3 Eucalyptus Urophylla    99
\end{verbatim}

\subsection{Ordenar por frequência}\label{ordenar-por-frequuxeancia}

\begin{Shaded}
\begin{Highlighting}[]
\CommentTok{\# Gênero mais comum primeiro}
\NormalTok{filmes }\SpecialCharTok{\%\textgreater{}\%}
  \FunctionTok{mutate}\NormalTok{(}\AttributeTok{genero =} \FunctionTok{fct\_infreq}\NormalTok{(genero)) }\SpecialCharTok{\%\textgreater{}\%}
  \FunctionTok{ggplot}\NormalTok{(}\FunctionTok{aes}\NormalTok{(}\AttributeTok{y =}\NormalTok{ genero)) }\SpecialCharTok{+}
  \FunctionTok{geom\_bar}\NormalTok{(}\AttributeTok{fill =} \StringTok{"\#18BC9C"}\NormalTok{) }\SpecialCharTok{+}
  \FunctionTok{labs}\NormalTok{(}\AttributeTok{title =} \StringTok{"Frequência de gêneros"}\NormalTok{, }\AttributeTok{x =} \StringTok{"Nº de filmes"}\NormalTok{, }\AttributeTok{y =} \ConstantTok{NULL}\NormalTok{)}
\end{Highlighting}
\end{Shaded}

\pandocbounded{\includegraphics[keepaspectratio]{04-dias/d3-manha_files/figure-pdf/unnamed-chunk-10-1.pdf}}

\begin{center}\rule{0.5\linewidth}{0.5pt}\end{center}

\section{\texorpdfstring{Datas com
\texttt{lubridate}}{Datas com lubridate}}\label{datas-com-lubridate}

Lidar com datas no R sempre foi doloroso --- até o \texttt{lubridate}.
Ele transforma texto em data com funções intuitivas.

\subsection{Parse de datas}\label{parse-de-datas}

\begin{Shaded}
\begin{Highlighting}[]
\FunctionTok{library}\NormalTok{(lubridate)}

\CommentTok{\# Formatos comuns no Brasil}
\FunctionTok{dmy}\NormalTok{(}\StringTok{"15/03/2023"}\NormalTok{)       }\CommentTok{\# dia/mês/ano}
\end{Highlighting}
\end{Shaded}

\begin{verbatim}
[1] "2023-03-15"
\end{verbatim}

\begin{Shaded}
\begin{Highlighting}[]
\FunctionTok{dmy}\NormalTok{(}\StringTok{"15 de março de 2023"}\NormalTok{)}
\end{Highlighting}
\end{Shaded}

\begin{verbatim}
[1] "2023-03-15"
\end{verbatim}

\begin{Shaded}
\begin{Highlighting}[]
\FunctionTok{ymd}\NormalTok{(}\StringTok{"2023{-}03{-}15"}\NormalTok{)       }\CommentTok{\# ano{-}mês{-}dia (ISO)}
\end{Highlighting}
\end{Shaded}

\begin{verbatim}
[1] "2023-03-15"
\end{verbatim}

\begin{Shaded}
\begin{Highlighting}[]
\FunctionTok{mdy}\NormalTok{(}\StringTok{"March 15, 2023"}\NormalTok{)   }\CommentTok{\# mês/dia/ano (inglês)}
\end{Highlighting}
\end{Shaded}

\begin{verbatim}
[1] "2023-03-15"
\end{verbatim}

\begin{Shaded}
\begin{Highlighting}[]
\CommentTok{\# Com hora}
\FunctionTok{ymd\_hms}\NormalTok{(}\StringTok{"2023{-}03{-}15 14:30:00"}\NormalTok{)}
\end{Highlighting}
\end{Shaded}

\begin{verbatim}
[1] "2023-03-15 14:30:00 UTC"
\end{verbatim}

\begin{Shaded}
\begin{Highlighting}[]
\FunctionTok{dmy\_hm}\NormalTok{(}\StringTok{"15/03/2023 14:30"}\NormalTok{)}
\end{Highlighting}
\end{Shaded}

\begin{verbatim}
[1] "2023-03-15 14:30:00 UTC"
\end{verbatim}

\subsection{Extrair componentes}\label{extrair-componentes}

\begin{Shaded}
\begin{Highlighting}[]
\NormalTok{data\_medicao }\OtherTok{\textless{}{-}} \FunctionTok{dmy}\NormalTok{(}\StringTok{"15/03/2023"}\NormalTok{)}

\FunctionTok{year}\NormalTok{(data\_medicao)     }\CommentTok{\# ano}
\end{Highlighting}
\end{Shaded}

\begin{verbatim}
[1] 2023
\end{verbatim}

\begin{Shaded}
\begin{Highlighting}[]
\FunctionTok{month}\NormalTok{(data\_medicao)    }\CommentTok{\# mês (número)}
\end{Highlighting}
\end{Shaded}

\begin{verbatim}
[1] 3
\end{verbatim}

\begin{Shaded}
\begin{Highlighting}[]
\FunctionTok{month}\NormalTok{(data\_medicao, }\AttributeTok{label =} \ConstantTok{TRUE}\NormalTok{)  }\CommentTok{\# mês (nome)}
\end{Highlighting}
\end{Shaded}

\begin{verbatim}
[1] Mar
12 Levels: Jan < Feb < Mar < Apr < May < Jun < Jul < Aug < Sep < ... < Dec
\end{verbatim}

\begin{Shaded}
\begin{Highlighting}[]
\FunctionTok{wday}\NormalTok{(data\_medicao, }\AttributeTok{label =} \ConstantTok{TRUE}\NormalTok{)   }\CommentTok{\# dia da semana}
\end{Highlighting}
\end{Shaded}

\begin{verbatim}
[1] Wed
Levels: Sun < Mon < Tue < Wed < Thu < Fri < Sat
\end{verbatim}

\begin{Shaded}
\begin{Highlighting}[]
\FunctionTok{quarter}\NormalTok{(data\_medicao)  }\CommentTok{\# trimestre}
\end{Highlighting}
\end{Shaded}

\begin{verbatim}
[1] 1
\end{verbatim}

\begin{Shaded}
\begin{Highlighting}[]
\FunctionTok{semester}\NormalTok{(data\_medicao) }\CommentTok{\# semestre}
\end{Highlighting}
\end{Shaded}

\begin{verbatim}
[1] 1
\end{verbatim}

\subsection{Aritmética com datas}\label{aritmuxe9tica-com-datas}

\begin{Shaded}
\begin{Highlighting}[]
\CommentTok{\# Adicionar dias}
\NormalTok{data\_medicao }\SpecialCharTok{+} \FunctionTok{days}\NormalTok{(}\DecValTok{30}\NormalTok{)}
\end{Highlighting}
\end{Shaded}

\begin{verbatim}
[1] "2023-04-14"
\end{verbatim}

\begin{Shaded}
\begin{Highlighting}[]
\CommentTok{\# Adicionar meses}
\NormalTok{data\_medicao }\SpecialCharTok{+} \FunctionTok{months}\NormalTok{(}\DecValTok{6}\NormalTok{)}
\end{Highlighting}
\end{Shaded}

\begin{verbatim}
[1] "2023-09-15"
\end{verbatim}

\begin{Shaded}
\begin{Highlighting}[]
\CommentTok{\# Adicionar anos}
\NormalTok{data\_medicao }\SpecialCharTok{+} \FunctionTok{years}\NormalTok{(}\DecValTok{1}\NormalTok{)}
\end{Highlighting}
\end{Shaded}

\begin{verbatim}
[1] "2024-03-15"
\end{verbatim}

\begin{Shaded}
\begin{Highlighting}[]
\CommentTok{\# Calcular idade do povoamento}
\NormalTok{data\_plantio }\OtherTok{\textless{}{-}} \FunctionTok{dmy}\NormalTok{(}\StringTok{"15/05/2017"}\NormalTok{)}
\NormalTok{intervalo }\OtherTok{\textless{}{-}} \FunctionTok{interval}\NormalTok{(data\_plantio, data\_medicao)}
\NormalTok{intervalo }\SpecialCharTok{/} \FunctionTok{years}\NormalTok{(}\DecValTok{1}\NormalTok{)    }\CommentTok{\# idade em anos decimais}
\end{Highlighting}
\end{Shaded}

\begin{verbatim}
[1] 5.832877
\end{verbatim}

\subsection{Exemplo prático: dados
climáticos}\label{exemplo-pruxe1tico-dados-climuxe1ticos}

\begin{Shaded}
\begin{Highlighting}[]
\NormalTok{clima }\OtherTok{\textless{}{-}} \FunctionTok{read\_csv}\NormalTok{(}\StringTok{"../data/clima.csv"}\NormalTok{)}

\NormalTok{clima }\SpecialCharTok{\%\textgreater{}\%}
  \FunctionTok{mutate}\NormalTok{(}
    \AttributeTok{data =} \FunctionTok{dmy}\NormalTok{(data),}
    \AttributeTok{mes =} \FunctionTok{month}\NormalTok{(data, }\AttributeTok{label =} \ConstantTok{TRUE}\NormalTok{, }\AttributeTok{abbr =} \ConstantTok{FALSE}\NormalTok{),}
    \AttributeTok{ano =} \FunctionTok{year}\NormalTok{(data)}
\NormalTok{  ) }\SpecialCharTok{\%\textgreater{}\%}
  \FunctionTok{group\_by}\NormalTok{(mes) }\SpecialCharTok{\%\textgreater{}\%}
  \FunctionTok{summarise}\NormalTok{(}
    \AttributeTok{temp\_media =} \FunctionTok{mean}\NormalTok{(temperatura, }\AttributeTok{na.rm =} \ConstantTok{TRUE}\NormalTok{),}
    \AttributeTok{precip\_total =} \FunctionTok{sum}\NormalTok{(precipitacao, }\AttributeTok{na.rm =} \ConstantTok{TRUE}\NormalTok{)}
\NormalTok{  )}
\end{Highlighting}
\end{Shaded}

\begin{verbatim}
# A tibble: 12 x 3
   mes       temp_media precip_total
   <ord>          <dbl>        <dbl>
 1 January         21.7         681.
 2 February        23.2         860.
 3 March           24.7         886.
 4 April           25.1         942.
 5 May             23.7         977.
 6 June            21.5         662.
 7 July            18.5         556.
 8 August          16.0         359.
 9 September       15.5         346.
10 October         15.2         305.
11 November        16.3         390.
12 December        18.7         593.
\end{verbatim}

\begin{Shaded}
\begin{Highlighting}[]
\CommentTok{\# Série temporal de temperatura}
\NormalTok{clima }\SpecialCharTok{\%\textgreater{}\%}
  \FunctionTok{mutate}\NormalTok{(}\AttributeTok{data =} \FunctionTok{dmy}\NormalTok{(data)) }\SpecialCharTok{\%\textgreater{}\%}
  \FunctionTok{ggplot}\NormalTok{(}\FunctionTok{aes}\NormalTok{(}\AttributeTok{x =}\NormalTok{ data, }\AttributeTok{y =}\NormalTok{ temperatura)) }\SpecialCharTok{+}
  \FunctionTok{geom\_line}\NormalTok{(}\AttributeTok{color =} \StringTok{"\#E74C3C"}\NormalTok{) }\SpecialCharTok{+}
  \FunctionTok{geom\_smooth}\NormalTok{(}\AttributeTok{se =} \ConstantTok{FALSE}\NormalTok{, }\AttributeTok{color =} \StringTok{"\#2C3E50"}\NormalTok{) }\SpecialCharTok{+}
  \FunctionTok{labs}\NormalTok{(}\AttributeTok{title =} \StringTok{"Temperatura ao longo do tempo"}\NormalTok{, }\AttributeTok{x =} \ConstantTok{NULL}\NormalTok{, }\AttributeTok{y =} \StringTok{"Temperatura (°C)"}\NormalTok{)}
\end{Highlighting}
\end{Shaded}

\pandocbounded{\includegraphics[keepaspectratio]{04-dias/d3-manha_files/figure-pdf/unnamed-chunk-16-1.pdf}}

\begin{center}\rule{0.5\linewidth}{0.5pt}\end{center}

\section{Exercícios}\label{exercuxedcios-11}

\subsection{Exercício 1}\label{exercuxedcio-1-7}

Na coluna \texttt{especie} do \texttt{inv\_florestal}, quantos registros
contêm ``Eucalyptus'' no nome?

\textbf{Resultado esperado:}

\begin{verbatim}
[1] 664
\end{verbatim}

\subsection{Exercício 2}\label{exercuxedcio-2-7}

Padronize os nomes das espécies removendo espaços extras e convertendo
para maiúsculas (título).

\textbf{Resultado esperado:}

\begin{verbatim}
[1] "Eucalyptus Grandis" "Pinus Taeda"
\end{verbatim}

\subsection{Exercício 3}\label{exercuxedcio-3-6}

Agrupe as espécies com menos de 100 árvores como ``Outras''. Conte a
frequência.

\textbf{Resultado esperado:}

\begin{verbatim}
# A tibble: 3 × 2
  especie_agrupada     n
  <fct>            <int>
1 Eucalyptus Grandis  664
2 Pinus Taeda         282
3 Outras               46
\end{verbatim}

\subsection{Exercício 4}\label{exercuxedcio-4-5}

No dataset \texttt{clima.xlsx}, parse a coluna de data e calcule a
temperatura média por mês (ordenada de janeiro a dezembro).

\textbf{Resultado esperado:}

\begin{verbatim}
# A tibble: 12 × 2
   mes       temp_media
   <ord>          <dbl>
 1 janeiro         24.5
 2 fevereiro       24.2
 3 março           23.1
...
12 dezembro        24.8
\end{verbatim}

\subsection{Exercício 5}\label{exercuxedcio-5-3}

Usando os dados de \texttt{clima}, crie um gráfico de barras da
precipitação total por mês.

\textbf{Resultado esperado:}

\emph{Gráfico com 12 barras azuis, maior precipitação no verão
(dez-jan).}

\subsection{Exercício 6}\label{exercuxedcio-6-2}

Parse a data \texttt{"15/05/2017"} e calcule quantos anos se passaram
até hoje.

\textbf{Resultado esperado:}

\begin{verbatim}
[1] 9.04
\end{verbatim}

\emph{(O valor exato depende da data atual.)}

\chapter{Dia 3 --- Tarde}\label{dia-3-tarde}

\textbf{Dia 3 --- Tarde \textbar{} 4 horas}

\section*{Objetivos}\label{objetivos-5}
\addcontentsline{toc}{section}{Objetivos}

\markright{Objetivos}

\begin{itemize}
\tightlist
\item
  Entender o conceito de dados organizados (tidy data)
\item
  Transformar dados entre formato largo e longo com
  \texttt{pivot\_longer()} e \texttt{pivot\_wider()}
\item
  Importar dados de múltiplas fontes (múltiplas abas Excel, múltiplos
  CSVs)
\end{itemize}

\begin{center}\rule{0.5\linewidth}{0.5pt}\end{center}

\section{Tidy Data}\label{tidy-data}

Dados organizados (\emph{tidy data}) seguem três regras:

\begin{enumerate}
\def\labelenumi{\arabic{enumi}.}
\tightlist
\item
  \textbf{Cada coluna é uma variável}
\item
  \textbf{Cada linha é uma observação}
\item
  \textbf{Cada célula é um valor}
\end{enumerate}

Este conceito, proposto por Hadley Wickham, é a base para todo o
\texttt{tidyverse}. Quando seus dados estão organizados, todas as
funções funcionam perfeitamente juntas.

\subsection{Formato largo vs formato
longo}\label{formato-largo-vs-formato-longo}

\textbf{Formato largo} (comum em planilhas Excel, visualização humana):

{\def\LTcaptype{none} % do not increment counter
\begin{longtable}[]{@{}llll@{}}
\toprule\noalign{}
parcela & dap\_2021 & dap\_2022 & dap\_2023 \\
\midrule\noalign{}
\endhead
\bottomrule\noalign{}
\endlastfoot
1 & 15.2 & 17.8 & 20.1 \\
2 & 14.8 & 16.9 & 19.5 \\
\end{longtable}
}

\textbf{Formato longo} (ideal para análise e ggplot2):

{\def\LTcaptype{none} % do not increment counter
\begin{longtable}[]{@{}lll@{}}
\toprule\noalign{}
parcela & ano & dap \\
\midrule\noalign{}
\endhead
\bottomrule\noalign{}
\endlastfoot
1 & 2021 & 15.2 \\
1 & 2022 & 17.8 \\
1 & 2023 & 20.1 \\
2 & 2021 & 14.8 \\
2 & 2022 & 16.9 \\
2 & 2023 & 19.5 \\
\end{longtable}
}

\subsection{\texorpdfstring{\texttt{pivot\_longer()} --- largo →
longo}{pivot\_longer() --- largo → longo}}\label{pivot_longer-largo-longo}

\begin{Shaded}
\begin{Highlighting}[]
\CommentTok{\# Simulando inventário contínuo (formato largo)}
\NormalTok{inv\_largo }\OtherTok{\textless{}{-}} \FunctionTok{tibble}\NormalTok{(}
  \AttributeTok{parcela =} \DecValTok{1}\SpecialCharTok{:}\DecValTok{3}\NormalTok{,}
  \AttributeTok{dap\_2021 =} \FunctionTok{c}\NormalTok{(}\FloatTok{15.2}\NormalTok{, }\FloatTok{14.8}\NormalTok{, }\FloatTok{13.5}\NormalTok{),}
  \AttributeTok{dap\_2022 =} \FunctionTok{c}\NormalTok{(}\FloatTok{17.8}\NormalTok{, }\FloatTok{16.9}\NormalTok{, }\FloatTok{15.2}\NormalTok{),}
  \AttributeTok{dap\_2023 =} \FunctionTok{c}\NormalTok{(}\FloatTok{20.1}\NormalTok{, }\FloatTok{19.5}\NormalTok{, }\FloatTok{17.8}\NormalTok{)}
\NormalTok{)}
\NormalTok{inv\_largo}
\end{Highlighting}
\end{Shaded}

\begin{verbatim}
# A tibble: 3 x 4
  parcela dap_2021 dap_2022 dap_2023
    <int>    <dbl>    <dbl>    <dbl>
1       1     15.2     17.8     20.1
2       2     14.8     16.9     19.5
3       3     13.5     15.2     17.8
\end{verbatim}

\begin{Shaded}
\begin{Highlighting}[]
\CommentTok{\# Converter para formato longo}
\NormalTok{inv\_largo }\SpecialCharTok{\%\textgreater{}\%} 
  \FunctionTok{pivot\_longer}\NormalTok{(}
    \AttributeTok{cols =} \FunctionTok{starts\_with}\NormalTok{(}\StringTok{"dap\_"}\NormalTok{),}
    \AttributeTok{names\_to =} \StringTok{"ano"}\NormalTok{,}
    \AttributeTok{values\_to =} \StringTok{"dap"}\NormalTok{,}
    \AttributeTok{names\_prefix =} \StringTok{"dap\_"}
\NormalTok{  ) }\SpecialCharTok{\%\textgreater{}\%} 
  \FunctionTok{mutate}\NormalTok{(}\AttributeTok{ano =} \FunctionTok{as.numeric}\NormalTok{(ano))}
\end{Highlighting}
\end{Shaded}

\begin{verbatim}
# A tibble: 9 x 3
  parcela   ano   dap
    <int> <dbl> <dbl>
1       1  2021  15.2
2       1  2022  17.8
3       1  2023  20.1
4       2  2021  14.8
5       2  2022  16.9
6       2  2023  19.5
7       3  2021  13.5
8       3  2022  15.2
9       3  2023  17.8
\end{verbatim}

\begin{Shaded}
\begin{Highlighting}[]
\CommentTok{\# Agora podemos analisar e plotar}
\NormalTok{inv\_longo }\OtherTok{\textless{}{-}}\NormalTok{ inv\_largo }\SpecialCharTok{\%\textgreater{}\%} 
  \FunctionTok{pivot\_longer}\NormalTok{(}\FunctionTok{starts\_with}\NormalTok{(}\StringTok{"dap\_"}\NormalTok{), }\AttributeTok{names\_to =} \StringTok{"ano"}\NormalTok{, }\AttributeTok{values\_to =} \StringTok{"dap"}\NormalTok{, }\AttributeTok{names\_prefix =} \StringTok{"dap\_"}\NormalTok{) }\SpecialCharTok{\%\textgreater{}\%} 
  \FunctionTok{mutate}\NormalTok{(}\AttributeTok{ano =} \FunctionTok{as.numeric}\NormalTok{(ano))}

\NormalTok{inv\_longo }\SpecialCharTok{\%\textgreater{}\%} 
  \FunctionTok{ggplot}\NormalTok{(}\FunctionTok{aes}\NormalTok{(}\AttributeTok{x =}\NormalTok{ ano, }\AttributeTok{y =}\NormalTok{ dap, }\AttributeTok{color =} \FunctionTok{factor}\NormalTok{(parcela))) }\SpecialCharTok{+}
  \FunctionTok{geom\_line}\NormalTok{(}\AttributeTok{size =} \DecValTok{1}\NormalTok{) }\SpecialCharTok{+}
  \FunctionTok{geom\_point}\NormalTok{(}\AttributeTok{size =} \DecValTok{2}\NormalTok{) }\SpecialCharTok{+}
  \FunctionTok{labs}\NormalTok{(}\AttributeTok{title =} \StringTok{"Evolução do DAP por parcela"}\NormalTok{, }\AttributeTok{x =} \StringTok{"Ano"}\NormalTok{, }\AttributeTok{y =} \StringTok{"DAP (cm)"}\NormalTok{, }\AttributeTok{color =} \StringTok{"Parcela"}\NormalTok{)}
\end{Highlighting}
\end{Shaded}

\pandocbounded{\includegraphics[keepaspectratio]{04-dias/d3-tarde_files/figure-pdf/unnamed-chunk-3-1.pdf}}

\subsection{\texorpdfstring{\texttt{pivot\_wider()} --- longo →
largo}{pivot\_wider() --- longo → largo}}\label{pivot_wider-longo-largo}

\begin{Shaded}
\begin{Highlighting}[]
\CommentTok{\# Tabela{-}resumo: volume por parcela × espécie}
\NormalTok{inv\_florestal }\SpecialCharTok{\%\textgreater{}\%} 
  \FunctionTok{group\_by}\NormalTok{(parcela, especie) }\SpecialCharTok{\%\textgreater{}\%} 
  \FunctionTok{summarise}\NormalTok{(}\AttributeTok{volume\_total =} \FunctionTok{sum}\NormalTok{(volume), }\AttributeTok{.groups =} \StringTok{"drop"}\NormalTok{) }\SpecialCharTok{\%\textgreater{}\%} 
  \FunctionTok{pivot\_wider}\NormalTok{(}
    \AttributeTok{names\_from =}\NormalTok{ especie,}
    \AttributeTok{values\_from =}\NormalTok{ volume\_total,}
    \AttributeTok{values\_fill =} \DecValTok{0}
\NormalTok{  )}
\end{Highlighting}
\end{Shaded}

\begin{verbatim}
# A tibble: 30 x 4
   parcela `Eucalyptus Grandis` `Eucalyptus Urophylla` `Pinus Taeda`
     <dbl>                <dbl>                  <dbl>         <dbl>
 1       1                 19.8                  4.33          12.6 
 2       2                 26.9                  8.17           5.08
 3       3                 26.3                  1.24           9.81
 4       4                 26.8                  2.00          10.3 
 5       5                 28.9                  3.24          10.9 
 6       6                 27.2                  2.73          12.2 
 7       7                 31.1                  2.22           5.00
 8       8                 13.3                  0.548          2.44
 9       9                 26.5                  1.98           3.95
10      10                 16.5                  2.28           5.01
# i 20 more rows
\end{verbatim}

\subsection{Separar e unir colunas}\label{separar-e-unir-colunas}

\begin{Shaded}
\begin{Highlighting}[]
\CommentTok{\# Código da parcela: "T1{-}P2" → talhão + parcela}
\FunctionTok{tibble}\NormalTok{(}\AttributeTok{codigo =} \FunctionTok{c}\NormalTok{(}\StringTok{"T1{-}P1"}\NormalTok{, }\StringTok{"T1{-}P2"}\NormalTok{, }\StringTok{"T2{-}P1"}\NormalTok{)) }\SpecialCharTok{\%\textgreater{}\%} 
  \FunctionTok{separate}\NormalTok{(codigo, }\AttributeTok{into =} \FunctionTok{c}\NormalTok{(}\StringTok{"talhao"}\NormalTok{, }\StringTok{"parcela"}\NormalTok{), }\AttributeTok{sep =} \StringTok{"{-}"}\NormalTok{)}
\end{Highlighting}
\end{Shaded}

\begin{verbatim}
# A tibble: 3 x 2
  talhao parcela
  <chr>  <chr>  
1 T1     P1     
2 T1     P2     
3 T2     P1     
\end{verbatim}

\begin{Shaded}
\begin{Highlighting}[]
\CommentTok{\# Unir colunas}
\FunctionTok{tibble}\NormalTok{(}\AttributeTok{talhao =} \StringTok{"T1"}\NormalTok{, }\AttributeTok{parcela =} \StringTok{"P1"}\NormalTok{) }\SpecialCharTok{\%\textgreater{}\%} 
  \FunctionTok{unite}\NormalTok{(}\StringTok{"codigo"}\NormalTok{, talhao, parcela, }\AttributeTok{sep =} \StringTok{"{-}"}\NormalTok{)}
\end{Highlighting}
\end{Shaded}

\begin{verbatim}
# A tibble: 1 x 1
  codigo
  <chr> 
1 T1-P1 
\end{verbatim}

\subsection{Preencher combinações e valores
ausentes}\label{preencher-combinauxe7uxf5es-e-valores-ausentes}

\begin{Shaded}
\begin{Highlighting}[]
\CommentTok{\# complete() — garante todas as combinações}
\NormalTok{inv\_florestal }\SpecialCharTok{\%\textgreater{}\%} 
  \FunctionTok{count}\NormalTok{(parcela, especie) }\SpecialCharTok{\%\textgreater{}\%} 
  \FunctionTok{complete}\NormalTok{(parcela, especie, }\AttributeTok{fill =} \FunctionTok{list}\NormalTok{(}\AttributeTok{n =} \DecValTok{0}\NormalTok{))}
\end{Highlighting}
\end{Shaded}

\begin{verbatim}
# A tibble: 90 x 3
   parcela especie                  n
     <dbl> <chr>                <int>
 1       1 Eucalyptus Grandis      21
 2       1 Eucalyptus Urophylla     4
 3       1 Pinus Taeda             15
 4       2 Eucalyptus Grandis      23
 5       2 Eucalyptus Urophylla     9
 6       2 Pinus Taeda             10
 7       3 Eucalyptus Grandis      26
 8       3 Eucalyptus Urophylla     2
 9       3 Pinus Taeda             19
10       4 Eucalyptus Grandis      26
# i 80 more rows
\end{verbatim}

\begin{center}\rule{0.5\linewidth}{0.5pt}\end{center}

\section{Importação de múltiplos
arquivos}\label{importauxe7uxe3o-de-muxfaltiplos-arquivos}

\subsection{Múltiplas abas de um
Excel}\label{muxfaltiplas-abas-de-um-excel}

\begin{Shaded}
\begin{Highlighting}[]
\CommentTok{\# Listar abas}
\FunctionTok{read\_csv}\NormalTok{(}\StringTok{"../data/talhoes.csv"}\NormalTok{)}
\end{Highlighting}
\end{Shaded}

\begin{verbatim}
# A tibble: 150 x 4
   talhao   parcela volume_m3 area_ha
   <chr>      <dbl>     <dbl>   <dbl>
 1 Talhão_1      23      143.    2.04
 2 Talhão_1      29      121.    3.01
 3 Talhão_1      18      155.    1.81
 4 Talhão_1       9      209.    2.24
 5 Talhão_1      24      280     3.3 
 6 Talhão_1       5      226.    1.87
 7 Talhão_1      14      218.    3.22
 8 Talhão_1      20      222.    2.96
 9 Talhão_1      10      254.    2.8 
10 Talhão_1      11      230.    1.86
# i 140 more rows
\end{verbatim}

\begin{Shaded}
\begin{Highlighting}[]
\CommentTok{\# Importar todas as abas e combinar}
\NormalTok{talhoes }\OtherTok{\textless{}{-}} \FunctionTok{read\_csv}\NormalTok{(}\StringTok{"../data/talhoes.csv"}\NormalTok{)}
\NormalTok{talhoes}
\end{Highlighting}
\end{Shaded}

\begin{verbatim}
# A tibble: 150 x 4
   talhao   parcela volume_m3 area_ha
   <chr>      <dbl>     <dbl>   <dbl>
 1 Talhão_1      23      143.    2.04
 2 Talhão_1      29      121.    3.01
 3 Talhão_1      18      155.    1.81
 4 Talhão_1       9      209.    2.24
 5 Talhão_1      24      280     3.3 
 6 Talhão_1       5      226.    1.87
 7 Talhão_1      14      218.    3.22
 8 Talhão_1      20      222.    2.96
 9 Talhão_1      10      254.    2.8 
10 Talhão_1      11      230.    1.86
# i 140 more rows
\end{verbatim}

\subsection{Múltiplos CSVs de uma
pasta}\label{muxfaltiplos-csvs-de-uma-pasta}

\begin{Shaded}
\begin{Highlighting}[]
\CommentTok{\# Listar arquivos e importar todos}
\NormalTok{arquivos }\OtherTok{\textless{}{-}} \FunctionTok{list.files}\NormalTok{(}\StringTok{"../data/parcelas\_ano/"}\NormalTok{, }\AttributeTok{pattern =} \StringTok{"*.csv"}\NormalTok{, }\AttributeTok{full.names =} \ConstantTok{TRUE}\NormalTok{)}
\NormalTok{arquivos}
\end{Highlighting}
\end{Shaded}

\begin{verbatim}
character(0)
\end{verbatim}

\begin{Shaded}
\begin{Highlighting}[]
\NormalTok{todos\_dados }\OtherTok{\textless{}{-}} \FunctionTok{map\_dfr}\NormalTok{(arquivos, read\_csv2)}
\FunctionTok{glimpse}\NormalTok{(todos\_dados)}
\end{Highlighting}
\end{Shaded}

\begin{verbatim}
Rows: 0
Columns: 0
\end{verbatim}

\begin{center}\rule{0.5\linewidth}{0.5pt}\end{center}

\section{Lidando com dados ausentes
(NA)}\label{lidando-com-dados-ausentes-na}

\begin{Shaded}
\begin{Highlighting}[]
\CommentTok{\# Identificar NAs}
\NormalTok{inv\_florestal }\SpecialCharTok{\%\textgreater{}\%} 
  \FunctionTok{summarise}\NormalTok{(}\FunctionTok{across}\NormalTok{(}\FunctionTok{everything}\NormalTok{(), }\SpecialCharTok{\textasciitilde{}} \FunctionTok{sum}\NormalTok{(}\FunctionTok{is.na}\NormalTok{(.)))) }\SpecialCharTok{\%\textgreater{}\%} 
  \FunctionTok{pivot\_longer}\NormalTok{(}\FunctionTok{everything}\NormalTok{(), }\AttributeTok{names\_to =} \StringTok{"coluna"}\NormalTok{, }\AttributeTok{values\_to =} \StringTok{"n\_na"}\NormalTok{)}
\end{Highlighting}
\end{Shaded}

\begin{verbatim}
# A tibble: 6 x 2
  coluna   n_na
  <chr>   <int>
1 parcela     0
2 arvore      0
3 especie     0
4 dap         0
5 altura      0
6 volume      0
\end{verbatim}

\begin{Shaded}
\begin{Highlighting}[]
\CommentTok{\# Remover linhas com NA}
\NormalTok{inv\_florestal }\SpecialCharTok{\%\textgreater{}\%} \FunctionTok{drop\_na}\NormalTok{(dap, altura)}
\end{Highlighting}
\end{Shaded}

\begin{verbatim}
# A tibble: 1,131 x 6
   parcela arvore especie              dap altura volume
     <dbl>  <dbl> <chr>              <dbl>  <dbl>  <dbl>
 1       1      1 Eucalyptus Grandis  21.7   46    0.876
 2       1      2 Pinus Taeda         20.5   44    0.699
 3       1      3 Eucalyptus Grandis  26.2   51.1  1.24 
 4       1      4 Pinus Taeda         12.6   38.7  0.217
 5       1      5 Eucalyptus Grandis  19.5   49.8  0.746
 6       1      6 Eucalyptus Grandis  17.9   44.3  0.596
 7       1      7 Pinus Taeda         23.5   49.8  1.02 
 8       1      8 Eucalyptus Grandis  27     52.9  1.56 
 9       1      9 Eucalyptus Grandis  24.5   48.9  1.11 
10       1     10 Pinus Taeda         27.5   50.2  1.57 
# i 1,121 more rows
\end{verbatim}

\begin{Shaded}
\begin{Highlighting}[]
\CommentTok{\# Substituir NA por valor}
\NormalTok{inv\_florestal }\SpecialCharTok{\%\textgreater{}\%} \FunctionTok{replace\_na}\NormalTok{(}\FunctionTok{list}\NormalTok{(}\AttributeTok{altura =} \DecValTok{0}\NormalTok{))}
\end{Highlighting}
\end{Shaded}

\begin{verbatim}
# A tibble: 1,131 x 6
   parcela arvore especie              dap altura volume
     <dbl>  <dbl> <chr>              <dbl>  <dbl>  <dbl>
 1       1      1 Eucalyptus Grandis  21.7   46    0.876
 2       1      2 Pinus Taeda         20.5   44    0.699
 3       1      3 Eucalyptus Grandis  26.2   51.1  1.24 
 4       1      4 Pinus Taeda         12.6   38.7  0.217
 5       1      5 Eucalyptus Grandis  19.5   49.8  0.746
 6       1      6 Eucalyptus Grandis  17.9   44.3  0.596
 7       1      7 Pinus Taeda         23.5   49.8  1.02 
 8       1      8 Eucalyptus Grandis  27     52.9  1.56 
 9       1      9 Eucalyptus Grandis  24.5   48.9  1.11 
10       1     10 Pinus Taeda         27.5   50.2  1.57 
# i 1,121 more rows
\end{verbatim}

\begin{center}\rule{0.5\linewidth}{0.5pt}\end{center}

\section{Exercícios}\label{exercuxedcios-12}

\subsection{Exercício 1}\label{exercuxedcio-1-8}

Dada a tabela no formato largo:

\begin{verbatim}
# A tibble: 3 × 4
  parcela dap_2021 dap_2022 dap_2023
    <int>    <dbl>    <dbl>    <dbl>
1       1     15.2     17.8     20.1
2       2     14.8     16.9     19.5
3       3     13.5     15.2     17.8
\end{verbatim}

Converta para formato longo com \texttt{pivot\_longer()}.

\textbf{Resultado esperado:}

\begin{verbatim}
# A tibble: 9 × 3
  parcela   ano   dap
    <int> <dbl> <dbl>
1       1  2021  15.2
2       1  2022  17.8
3       1  2023  20.1
4       2  2021  14.8
5       2  2022  16.9
6       2  2023  19.5
7       3  2021  13.5
8       3  2022  15.2
9       3  2023  17.8
\end{verbatim}

\subsection{Exercício 2}\label{exercuxedcio-2-8}

Calcule o volume total por \texttt{parcela} × \texttt{especie} no
\texttt{inv\_florestal}. Use \texttt{pivot\_wider()} para criar uma
tabela com espécies nas colunas.

\textbf{Resultado esperado:}

\begin{verbatim}
# A tibble: 30 × 3
   parcela `Eucalyptus Grandis` `Pinus Taeda`
     <dbl>                <dbl>         <dbl>
 1       1                10.3          6.8
 2       2                11.2          7.1
 3       3                 8.9          5.4
\end{verbatim}

\subsection{Exercício 3}\label{exercuxedcio-3-7}

Importe todas as abas do arquivo \texttt{talhoes.xlsx} e combine em um
único data frame. Quantas linhas no total?

\textbf{Resultado esperado:}

\begin{verbatim}
[1] 150
\end{verbatim}

\subsection{Exercício 4}\label{exercuxedcio-4-6}

Verifique com \texttt{complete()} se toda parcela tem registro para toda
espécie. Preencha combinações ausentes com volume 0.

\chapter{Dia 4 --- Manhã}\label{dia-4-manhuxe3}

\textbf{Dia 4 --- Manhã \textbar{} 4 horas}

\section*{Objetivos}\label{objetivos-6}
\addcontentsline{toc}{section}{Objetivos}

\markright{Objetivos}

\begin{itemize}
\tightlist
\item
  Escrever funções próprias para automatizar tarefas repetitivas
\item
  Usar \texttt{purrr::map()} para aplicar funções a múltiplos elementos
\item
  Aninhar dados e ajustar modelos por grupo
\end{itemize}

\begin{center}\rule{0.5\linewidth}{0.5pt}\end{center}

\section{Por que criar funções?}\label{por-que-criar-funuxe7uxf5es}

Se você copiou e colou um código mais de 3 vezes, \textbf{transforme em
função}. O princípio DRY (\emph{Don't Repeat Yourself}) evita erros e
torna seu código reutilizável.

\begin{Shaded}
\begin{Highlighting}[]
\CommentTok{\# SEM função: copiar e colar para cada parcela (ruim!)}
\CommentTok{\# volume\_p1 \textless{}{-} sum(inv\_florestal$volume[inv\_florestal$parcela == 1])}
\CommentTok{\# volume\_p2 \textless{}{-} sum(inv\_florestal$volume[inv\_florestal$parcela == 2])}
\CommentTok{\# ... repetir 30 vezes}

\CommentTok{\# COM função: escreva uma vez, reutilize sempre}
\NormalTok{volume\_parcela }\OtherTok{\textless{}{-}} \ControlFlowTok{function}\NormalTok{(dados, num\_parcela) \{}
\NormalTok{  dados }\SpecialCharTok{\%\textgreater{}\%} 
    \FunctionTok{filter}\NormalTok{(parcela }\SpecialCharTok{==}\NormalTok{ num\_parcela) }\SpecialCharTok{\%\textgreater{}\%} 
    \FunctionTok{summarise}\NormalTok{(}\AttributeTok{volume\_total =} \FunctionTok{sum}\NormalTok{(volume))}
\NormalTok{\}}
\end{Highlighting}
\end{Shaded}

\begin{center}\rule{0.5\linewidth}{0.5pt}\end{center}

\section{Anatomia de uma função}\label{anatomia-de-uma-funuxe7uxe3o}

\begin{Shaded}
\begin{Highlighting}[]
\NormalTok{nome\_da\_funcao }\OtherTok{\textless{}{-}} \ControlFlowTok{function}\NormalTok{(argumento1, }\AttributeTok{argumento2 =}\NormalTok{ valor\_default) \{}
  \CommentTok{\# corpo da função}
\NormalTok{  resultado }\OtherTok{\textless{}{-}}\NormalTok{ argumento1 }\SpecialCharTok{+}\NormalTok{ argumento2}
  \FunctionTok{return}\NormalTok{(resultado)}
\NormalTok{\}}
\end{Highlighting}
\end{Shaded}

\subsection{Exemplos aplicados}\label{exemplos-aplicados}

\begin{Shaded}
\begin{Highlighting}[]
\CommentTok{\# Função para área basal (m²)}
\NormalTok{area\_basal }\OtherTok{\textless{}{-}} \ControlFlowTok{function}\NormalTok{(dap) \{}
\NormalTok{  pi }\SpecialCharTok{*}\NormalTok{ (dap }\SpecialCharTok{/} \DecValTok{200}\NormalTok{)}\SpecialCharTok{\^{}}\DecValTok{2}
\NormalTok{\}}

\FunctionTok{area\_basal}\NormalTok{(}\AttributeTok{dap =} \DecValTok{25}\NormalTok{)}
\end{Highlighting}
\end{Shaded}

\begin{verbatim}
[1] 0.04908739
\end{verbatim}

\begin{Shaded}
\begin{Highlighting}[]
\FunctionTok{area\_basal}\NormalTok{(}\FunctionTok{c}\NormalTok{(}\DecValTok{15}\NormalTok{, }\DecValTok{20}\NormalTok{, }\DecValTok{25}\NormalTok{, }\DecValTok{30}\NormalTok{))}
\end{Highlighting}
\end{Shaded}

\begin{verbatim}
[1] 0.01767146 0.03141593 0.04908739 0.07068583
\end{verbatim}

\begin{Shaded}
\begin{Highlighting}[]
\CommentTok{\# Função para volume cilíndrico}
\NormalTok{volume }\OtherTok{\textless{}{-}} \ControlFlowTok{function}\NormalTok{(dap, altura, }\AttributeTok{ff =} \FloatTok{0.5}\NormalTok{) \{}
  \FunctionTok{area\_basal}\NormalTok{(dap) }\SpecialCharTok{*}\NormalTok{ altura }\SpecialCharTok{*}\NormalTok{ ff}
\NormalTok{\}}

\FunctionTok{volume}\NormalTok{(}\DecValTok{25}\NormalTok{, }\DecValTok{30}\NormalTok{)}
\end{Highlighting}
\end{Shaded}

\begin{verbatim}
[1] 0.7363108
\end{verbatim}

\begin{Shaded}
\begin{Highlighting}[]
\FunctionTok{volume}\NormalTok{(}\DecValTok{25}\NormalTok{, }\DecValTok{30}\NormalTok{, }\AttributeTok{ff =} \FloatTok{0.45}\NormalTok{)}
\end{Highlighting}
\end{Shaded}

\begin{verbatim}
[1] 0.6626797
\end{verbatim}

\begin{Shaded}
\begin{Highlighting}[]
\CommentTok{\# Função para classificar sítio}
\NormalTok{classificar\_sitio }\OtherTok{\textless{}{-}} \ControlFlowTok{function}\NormalTok{(altura\_dominante, idade) \{}
  \ControlFlowTok{if}\NormalTok{ (altura\_dominante }\SpecialCharTok{\textgreater{}} \DecValTok{30}\NormalTok{) }\FunctionTok{return}\NormalTok{(}\StringTok{"Bom"}\NormalTok{)}
  \ControlFlowTok{if}\NormalTok{ (altura\_dominante }\SpecialCharTok{\textgreater{}} \DecValTok{22}\NormalTok{) }\FunctionTok{return}\NormalTok{(}\StringTok{"Médio"}\NormalTok{)}
  \FunctionTok{return}\NormalTok{(}\StringTok{"Ruim"}\NormalTok{)}
\NormalTok{\}}

\FunctionTok{classificar\_sitio}\NormalTok{(}\DecValTok{33}\NormalTok{, }\DecValTok{7}\NormalTok{)}
\end{Highlighting}
\end{Shaded}

\begin{verbatim}
[1] "Bom"
\end{verbatim}

\begin{Shaded}
\begin{Highlighting}[]
\FunctionTok{classificar\_sitio}\NormalTok{(}\DecValTok{25}\NormalTok{, }\DecValTok{7}\NormalTok{)}
\end{Highlighting}
\end{Shaded}

\begin{verbatim}
[1] "Médio"
\end{verbatim}

\begin{Shaded}
\begin{Highlighting}[]
\FunctionTok{classificar\_sitio}\NormalTok{(}\DecValTok{18}\NormalTok{, }\DecValTok{7}\NormalTok{)}
\end{Highlighting}
\end{Shaded}

\begin{verbatim}
[1] "Ruim"
\end{verbatim}

\begin{center}\rule{0.5\linewidth}{0.5pt}\end{center}

\section{\texorpdfstring{Iteração com
\texttt{purrr}}{Iteração com purrr}}\label{iterauxe7uxe3o-com-purrr}

O pacote \texttt{purrr} oferece uma família de funções \texttt{map()}
para aplicar uma função a cada elemento de um vetor ou lista --- sem
loops explícitos.

\subsection{\texorpdfstring{\texttt{map()} --- a base de
tudo}{map() --- a base de tudo}}\label{map-a-base-de-tudo}

\begin{Shaded}
\begin{Highlighting}[]
\CommentTok{\# Calcular volume para várias árvores}
\NormalTok{daps }\OtherTok{\textless{}{-}} \FunctionTok{c}\NormalTok{(}\DecValTok{15}\NormalTok{, }\DecValTok{20}\NormalTok{, }\DecValTok{25}\NormalTok{, }\DecValTok{30}\NormalTok{, }\DecValTok{35}\NormalTok{)}
\NormalTok{alturas }\OtherTok{\textless{}{-}} \FunctionTok{c}\NormalTok{(}\DecValTok{18}\NormalTok{, }\DecValTok{22}\NormalTok{, }\DecValTok{28}\NormalTok{, }\DecValTok{32}\NormalTok{, }\DecValTok{34}\NormalTok{)}

\FunctionTok{map2\_dbl}\NormalTok{(daps, alturas, volume)}
\end{Highlighting}
\end{Shaded}

\begin{verbatim}
[1] 0.1590431 0.3455752 0.6872234 1.1309734 1.6355917
\end{verbatim}

\subsection{Tipos de saída}\label{tipos-de-sauxedda}

\begin{Shaded}
\begin{Highlighting}[]
\CommentTok{\# map\_dbl() → vetor numérico}
\CommentTok{\# map\_chr() → vetor de texto}
\CommentTok{\# map\_dfr() → data frame (empilhado)}
\CommentTok{\# map\_lgl() → vetor lógico}

\CommentTok{\# Exemplo: aplicar área\_basal a uma coluna}
\NormalTok{inv\_florestal }\SpecialCharTok{\%\textgreater{}\%} 
  \FunctionTok{mutate}\NormalTok{(}\AttributeTok{ab =} \FunctionTok{map\_dbl}\NormalTok{(dap, area\_basal)) }\SpecialCharTok{\%\textgreater{}\%} 
  \FunctionTok{select}\NormalTok{(dap, ab) }\SpecialCharTok{\%\textgreater{}\%} 
  \FunctionTok{head}\NormalTok{(}\DecValTok{5}\NormalTok{)}
\end{Highlighting}
\end{Shaded}

\begin{verbatim}
# A tibble: 5 x 2
    dap     ab
  <dbl>  <dbl>
1  21.7 0.0370
2  20.5 0.0330
3  26.2 0.0539
4  12.6 0.0125
5  19.5 0.0299
\end{verbatim}

\subsection{\texorpdfstring{\texttt{map\_dfr()} --- empilhar
resultados}{map\_dfr() --- empilhar resultados}}\label{map_dfr-empilhar-resultados}

\begin{Shaded}
\begin{Highlighting}[]
\CommentTok{\# Calcular estatísticas por parcela e empilhar}
\NormalTok{parcelas }\OtherTok{\textless{}{-}} \FunctionTok{unique}\NormalTok{(inv\_florestal}\SpecialCharTok{$}\NormalTok{parcela)}

\NormalTok{resultados }\OtherTok{\textless{}{-}} \FunctionTok{map\_dfr}\NormalTok{(parcelas, }\ControlFlowTok{function}\NormalTok{(p) \{}
\NormalTok{  inv\_florestal }\SpecialCharTok{\%\textgreater{}\%} 
    \FunctionTok{filter}\NormalTok{(parcela }\SpecialCharTok{==}\NormalTok{ p) }\SpecialCharTok{\%\textgreater{}\%} 
    \FunctionTok{summarise}\NormalTok{(}
      \AttributeTok{parcela =}\NormalTok{ p,}
      \AttributeTok{dap\_medio =} \FunctionTok{mean}\NormalTok{(dap),}
      \AttributeTok{altura\_media =} \FunctionTok{mean}\NormalTok{(altura),}
      \AttributeTok{volume\_total =} \FunctionTok{sum}\NormalTok{(volume)}
\NormalTok{    )}
\NormalTok{\})}

\FunctionTok{head}\NormalTok{(resultados, }\DecValTok{5}\NormalTok{)}
\end{Highlighting}
\end{Shaded}

\begin{verbatim}
# A tibble: 5 x 4
  parcela dap_medio altura_media volume_total
    <dbl>     <dbl>        <dbl>        <dbl>
1       1      20.7         46.1         36.7
2       2      20.6         45.7         40.1
3       3      19.5         44.7         37.3
4       4      20.5         45.7         39.1
5       5      22.0         46.8         43.0
\end{verbatim}

\begin{center}\rule{0.5\linewidth}{0.5pt}\end{center}

\section{Aninhamento: modelos por
grupo}\label{aninhamento-modelos-por-grupo}

Uma das técnicas mais poderosas do \texttt{purrr}: aninhar dados por
grupo e aplicar modelos a cada grupo.

\begin{Shaded}
\begin{Highlighting}[]
\CommentTok{\# Modelo hipsométrico por espécie}
\NormalTok{modelos\_especie }\OtherTok{\textless{}{-}}\NormalTok{ inv\_florestal }\SpecialCharTok{\%\textgreater{}\%} 
  \FunctionTok{group\_by}\NormalTok{(especie) }\SpecialCharTok{\%\textgreater{}\%} 
  \FunctionTok{nest}\NormalTok{() }\SpecialCharTok{\%\textgreater{}\%} 
  \FunctionTok{mutate}\NormalTok{(}
    \AttributeTok{modelo =} \FunctionTok{map}\NormalTok{(data, }\SpecialCharTok{\textasciitilde{}} \FunctionTok{lm}\NormalTok{(altura }\SpecialCharTok{\textasciitilde{}} \FunctionTok{log}\NormalTok{(dap), }\AttributeTok{data =}\NormalTok{ .x)),}
    \AttributeTok{r2 =} \FunctionTok{map\_dbl}\NormalTok{(modelo, }\SpecialCharTok{\textasciitilde{}} \FunctionTok{summary}\NormalTok{(.x)}\SpecialCharTok{$}\NormalTok{r.squared),}
    \AttributeTok{coef =} \FunctionTok{map}\NormalTok{(modelo, broom}\SpecialCharTok{::}\NormalTok{tidy)}
\NormalTok{  )}

\NormalTok{modelos\_especie }\SpecialCharTok{\%\textgreater{}\%} \FunctionTok{select}\NormalTok{(especie, r2)}
\end{Highlighting}
\end{Shaded}

\begin{verbatim}
# A tibble: 3 x 2
# Groups:   especie [3]
  especie                 r2
  <chr>                <dbl>
1 Eucalyptus Grandis   0.889
2 Pinus Taeda          0.928
3 Eucalyptus Urophylla 0.906
\end{verbatim}

\begin{Shaded}
\begin{Highlighting}[]
\CommentTok{\# Extrair coeficientes}
\NormalTok{modelos\_especie }\SpecialCharTok{\%\textgreater{}\%} 
  \FunctionTok{select}\NormalTok{(especie, coef) }\SpecialCharTok{\%\textgreater{}\%} 
  \FunctionTok{unnest}\NormalTok{(coef)}
\end{Highlighting}
\end{Shaded}

\begin{verbatim}
# A tibble: 6 x 6
# Groups:   especie [3]
  especie              term        estimate std.error statistic   p.value
  <chr>                <chr>          <dbl>     <dbl>     <dbl>     <dbl>
1 Eucalyptus Grandis   (Intercept)    0.654     0.632     1.04  3.01e-  1
2 Eucalyptus Grandis   log(dap)      15.2       0.206    73.6   0        
3 Pinus Taeda          (Intercept)    0.484     0.641     0.754 4.51e-  1
4 Pinus Taeda          log(dap)      15.3       0.228    67.0   4.91e-201
5 Eucalyptus Urophylla (Intercept)    2.99      1.44      2.08  4.01e-  2
6 Eucalyptus Urophylla log(dap)      14.4       0.473    30.5   1.58e- 51
\end{verbatim}

\begin{Shaded}
\begin{Highlighting}[]
\CommentTok{\# Plotar curvas ajustadas}
\NormalTok{inv\_florestal }\SpecialCharTok{\%\textgreater{}\%} 
  \FunctionTok{ggplot}\NormalTok{(}\FunctionTok{aes}\NormalTok{(}\AttributeTok{x =}\NormalTok{ dap, }\AttributeTok{y =}\NormalTok{ altura, }\AttributeTok{color =}\NormalTok{ especie)) }\SpecialCharTok{+}
  \FunctionTok{geom\_point}\NormalTok{(}\AttributeTok{alpha =} \FloatTok{0.3}\NormalTok{) }\SpecialCharTok{+}
  \FunctionTok{geom\_smooth}\NormalTok{(}\AttributeTok{method =} \StringTok{"lm"}\NormalTok{, }\AttributeTok{formula =}\NormalTok{ y }\SpecialCharTok{\textasciitilde{}} \FunctionTok{log}\NormalTok{(x), }\AttributeTok{se =} \ConstantTok{FALSE}\NormalTok{, }\AttributeTok{linewidth =} \FloatTok{1.5}\NormalTok{) }\SpecialCharTok{+}
  \FunctionTok{labs}\NormalTok{(}\AttributeTok{title =} \StringTok{"Curvas hipsométricas ajustadas por espécie"}\NormalTok{, }\AttributeTok{x =} \StringTok{"DAP (cm)"}\NormalTok{, }\AttributeTok{y =} \StringTok{"Altura (m)"}\NormalTok{)}
\end{Highlighting}
\end{Shaded}

\pandocbounded{\includegraphics[keepaspectratio]{04-dias/d4-manha_files/figure-pdf/unnamed-chunk-11-1.pdf}}

\subsection{\texorpdfstring{\texttt{walk()} --- Efeitos
colaterais}{walk() --- Efeitos colaterais}}\label{walk-efeitos-colaterais}

\texttt{walk()} é como \texttt{map()}, mas não retorna resultado ---
ideal para salvar arquivos:

\begin{Shaded}
\begin{Highlighting}[]
\CommentTok{\# Gerar e salvar um gráfico por parcela}
\NormalTok{inv\_florestal }\SpecialCharTok{\%\textgreater{}\%} 
  \FunctionTok{group\_by}\NormalTok{(parcela) }\SpecialCharTok{\%\textgreater{}\%} 
  \FunctionTok{group\_walk}\NormalTok{(}\SpecialCharTok{\textasciitilde{}}\NormalTok{ \{}
\NormalTok{    p }\OtherTok{\textless{}{-}} \FunctionTok{ggplot}\NormalTok{(.x, }\FunctionTok{aes}\NormalTok{(}\AttributeTok{x =}\NormalTok{ dap, }\AttributeTok{y =}\NormalTok{ altura)) }\SpecialCharTok{+}
      \FunctionTok{geom\_point}\NormalTok{() }\SpecialCharTok{+} \FunctionTok{geom\_smooth}\NormalTok{(}\AttributeTok{method =} \StringTok{"lm"}\NormalTok{, }\AttributeTok{formula =}\NormalTok{ y }\SpecialCharTok{\textasciitilde{}} \FunctionTok{log}\NormalTok{(x), }\AttributeTok{se =} \ConstantTok{FALSE}\NormalTok{) }\SpecialCharTok{+}
      \FunctionTok{labs}\NormalTok{(}\AttributeTok{title =} \FunctionTok{paste}\NormalTok{(}\StringTok{"Parcela"}\NormalTok{, .y}\SpecialCharTok{$}\NormalTok{parcela))}
    \FunctionTok{ggsave}\NormalTok{(}\FunctionTok{paste0}\NormalTok{(}\StringTok{"parcela\_"}\NormalTok{, .y}\SpecialCharTok{$}\NormalTok{parcela, }\StringTok{".png"}\NormalTok{), p, }\AttributeTok{width =} \DecValTok{6}\NormalTok{, }\AttributeTok{height =} \DecValTok{4}\NormalTok{)}
\NormalTok{  \})}
\end{Highlighting}
\end{Shaded}

\begin{center}\rule{0.5\linewidth}{0.5pt}\end{center}

\section{Exercícios}\label{exercuxedcios-13}

\subsection{Exercício 1}\label{exercuxedcio-1-9}

Escreva uma função \texttt{cv(x)} que calcula o coeficiente de variação
(\texttt{sd(x)\ /\ mean(x)\ *\ 100}). Teste com
\texttt{c(15,\ 20,\ 25,\ 30,\ 35)}.

\textbf{Resultado esperado:}

\begin{verbatim}
[1] 31.62278
\end{verbatim}

\subsection{Exercício 2}\label{exercuxedcio-2-9}

Escreva uma função \texttt{resumo\_parcela(dados,\ num\_parcela)} que
retorna n° de árvores, DAP médio e volume total da parcela. Teste com a
parcela 1.

\textbf{Resultado esperado:}

\begin{verbatim}
# A tibble: 1 × 3
      n dap_medio volume_total
  <int>     <dbl>        <dbl>
1    33      20.3         17.2
\end{verbatim}

\subsection{Exercício 3}\label{exercuxedcio-3-8}

Use \texttt{map\_dbl()} para aplicar a função \texttt{volume()} a um
vetor de DAPs \texttt{c(10,\ 15,\ 20,\ 25,\ 30)} com altura fixa de 25
m.

\textbf{Resultado esperado:}

\begin{verbatim}
[1] 0.098 0.221 0.393 0.614 0.884
\end{verbatim}

\subsection{Exercício 4}\label{exercuxedcio-4-7}

Aninhe os dados por \texttt{parcela} e ajuste
\texttt{lm(altura\ \textasciitilde{}\ log(dap))} para cada parcela. Qual
parcela tem o maior R²?

\textbf{Resultado esperado:}

\begin{verbatim}
# A tibble: 30 × 2
   parcela    r2
     <int> <dbl>
 1       7 0.892
 2      12 0.881
 3      15 0.876
...
\end{verbatim}

\subsection{Exercício 5}\label{exercuxedcio-5-4}

Gere e salve um histograma do DAP para cada espécie usando
\texttt{group\_walk()}.

\chapter{Dia 4 --- Tarde}\label{dia-4-tarde}

\textbf{Dia 4 --- Tarde \textbar{} 4 horas}

\section*{Objetivos}\label{objetivos-7}
\addcontentsline{toc}{section}{Objetivos}

\markright{Objetivos}

\begin{itemize}
\tightlist
\item
  Ajustar modelos de regressão linear com \texttt{lm()}
\item
  Interpretar resultados com \texttt{summary()} e \texttt{broom}
\item
  Realizar ANOVA e comparações múltiplas (Tukey)
\item
  Produzir um relatório técnico completo de inventário florestal
\end{itemize}

\begin{center}\rule{0.5\linewidth}{0.5pt}\end{center}

\section{Regressão linear}\label{regressuxe3o-linear}

A função \texttt{lm()} (linear model) ajusta modelos de regressão. Vamos
aplicá-la a dois contextos florestais clássicos.

\subsection{Relação hipsométrica (altura \textasciitilde{}
DAP)}\label{relauxe7uxe3o-hipsomuxe9trica-altura-dap}

\begin{Shaded}
\begin{Highlighting}[]
\CommentTok{\# Ajuste do modelo}
\NormalTok{mod\_hips }\OtherTok{\textless{}{-}} \FunctionTok{lm}\NormalTok{(altura }\SpecialCharTok{\textasciitilde{}} \FunctionTok{log}\NormalTok{(dap), }\AttributeTok{data =}\NormalTok{ inv\_florestal)}
\FunctionTok{summary}\NormalTok{(mod\_hips)}
\end{Highlighting}
\end{Shaded}

\begin{verbatim}

Call:
lm(formula = altura ~ log(dap), data = inv_florestal)

Residuals:
    Min      1Q  Median      3Q     Max 
-6.9870 -1.3747  0.0588  1.3301  6.2382 

Coefficients:
            Estimate Std. Error t value Pr(>|t|)    
(Intercept)   0.7887     0.4150     1.9   0.0577 .  
log(dap)     15.1632     0.1391   109.0   <2e-16 ***
---
Signif. codes:  0 '***' 0.001 '**' 0.01 '*' 0.05 '.' 0.1 ' ' 1

Residual standard error: 1.946 on 1129 degrees of freedom
Multiple R-squared:  0.9133,    Adjusted R-squared:  0.9132 
F-statistic: 1.189e+04 on 1 and 1129 DF,  p-value: < 2.2e-16
\end{verbatim}

\textbf{Interpretação do \texttt{summary()}:}

\begin{itemize}
\tightlist
\item
  \textbf{Estimate:} coeficientes da equação. Para cada aumento de 1
  unidade em \texttt{log(dap)}, a altura aumenta \texttt{b1} metros.
\item
  \textbf{R²:} proporção da variabilidade da altura explicada pelo DAP.
  0.78 = o DAP explica 78\% da variação da altura.
\item
  \textbf{p-valor:} probabilidade do coeficiente ser zero. p \textless{}
  0.05 = significativo.
\end{itemize}

\begin{Shaded}
\begin{Highlighting}[]
\CommentTok{\# Extrair resultados com broom}
\FunctionTok{tidy}\NormalTok{(mod\_hips)     }\CommentTok{\# tabela de coeficientes}
\end{Highlighting}
\end{Shaded}

\begin{verbatim}
# A tibble: 2 x 5
  term        estimate std.error statistic p.value
  <chr>          <dbl>     <dbl>     <dbl>   <dbl>
1 (Intercept)    0.789     0.415      1.90  0.0577
2 log(dap)      15.2       0.139    109.    0     
\end{verbatim}

\begin{Shaded}
\begin{Highlighting}[]
\FunctionTok{glance}\NormalTok{(mod\_hips)   }\CommentTok{\# estatísticas do modelo (R², AIC, etc.)}
\end{Highlighting}
\end{Shaded}

\begin{verbatim}
# A tibble: 1 x 12
  r.squared adj.r.squared sigma statistic p.value    df logLik   AIC   BIC
      <dbl>         <dbl> <dbl>     <dbl>   <dbl> <dbl>  <dbl> <dbl> <dbl>
1     0.913         0.913  1.95    11886.       0     1 -2357. 4719. 4734.
# i 3 more variables: deviance <dbl>, df.residual <int>, nobs <int>
\end{verbatim}

\begin{Shaded}
\begin{Highlighting}[]
\CommentTok{\# Adicionar valores preditos ao dataset}
\NormalTok{inv\_com\_previsao }\OtherTok{\textless{}{-}} \FunctionTok{augment}\NormalTok{(mod\_hips, }\AttributeTok{data =}\NormalTok{ inv\_florestal)}

\FunctionTok{ggplot}\NormalTok{(inv\_com\_previsao, }\FunctionTok{aes}\NormalTok{(}\AttributeTok{x =}\NormalTok{ dap, }\AttributeTok{y =}\NormalTok{ altura)) }\SpecialCharTok{+}
  \FunctionTok{geom\_point}\NormalTok{(}\AttributeTok{alpha =} \FloatTok{0.3}\NormalTok{) }\SpecialCharTok{+}
  \FunctionTok{geom\_line}\NormalTok{(}\FunctionTok{aes}\NormalTok{(}\AttributeTok{y =}\NormalTok{ .fitted), }\AttributeTok{color =} \StringTok{"\#E74C3C"}\NormalTok{, }\AttributeTok{linewidth =} \DecValTok{1}\NormalTok{) }\SpecialCharTok{+}
  \FunctionTok{labs}\NormalTok{(}\AttributeTok{title =} \StringTok{"Modelo hipsométrico ajustado"}\NormalTok{, }\AttributeTok{x =} \StringTok{"DAP (cm)"}\NormalTok{, }\AttributeTok{y =} \StringTok{"Altura (m)"}\NormalTok{)}
\end{Highlighting}
\end{Shaded}

\pandocbounded{\includegraphics[keepaspectratio]{04-dias/d4-tarde_files/figure-pdf/unnamed-chunk-3-1.pdf}}

\subsection{Modelo de volume}\label{modelo-de-volume}

\begin{Shaded}
\begin{Highlighting}[]
\CommentTok{\# Modelo de Schumacher{-}Hall (log{-}log)}
\NormalTok{mod\_vol }\OtherTok{\textless{}{-}} \FunctionTok{lm}\NormalTok{(}\FunctionTok{log}\NormalTok{(volume) }\SpecialCharTok{\textasciitilde{}} \FunctionTok{log}\NormalTok{(dap) }\SpecialCharTok{+} \FunctionTok{log}\NormalTok{(altura), }\AttributeTok{data =}\NormalTok{ inv\_florestal)}
\FunctionTok{summary}\NormalTok{(mod\_vol)}
\end{Highlighting}
\end{Shaded}

\begin{verbatim}

Call:
lm(formula = log(volume) ~ log(dap) + log(altura), data = inv_florestal)

Residuals:
     Min       1Q   Median       3Q      Max 
-0.37211 -0.07042  0.00353  0.07241  0.32412 

Coefficients:
             Estimate Std. Error t value Pr(>|t|)    
(Intercept) -10.24013    0.16998  -60.24   <2e-16 ***
log(dap)      1.97437    0.02449   80.62   <2e-16 ***
log(altura)   1.04262    0.06238   16.71   <2e-16 ***
---
Signif. codes:  0 '***' 0.001 '**' 0.01 '*' 0.05 '.' 0.1 ' ' 1

Residual standard error: 0.1034 on 1128 degrees of freedom
Multiple R-squared:  0.9891,    Adjusted R-squared:  0.9891 
F-statistic: 5.131e+04 on 2 and 1128 DF,  p-value: < 2.2e-16
\end{verbatim}

\begin{Shaded}
\begin{Highlighting}[]
\CommentTok{\# R² de 0.96 — excelente ajuste!}
\end{Highlighting}
\end{Shaded}

\begin{Shaded}
\begin{Highlighting}[]
\CommentTok{\# Diagnóstico de resíduos}
\FunctionTok{par}\NormalTok{(}\AttributeTok{mfrow =} \FunctionTok{c}\NormalTok{(}\DecValTok{2}\NormalTok{, }\DecValTok{2}\NormalTok{))}
\FunctionTok{plot}\NormalTok{(mod\_vol)}
\end{Highlighting}
\end{Shaded}

\pandocbounded{\includegraphics[keepaspectratio]{04-dias/d4-tarde_files/figure-pdf/unnamed-chunk-5-1.pdf}}

\begin{Shaded}
\begin{Highlighting}[]
\FunctionTok{par}\NormalTok{(}\AttributeTok{mfrow =} \FunctionTok{c}\NormalTok{(}\DecValTok{1}\NormalTok{, }\DecValTok{1}\NormalTok{))}
\end{Highlighting}
\end{Shaded}

\begin{tcolorbox}[enhanced jigsaw, arc=.35mm, bottomrule=.15mm, bottomtitle=1mm, breakable, colback=white, colbacktitle=quarto-callout-tip-color!10!white, colframe=quarto-callout-tip-color-frame, coltitle=black, left=2mm, leftrule=.75mm, opacityback=0, opacitybacktitle=0.6, rightrule=.15mm, title=\textcolor{quarto-callout-tip-color}{\faLightbulb}\hspace{0.5em}{Dica de modelagem florestal}, titlerule=0mm, toprule=.15mm, toptitle=1mm]

Sempre transforme \texttt{volume} com \texttt{log()} antes de modelar. O
volume tem relação exponencial com DAP e altura --- o log lineariza essa
relação, melhorando o ajuste. Para obter o volume na escala original,
use \texttt{exp()} na predição.

\end{tcolorbox}

\begin{center}\rule{0.5\linewidth}{0.5pt}\end{center}

\section{ANOVA --- comparando grupos}\label{anova-comparando-grupos}

A ANOVA (\emph{Analysis of Variance}) testa se há diferença
significativa entre as médias de três ou mais grupos.

\begin{Shaded}
\begin{Highlighting}[]
\CommentTok{\# Volume difere entre espécies?}
\NormalTok{mod\_aov }\OtherTok{\textless{}{-}} \FunctionTok{aov}\NormalTok{(volume }\SpecialCharTok{\textasciitilde{}}\NormalTok{ especie, }\AttributeTok{data =}\NormalTok{ inv\_florestal)}
\FunctionTok{summary}\NormalTok{(mod\_aov)}
\end{Highlighting}
\end{Shaded}

\begin{verbatim}
              Df Sum Sq Mean Sq F value Pr(>F)    
especie        2   35.9  17.965   39.69 <2e-16 ***
Residuals   1128  510.5   0.453                   
---
Signif. codes:  0 '***' 0.001 '**' 0.01 '*' 0.05 '.' 0.1 ' ' 1
\end{verbatim}

\begin{Shaded}
\begin{Highlighting}[]
\CommentTok{\# Quais espécies diferem entre si?}
\FunctionTok{TukeyHSD}\NormalTok{(mod\_aov)}
\end{Highlighting}
\end{Shaded}

\begin{verbatim}
  Tukey multiple comparisons of means
    95% family-wise confidence level

Fit: aov(formula = volume ~ especie, data = inv_florestal)

$especie
                                                 diff        lwr        upr
Eucalyptus Urophylla-Eucalyptus Grandis  0.0002599707 -0.1695454  0.1700654
Pinus Taeda-Eucalyptus Grandis          -0.3855308345 -0.4893433 -0.2817184
Pinus Taeda-Eucalyptus Urophylla        -0.3857908052 -0.5655153 -0.2060663
                                            p adj
Eucalyptus Urophylla-Eucalyptus Grandis 0.9999929
Pinus Taeda-Eucalyptus Grandis          0.0000000
Pinus Taeda-Eucalyptus Urophylla        0.0000016
\end{verbatim}

\begin{Shaded}
\begin{Highlighting}[]
\CommentTok{\# Visualização: boxplot + letras de Tukey}
\NormalTok{inv\_florestal }\SpecialCharTok{\%\textgreater{}\%} 
  \FunctionTok{ggplot}\NormalTok{(}\FunctionTok{aes}\NormalTok{(}\AttributeTok{x =}\NormalTok{ especie, }\AttributeTok{y =}\NormalTok{ volume, }\AttributeTok{fill =}\NormalTok{ especie)) }\SpecialCharTok{+}
  \FunctionTok{geom\_boxplot}\NormalTok{() }\SpecialCharTok{+}
  \FunctionTok{labs}\NormalTok{(}\AttributeTok{title =} \StringTok{"Volume por espécie"}\NormalTok{, }\AttributeTok{x =} \ConstantTok{NULL}\NormalTok{, }\AttributeTok{y =} \FunctionTok{expression}\NormalTok{(Volume}\SpecialCharTok{\textasciitilde{}}\NormalTok{(m}\SpecialCharTok{\^{}}\DecValTok{3}\NormalTok{)),}
       \AttributeTok{subtitle =} \StringTok{"Letras diferentes = diferença significativa (Tukey, p \textless{} 0.05)"}\NormalTok{) }\SpecialCharTok{+}
  \FunctionTok{theme}\NormalTok{(}\AttributeTok{legend.position =} \StringTok{"none"}\NormalTok{)}
\end{Highlighting}
\end{Shaded}

\pandocbounded{\includegraphics[keepaspectratio]{04-dias/d4-tarde_files/figure-pdf/unnamed-chunk-8-1.pdf}}

\subsection{Delineamento experimental}\label{delineamento-experimental}

\begin{Shaded}
\begin{Highlighting}[]
\NormalTok{experimento }\OtherTok{\textless{}{-}} \FunctionTok{read\_csv}\NormalTok{(}\StringTok{"../data/experimento.csv"}\NormalTok{)}
\FunctionTok{glimpse}\NormalTok{(experimento)}
\end{Highlighting}
\end{Shaded}

\begin{verbatim}
Rows: 16
Columns: 3
$ bloco      <dbl> 1, 1, 1, 1, 2, 2, 2, 2, 3, 3, 3, 3, 4, 4, 4, 4
$ tratamento <chr> "Controle", "Adubação", "Irrigação", "Adub+Irrig", "Control~
$ resposta   <dbl> 24.1, 27.9, 33.4, 36.6, 19.1, 25.6, 31.3, 37.7, 30.6, 35.4,~
\end{verbatim}

\begin{Shaded}
\begin{Highlighting}[]
\CommentTok{\# ANOVA para delineamento experimental}
\NormalTok{mod\_exp }\OtherTok{\textless{}{-}} \FunctionTok{aov}\NormalTok{(resposta }\SpecialCharTok{\textasciitilde{}}\NormalTok{ tratamento, }\AttributeTok{data =}\NormalTok{ experimento)}
\FunctionTok{summary}\NormalTok{(mod\_exp)}
\end{Highlighting}
\end{Shaded}

\begin{verbatim}
            Df Sum Sq Mean Sq F value  Pr(>F)   
tratamento   3  245.2   81.74   6.476 0.00745 **
Residuals   12  151.5   12.62                   
---
Signif. codes:  0 '***' 0.001 '**' 0.01 '*' 0.05 '.' 0.1 ' ' 1
\end{verbatim}

\begin{Shaded}
\begin{Highlighting}[]
\CommentTok{\# Comparações múltiplas}
\FunctionTok{TukeyHSD}\NormalTok{(mod\_exp) }\SpecialCharTok{\%\textgreater{}\%} 
\NormalTok{  broom}\SpecialCharTok{::}\FunctionTok{tidy}\NormalTok{() }\SpecialCharTok{\%\textgreater{}\%} 
  \FunctionTok{filter}\NormalTok{(adj.p.value }\SpecialCharTok{\textless{}} \FloatTok{0.05}\NormalTok{)}
\end{Highlighting}
\end{Shaded}

\begin{verbatim}
# A tibble: 2 x 7
  term       contrast         null.value estimate conf.low conf.high adj.p.value
  <chr>      <chr>                 <dbl>    <dbl>    <dbl>     <dbl>       <dbl>
1 tratamento Controle-Adub+I~          0    -10.4  -17.9       -2.94     0.00648
2 tratamento Irrigação-Contr~          0      7.9    0.442     15.4      0.0369 
\end{verbatim}

\begin{center}\rule{0.5\linewidth}{0.5pt}\end{center}

\section{Comunicando resultados}\label{comunicando-resultados}

\subsection{\texorpdfstring{Tabela formatada com
\texttt{knitr::kable()}}{Tabela formatada com knitr::kable()}}\label{tabela-formatada-com-knitrkable}

\begin{Shaded}
\begin{Highlighting}[]
\NormalTok{inv\_florestal }\SpecialCharTok{\%\textgreater{}\%} 
  \FunctionTok{group\_by}\NormalTok{(especie) }\SpecialCharTok{\%\textgreater{}\%} 
  \FunctionTok{summarise}\NormalTok{(}
    \StringTok{\textasciigrave{}}\AttributeTok{N° árvores}\StringTok{\textasciigrave{}} \OtherTok{=} \FunctionTok{n}\NormalTok{(),}
    \StringTok{\textasciigrave{}}\AttributeTok{DAP médio (cm)}\StringTok{\textasciigrave{}} \OtherTok{=} \FunctionTok{round}\NormalTok{(}\FunctionTok{mean}\NormalTok{(dap), }\DecValTok{1}\NormalTok{),}
    \StringTok{\textasciigrave{}}\AttributeTok{Altura média (m)}\StringTok{\textasciigrave{}} \OtherTok{=} \FunctionTok{round}\NormalTok{(}\FunctionTok{mean}\NormalTok{(altura), }\DecValTok{1}\NormalTok{),}
    \StringTok{\textasciigrave{}}\AttributeTok{Volume total (m³)}\StringTok{\textasciigrave{}} \OtherTok{=} \FunctionTok{round}\NormalTok{(}\FunctionTok{sum}\NormalTok{(volume), }\DecValTok{2}\NormalTok{)}
\NormalTok{  ) }\SpecialCharTok{\%\textgreater{}\%} 
\NormalTok{  knitr}\SpecialCharTok{::}\FunctionTok{kable}\NormalTok{(}
    \AttributeTok{caption =} \StringTok{"Resumo do inventário florestal por espécie"}\NormalTok{,}
    \AttributeTok{align =} \StringTok{"c"}
\NormalTok{  )}
\end{Highlighting}
\end{Shaded}

\begin{longtable}[]{@{}
  >{\centering\arraybackslash}p{(\linewidth - 8\tabcolsep) * \real{0.2529}}
  >{\centering\arraybackslash}p{(\linewidth - 8\tabcolsep) * \real{0.1379}}
  >{\centering\arraybackslash}p{(\linewidth - 8\tabcolsep) * \real{0.1839}}
  >{\centering\arraybackslash}p{(\linewidth - 8\tabcolsep) * \real{0.2069}}
  >{\centering\arraybackslash}p{(\linewidth - 8\tabcolsep) * \real{0.2184}}@{}}
\caption{Resumo do inventário florestal por espécie}\tabularnewline
\toprule\noalign{}
\begin{minipage}[b]{\linewidth}\centering
especie
\end{minipage} & \begin{minipage}[b]{\linewidth}\centering
N° árvores
\end{minipage} & \begin{minipage}[b]{\linewidth}\centering
DAP médio (cm)
\end{minipage} & \begin{minipage}[b]{\linewidth}\centering
Altura média (m)
\end{minipage} & \begin{minipage}[b]{\linewidth}\centering
Volume total (m³)
\end{minipage} \\
\midrule\noalign{}
\endfirsthead
\toprule\noalign{}
\begin{minipage}[b]{\linewidth}\centering
especie
\end{minipage} & \begin{minipage}[b]{\linewidth}\centering
N° árvores
\end{minipage} & \begin{minipage}[b]{\linewidth}\centering
DAP médio (cm)
\end{minipage} & \begin{minipage}[b]{\linewidth}\centering
Altura média (m)
\end{minipage} & \begin{minipage}[b]{\linewidth}\centering
Volume total (m³)
\end{minipage} \\
\midrule\noalign{}
\endhead
\bottomrule\noalign{}
\endlastfoot
Eucalyptus Grandis & 682 & 22.1 & 46.9 & 710.81 \\
Eucalyptus Urophylla & 99 & 21.7 & 46.4 & 103.21 \\
Pinus Taeda & 350 & 17.6 & 42.9 & 229.85 \\
\end{longtable}

\begin{center}\rule{0.5\linewidth}{0.5pt}\end{center}

\section{Projeto Final}\label{projeto-final}

\begin{tcolorbox}[enhanced jigsaw, arc=.35mm, bottomrule=.15mm, bottomtitle=1mm, breakable, colback=white, colbacktitle=quarto-callout-important-color!10!white, colframe=quarto-callout-important-color-frame, coltitle=black, left=2mm, leftrule=.75mm, opacityback=0, opacitybacktitle=0.6, rightrule=.15mm, title=\textcolor{quarto-callout-important-color}{\faExclamation}\hspace{0.5em}{Seu desafio final}, titlerule=0mm, toprule=.15mm, toptitle=1mm]

Você recebeu um dataset de inventário florestal. Sua missão é produzir
um \textbf{mini-relatório técnico} que responda às perguntas abaixo. Use
tudo que aprendeu nos 4 dias.

\end{tcolorbox}

\subsection{Roteiro}\label{roteiro}

\begin{enumerate}
\def\labelenumi{\arabic{enumi}.}
\tightlist
\item
  \textbf{Importar} o dataset \texttt{inv\_florestal.csv}
\item
  \textbf{Calcular} área basal e volume por árvore
\item
  \textbf{Explorar} graficamente: distribuição diamétrica, relação
  hipsométrica, volume por espécie
\item
  \textbf{Modelar} a relação altura \textasciitilde{} DAP e volume
  \textasciitilde{} DAP + altura
\item
  \textbf{Comparar} volume entre espécies via ANOVA + Tukey
\item
  \textbf{Produzir} uma tabela-resumo formatada com \texttt{gt}
\end{enumerate}

\subsection{Entregáveis}\label{entreguxe1veis}

\begin{itemize}
\tightlist
\item
  Um gráfico de distribuição diamétrica
\item
  Um gráfico de relação hipsométrica com curva ajustada
\item
  Um boxplot de volume por espécie com letras de Tukey
\item
  Uma tabela-resumo com estatísticas por espécie
\item
  A equação do modelo de volume com R²
\end{itemize}

\begin{center}\rule{0.5\linewidth}{0.5pt}\end{center}

\section{Exercícios}\label{exercuxedcios-14}

\subsection{Exercício 1}\label{exercuxedcio-1-10}

Ajuste
\texttt{lm(altura\ \textasciitilde{}\ log(dap),\ data\ =\ inv\_florestal)}.
Qual o R² do modelo?

\textbf{Resultado esperado:}

\begin{verbatim}
[1] 0.78
\end{verbatim}

\subsection{Exercício 2}\label{exercuxedcio-2-10}

Ajuste
\texttt{lm(log(volume)\ \textasciitilde{}\ log(dap)\ +\ log(altura),\ data\ =\ inv\_florestal)}.
Qual variável tem maior influência (maior coeficiente padronizado)?

\textbf{Resultado esperado:}

\begin{verbatim}
# A tibble: 3 × 2
  term         estimate
  <chr>           <dbl>
1 (Intercept)    -9.59 
2 log(dap)        1.89 
3 log(altura)     0.90
\end{verbatim}

\emph{log(dap) tem o maior coeficiente (1.89).}

\subsection{Exercício 3}\label{exercuxedcio-3-9}

Faça uma ANOVA:
\texttt{aov(volume\ \textasciitilde{}\ especie,\ data\ =\ inv\_florestal)}.
As espécies diferem significativamente em volume?

\textbf{Resultado esperado:}

\begin{verbatim}
             Df Sum Sq Mean Sq F value Pr(>F)    
especie       1   5.93   5.929   40.52 <2e-16 ***
Residuals   990 144.84   0.146
\end{verbatim}

\emph{Sim, p \textless{} 0.001 --- há diferença significativa entre as
espécies.}

\subsection{Exercício 4}\label{exercuxedcio-4-8}

Use \texttt{TukeyHSD()} para verificar qual espécie tem maior volume.
Qual a diferença média?

\textbf{Resultado esperado:}

\begin{verbatim}
                                      diff      lwr      upr p adj
Pinus Taeda-Eucalyptus Grandis -0.1664149 -0.21786 -0.11497     0
\end{verbatim}

\emph{Eucalyptus Grandis tem em média 0.17 m³ a mais que Pinus Taeda.}

\subsection{Exercício 5 (Projeto
final)}\label{exercuxedcio-5-projeto-final}

Produza uma tabela-resumo com \texttt{gt()} contendo: n° de árvores, DAP
médio, altura média, volume total e volume médio por espécie.

\textbf{Resultado esperado:}

\emph{Tabela formatada com 2 linhas (uma por espécie) e 5 colunas.}


\backmatter


\end{document}
